\chapter{Sequences, series, limits}
\label{vol02:ch04}

\epigraph{All science is dominated by the idea of approximation.}{Bertrand
Russell, \emph{The Scientific Outlook} (1931), Ch.~3.}

\chapmeta{Half-life: foundational. Archetypes: scaling (Taylor expansion
as the formal scaling near a working point); uncertainty (the engineer's
first formal handling of approximation error). Exercise target: 30.
Project track: analysis only. Hazard class: none. Reader path default:
standard, with sections explicitly tagged. Named cases: none.}

\noindent
A sequence is a list of numbers indexed by the integers; a series is the
running sum of a sequence; a limit is the value the sequence or its
partial sum approaches as the index runs. The engineering content of
these three definitions is the question that runs through the chapter:
when the engineer truncates a series after a finite number of terms, or
evaluates a function near a working point, how large is the residual,
and what controls its size? This chapter answers that question with the
machinery of geometric and Taylor series, with the convergence tests
engineers use, and with limits as the formal language of behaviour at
the edge of an operating range. The chapter is the bridge into calculus.
Calculus formally arrives in Chapter~5; this chapter prepares the reader
to read its derivations as expansions of objects already understood.

\section{Sequences and their behaviour at infinity}\pathtag{core}

A \engterm{sequence} is a function from the natural numbers $\N$ (or
$\N \cup \{0\}$) to a set, usually $\R$ or $\C$. We write the sequence
as $(a_{n})_{n \geq 1}$ or simply $a_{1}, a_{2}, a_{3}, \ldots$. The
index $n$ is the sequence's position; the value $a_{n}$ is the
sequence's $n$-th term.

Engineering sequences arrive from two sources. The first is iteration:
a numerical method produces $x_{n+1} = g(x_{n})$ from a starting value
$x_{0}$, and the question is whether $x_{n}$ approaches a useful answer
as $n$ grows. The second is sampling: a continuous quantity is recorded
at times $t_{n}$, producing the sequence $a_{n} = f(t_{n})$, and the
question is what the long-time average or the behaviour at large $n$
tells us about the underlying process. Both situations turn on the
same question: does the sequence settle, and to what?

\subsection{Convergence}

\engterm{Convergence}. A sequence $(a_{n})$ \engterm{converges} to a
limit $L$ if, for every tolerance $\varepsilon > 0$, there exists an
index $N$ such that $\abs{a_{n} - L} < \varepsilon$ for all $n \geq N$.
We write $a_{n} \to L$ or $\lim_{n \to \infty} a_{n} = L$. A sequence
that does not converge is said to \engterm{diverge}.

The definition has the formal flavour of an epsilon argument; we will
return to that style in section 4.5 when we treat limits of functions.
The working content for sequences is simpler. The reader picks how close
the sequence has to get; the definition guarantees an index past which
the sequence stays at least that close, forever. If the reader can write
down such an $N$ for any $\varepsilon$, the sequence converges. If the
reader can show that no such $N$ exists for some $\varepsilon$, it does
not.

\engterm{Bounded}. A sequence is \engterm{bounded} if there is a single
$M$ with $\abs{a_{n}} \leq M$ for all $n$. Bounded does not imply
convergent: $a_{n} = (-1)^{n}$ is bounded between $-1$ and $1$ but
oscillates forever. Convergent does imply bounded: a sequence that
eventually sits within $\varepsilon$ of $L$ has all but finitely many
terms in $[L - \varepsilon, L + \varepsilon]$, and the finitely many
remaining terms add a finite bound.

\engterm{Monotone}. A sequence is \engterm{monotone increasing} if
$a_{n+1} \geq a_{n}$ for all $n$, and \engterm{monotone decreasing} if
$a_{n+1} \leq a_{n}$. The \engterm{monotone convergence theorem} says
that a bounded monotone sequence converges. The proof depends on the
\engterm{completeness} of $\R$: every nonempty bounded subset of the
real numbers has a least upper bound. The engineering use of the
theorem is direct. If a numerical iteration produces a sequence that
is bounded and monotone (each step makes the residual smaller, and the
residuals do not run off to infinity), the iteration converges, even
when we cannot compute the limit directly.

\engterm{Cauchy}. A sequence is \engterm{Cauchy} if its terms become
arbitrarily close to one another: for every $\varepsilon > 0$ there is
an $N$ such that $\abs{a_{m} - a_{n}} < \varepsilon$ for all
$m, n \geq N$. In $\R$, a sequence converges if and only if it is
Cauchy. The distinction matters in practice because the Cauchy
condition can be checked without knowing the limit. A numerical
iteration whose successive iterates differ by less than the tolerance,
and continue to do so, has met the Cauchy criterion and may be
stopped. The stopping rule $\abs{a_{n+1} - a_{n}} < \varepsilon$ is the
Cauchy criterion truncated to consecutive terms, and it can mislead
when the gaps shrink while the running total still drifts, exactly the
harmonic-series trap of section 4.6.

A short worked check of the monotone convergence theorem fixes the
machinery. Consider the iteration $a_{n+1} = \tfrac{1}{2}(a_{n} +
2/a_{n})$ with $a_{0} = 2$, the Babylonian rule for $\sqrt{2}$. Each
iterate is positive, and the arithmetic-geometric-mean inequality gives
$a_{n+1} = \tfrac{1}{2}(a_{n} + 2/a_{n}) \geq \sqrt{a_{n} \cdot 2/a_{n}}
= \sqrt{2}$, so every term from $a_{1}$ on is bounded below by
$\sqrt{2}$. The sequence is also decreasing for $n \geq 1$: the
difference $a_{n+1} - a_{n} = \tfrac{1}{2}(2/a_{n} - a_{n}) =
(2 - a_{n}^{2})/(2 a_{n}) \leq 0$ once $a_{n} \geq \sqrt{2}$. A
sequence that is decreasing and bounded below converges, by the
monotone convergence theorem, and the limit $L$ satisfies $L =
\tfrac{1}{2}(L + 2/L)$, that is $L = \sqrt{2}$. The argument certifies
convergence and identifies the limit without ever computing the decimal
expansion. The numerical run of this iteration is the subject of a
simulation exercise; the present point is that boundedness and
monotonicity, both provable by hand, settle the convergence question.

\subsection{Standard sequences and their limits}

A small inventory carries most engineering work.

\begin{itemize}[leftmargin=*]
  \item \engterm{Geometric}: $a_{n} = r^{n}$. Converges to $0$ for
    $\abs{r} < 1$; converges to $1$ for $r = 1$; diverges otherwise.
    The engineering quantity behind this sequence is repeated
    multiplication by a fixed factor: capacitor voltage discharging
    through a resistor in equal time steps, a dropped fraction of an
    investment compounded each period, a numerical iteration whose
    error shrinks by a fixed ratio.
  \item \engterm{Power-law}: $a_{n} = n^{p}$. Converges to $0$ for
    $p < 0$, to $1$ for $p = 0$, and diverges for $p > 0$. The
    \engterm{harmonic} sequence $a_{n} = 1/n$ converges to $0$ but, as
    we will see, its partial sums diverge.
  \item \engterm{Logarithmic}: $a_{n} = \log n / n$ converges to $0$.
    Logarithms grow more slowly than any positive power.
  \item \engterm{Factorial}: $a_{n} = c^{n} / n!$ converges to $0$ for
    every constant $c$. The factorial dominates exponential growth.
\end{itemize}

The relative growth rates compose into a working hierarchy: $\log n
\ll n^{p} \ll c^{n} \ll n!$ as $n \to \infty$, for any $p > 0$ and any
$c > 1$. We use $\ll$ to mean that the ratio of left-hand to right-hand
tends to zero. The hierarchy will reappear in algorithm analysis
(Volume II Chapter 17) and in the asymptotic behaviour of physical
expansions later in this chapter.

The hierarchy is worth a worked confirmation, because the engineer who
trusts it without seeing it justified once will misapply it under
pressure. Take the claim $c^{n} \ll n!$ for $c = 10$. The ratio
$10^{n}/n!$ is itself a sequence; its term-to-term ratio is
$10^{n+1}/(n+1)! \div 10^{n}/n! = 10/(n+1)$. The ratio drops below $1$
once $n \geq 10$, so the sequence $10^{n}/n!$ increases up to $n = 9$,
peaks near $n = 9$ at about $2.8 \times 10^{3}$, and decreases
thereafter, each step multiplying by a factor below $1$. Past the peak
the sequence is geometric-like with a shrinking ratio, so it tends to
zero. The factorial wins, but only after the exponential has had its
run. An engineer who truncates an exponential-over-factorial series too
early, before the peak, sees terms still growing and may wrongly
conclude the series diverges. The ratio test reads the post-peak
behaviour and returns the correct verdict; the eye, watching the
pre-peak terms, does not.

\begin{figure}[ht]
\centering
\begin{tikzpicture}[x=1cm, y=1cm, font=\small]

% Axes
\draw[->, thick] (-0.2, 0) -- (10.5, 0) node[right] {$n$};
\draw[->, thick] (0, -0.2) -- (0, 4.7) node[above] {$|r|^{n}$ (log scale)};

% y-ticks: log10 from 0 to -4
\foreach \y/\lbl in {4/$1$, 3/$10^{-1}$, 2/$10^{-2}$, 1/$10^{-3}$, 0.2/$10^{-4}$} {
  \draw (-0.1, \y) -- (0.1, \y) node[left, xshift=-0.15cm] {\lbl};
}
% x-ticks
\foreach \x in {2, 4, 6, 8, 10} {
  \draw (\x, -0.1) -- (\x, 0.1) node[below, yshift=-0.15cm] {\x};
}

% r = 0.5: log10(0.5^n) = -0.301 n; map y = 4 + 4*log10(value), i.e. y_n = 4 - 0.301 n * (1)
% Use y = 4 + log10(0.5^n) where each decade is 1 cm vertically; but we have 4cm spanning 4 decades
% So plotting points: y(n) = 4 + log10(r^n)*1, but careful: at log=-4, y=0.2 (close to 0)
% Use a coordinate map: y_unit = 4 - n*|log10(r)|
\draw[thick, engblue, mark=*, mark size=1pt] plot coordinates {
  (0,4) (1,3.699) (2,3.398) (3,3.097) (4,2.796) (5,2.495) (6,2.194) (7,1.893) (8,1.592) (9,1.291) (10,0.990)
};
\node[engblue, anchor=west] at (10.1, 0.99) {$r=0.5$};

% r = 0.9: log10(0.9^n) = -0.0458 n
\draw[thick, engred, mark=square*, mark size=1pt] plot coordinates {
  (0,4) (1,3.954) (2,3.908) (3,3.863) (4,3.817) (5,3.771) (6,3.725) (7,3.679) (8,3.634) (9,3.588) (10,3.542)
};
\node[engred, anchor=west] at (10.1, 3.54) {$r=0.9$};

% r = 0.1: log10(0.1^n) = -n; only first few visible
\draw[thick, engblack, mark=triangle*, mark size=1pt] plot coordinates {
  (0,4) (1,3) (2,2) (3,1) (4,0.2)
};
\node[engblack, anchor=west] at (4.1, 0.4) {$r=0.1$};

% Tolerance line at 10^-3
\draw[dashed, gray] (0, 1) -- (10.3, 1) node[right, gray] {$10^{-3}$};

\end{tikzpicture}
\caption[Geometric-sequence decay on a log scale]{Three geometric
sequences $|r|^{n}$ plotted on a log scale. The decay rate depends
linearly on $\log|r|$, so a sequence with $r = 0.1$ crosses a
$10^{-3}$ tolerance in three iterations, a sequence with $r = 0.5$ in
ten, and a sequence with $r = 0.9$ in sixty-six. Iteration counts to
reach a fixed tolerance scale as $-\log(\text{tolerance})/\log(1/|r|)$.
This is the visual content behind the second estimation block on
convergence-rate budgets.}
\label{fig:vol02:ch04:geometric-convergence}
\end{figure}


The figure makes the engineering content of the geometric sequence
explicit. A residual that decays with ratio $r$ requires roughly
$\lceil \log(\text{tolerance}) / \log r \rceil$ iterations to fall
below a target. Tightening the tolerance by a factor of ten costs
$\lceil 1 / \log_{10}(1/|r|) \rceil$ extra iterations. The cost is
manageable for $r \leq 0.5$ and brutal for $r \geq 0.9$. The same
arithmetic applies to a discrete-time controller whose tracking
error decays by a factor $r$ per sample, to a stochastic-gradient
optimiser with effective convergence rate $r$, and to any iteration
whose error contracts linearly.

\subsection{Recursive sequences}

Many engineering sequences are defined recursively: $a_{n+1} = g(a_{n})$
with a starting value $a_{0}$. Three patterns occur often enough to
deserve naming.

\engterm{Linear recursion}: $a_{n+1} = r a_{n} + c$. The fixed point is
$a^{*} = c / (1 - r)$ when $r \neq 1$. The deviation $b_{n} = a_{n} -
a^{*}$ satisfies $b_{n+1} = r b_{n}$, so $b_{n} = r^{n} b_{0}$. The
sequence converges to $a^{*}$ for $\abs{r} < 1$ and diverges for
$\abs{r} > 1$. The case $r = 1$ is a random walk in disguise: $a_{n} =
a_{0} + n c$ grows or decays linearly. This recursion appears in
discrete-time models of cooling, tank filling, and savings.

\engterm{Newton's iteration}: $a_{n+1} = a_{n} - f(a_{n}) / f'(a_{n})$
for finding a root of $f$. The iteration converges quadratically near
a simple root: the new error is approximately the square of the old
error, divided by twice the derivative at the root. We treat this
formally in Volume II Chapter 16; the present-chapter point is that the
question "does this iteration converge?" is a question about the
sequence the iteration produces.

\engterm{Fixed-point iteration}: $a_{n+1} = g(a_{n})$ converges to a
fixed point $a^{*} = g(a^{*})$ when $\abs{g'(a^{*})} < 1$ and the
starting value is close enough. The convergence is linear with rate
$\abs{g'(a^{*})}$. Engineering iterations from heat-transfer
relaxation to PageRank are fixed-point iterations in disguise.

A worked fixed-point example shows the rate-prediction working. The
Colebrook equation for the Darcy friction factor $f$ in a turbulent
pipe is implicit, but a common engineering rearrangement solves the
nearby problem $a_{n+1} = g(a_{n})$ with $g(a) = \cos a$, the iteration
that locates the Dottie number, the unique real solution of $a =
\cos a$. Starting from $a_{0} = 1$ (radian), the iterates run
$1.0000$, $0.5403$, $0.8576$, $0.6543$, $0.7935$, $0.7014$, $0.7639$,
$0.7221$, and settle toward $a^{*} = 0.739085$. The fixed point
satisfies $g'(a^{*}) = -\sin(0.739085) = -0.6736$, and $\abs{g'} =
0.6736 < 1$, so the iteration converges, linearly, with each error
multiplied by about $0.674$ per step. The number of iterations to
shrink the error by a factor of $10^{6}$ is therefore $6 /
\log_{10}(1/0.674) \approx 35$. The alternating overshoot in the
iterate list is the signature of a negative $g'$: successive iterates
straddle the fixed point, a cobweb that spirals inward. The figure
shows the cobweb for this case.

\begin{figure}[ht]
\centering
\begin{tikzpicture}[x=7cm, y=7cm, font=\small]

% Cobweb for a_{n+1} = cos(a_n), fixed point at ~0.739085 (Dottie number).
% Domain shown: a in [0.4, 1.05]. Map directly (x,y in same units).
% Axes
\draw[->, thick] (0.38, 0.38) -- (1.08, 0.38) node[right] {$a_{n}$};
\draw[->, thick] (0.38, 0.38) -- (0.38, 1.08) node[above] {$a_{n+1}$};

% ticks
\foreach \t in {0.5, 0.7, 0.9} {
  \draw (\t, 0.375) -- (\t, 0.385) node[below, yshift=-0.12cm] {\t};
  \draw (0.375, \t) -- (0.385, \t) node[left, xshift=-0.12cm] {\t};
}

% Line y = a
\draw[thick, engblack] (0.4, 0.4) -- (1.05, 1.05) node[right, pos=0.95] {$y=a$};

% Curve y = cos(a) over [0.4, 1.05]
\draw[thick, engblue] plot[domain=0.4:1.05, samples=60] (\x, {cos(\x r)})
  node[right, pos=0.0, yshift=0.25cm] {$y=\cos a$};

% Cobweb path: start a0=1.0
% iterates: 1.0 -> cos = 0.5403 ; 0.5403 -> 0.8576 ; 0.8576 -> 0.6543 ; 0.6543 -> 0.7935
\draw[thick, engred] (1.0, 0.4) -- (1.0, 0.5403);      % vertical to curve
\draw[thick, engred] (1.0, 0.5403) -- (0.5403, 0.5403); % horizontal to y=a
\draw[thick, engred] (0.5403, 0.5403) -- (0.5403, 0.8576);
\draw[thick, engred] (0.5403, 0.8576) -- (0.8576, 0.8576);
\draw[thick, engred] (0.8576, 0.8576) -- (0.8576, 0.6543);
\draw[thick, engred] (0.8576, 0.6543) -- (0.6543, 0.6543);
\draw[thick, engred] (0.6543, 0.6543) -- (0.6543, 0.7935);
\draw[thick, engred] (0.6543, 0.7935) -- (0.7935, 0.7935);

% fixed point marker
\fill[engblack] (0.739085, 0.739085) circle (0.4pt);
\node[anchor=south west] at (0.739085, 0.739085) {\footnotesize $a^{*}=0.7391$};

\node[engblue, anchor=north] at (0.5, 0.42) {\footnotesize $a_{0}=1$};

\end{tikzpicture}
\caption[Cobweb diagram for the iteration $a_{n+1}=\cos a_{n}$]{Cobweb
diagram for the fixed-point iteration $a_{n+1} = \cos a_{n}$, which
converges to the Dottie number $a^{*} = 0.739085$, the unique real
solution of $a = \cos a$. The iteration alternates above and below the
fixed point, the signature of a negative slope $g'(a^{*}) = -\sin a^{*}
\approx -0.674$ at the crossing. Because $\abs{g'(a^{*})} < 1$, the
cobweb spirals inward and the iteration converges linearly, gaining
about $0.17$ decimal digit per step.}
\label{fig:vol02:ch04:cobweb}
\end{figure}


The cobweb diagram of \autoref{fig:vol02:ch04:cobweb} reads the
convergence geometrically. The iteration steps vertically to the curve
$y = g(a)$, then horizontally to the line $y = a$, and repeats. A slope
of magnitude below $1$ at the fixed point pulls the cobweb inward; a
slope of magnitude above $1$ pushes it out, and the iteration diverges
no matter how close the start. The visual criterion ($\abs{\text{slope}}
< 1$ at the crossing) is the same $\abs{g'(a^{*})} < 1$ condition stated
above, made legible. The script
\texttt{code/fixed\_point\_cobweb.py} runs the iteration and prints the
empirical contraction ratio $e_{n+1}/e_{n}$, which settles to $-0.674$
as predicted; the run is tabulated in
\texttt{data/fixed\_point\_cobweb.csv}.

\subsection{The big-O and little-o notations}

The engineering notation for "this term is small in this regime" is
the \engterm{big-O} family. Given two sequences $(a_{n})$ and
$(b_{n})$ with $b_{n} > 0$ eventually:

\begin{itemize}[leftmargin=*]
  \item $a_{n} = O(b_{n})$ means $\abs{a_{n}} \leq C b_{n}$ for some
    constant $C$ and all sufficiently large $n$. The sequence $a_{n}$
    is bounded by a constant multiple of $b_{n}$.
  \item $a_{n} = o(b_{n})$ means $a_{n} / b_{n} \to 0$. The sequence
    $a_{n}$ is asymptotically negligible compared to $b_{n}$.
  \item $a_{n} \sim b_{n}$ means $a_{n} / b_{n} \to 1$. The two
    sequences are asymptotically equal.
\end{itemize}

The same notation applies to functions of a continuous variable, with
$x \to x_{0}$ or $x \to \infty$ replacing $n \to \infty$. Engineers use
these symbols to compress statements like "the truncation error is
proportional to $h^{2}$ as the step $h$ shrinks" into $\text{error} =
O(h^{2})$. Volume I taught the reader to carry a relative uncertainty
through every calculation; the big-O family is the formal home of that
habit when the smallness depends on a parameter.

\begin{archetype}[Scaling, formalised as expansion near a working point]
Volume I introduced scaling as the engineer's habit of asking how a
quantity changes when an input is multiplied by a factor. This chapter
sharpens the habit. When a function $f(x)$ is studied near a working
point $x_{0}$, its behaviour to leading order is captured by a small
finite number of terms in a Taylor series. The first-order term scales
linearly with the deviation $h = x - x_{0}$; the second-order term
scales as $h^{2}$; the residual scales as $h^{n+1}$ when $n$ terms are
kept. Scaling becomes the apparatus by which the engineer estimates
how good an approximation is and when to stop including more terms.
\end{archetype}

\begin{estimation}
\textbf{Question.} A numerical iteration produces residuals $r_{0} =
1.0$, $r_{1} = 0.5$, $r_{2} = 0.25$, $r_{3} = 0.125$. Estimate the
number of further iterations needed to reduce the residual to $10^{-6}$.

\textbf{Estimate.} The residuals halve at each step, so the sequence is
geometric with ratio $r = 0.5$. To reduce by a factor of $10^{6}$ from
the current $r_{3} = 0.125$, we need an additional factor of about
$10^{6} / 8 \approx 10^{5}$. Since $0.5^{17} \approx 7.6 \times 10^{-6}$
and $0.5^{20} \approx 9.5 \times 10^{-7}$, about $17$ to $20$ further
iterations should suffice. A working budget of $20$ iterations leaves
some margin.
\end{estimation}

\section{Series: geometric, harmonic, alternating}\pathtag{core}

Given a sequence $(a_{n})_{n \geq 1}$, the \engterm{partial sum} is
$S_{N} = \sum_{n=1}^{N} a_{n}$. The \engterm{series} is the limit of
the partial sums:
\[
\sum_{n=1}^{\infty} a_{n} = \lim_{N \to \infty} S_{N},
\]
when this limit exists. If it does, the series \engterm{converges}; if
not, it \engterm{diverges}.

A series is a sequence (the partial sums) that happens to be defined as
a sum. The convergence question for a series is the convergence question
for its sequence of partial sums. Three series carry most of the early
engineering content.

\subsection{The geometric series}

For $\abs{r} < 1$,
\[
\sum_{n=0}^{\infty} r^{n} = \frac{1}{1 - r}.
\]
The partial sum has a closed form. Multiply $S_{N} = 1 + r + r^{2} +
\cdots + r^{N}$ by $r$ and subtract: $(1 - r) S_{N} = 1 - r^{N+1}$, so
$S_{N} = (1 - r^{N+1}) / (1 - r)$. As $N \to \infty$, $r^{N+1} \to 0$
when $\abs{r} < 1$, giving the stated limit. For $\abs{r} \geq 1$, the
partial sums do not settle (for $r = 1$ they grow as $N + 1$; for
$r = -1$ they oscillate between $0$ and $1$; for $\abs{r} > 1$ they
grow without bound).

The geometric series is the workhorse of engineering. It models
discounted cash flow, signal repetition, multipath echoes, the
steady-state of a leaky integrator, the response of a recursive filter.
The closed form $1/(1 - r)$ also gives a useful first-order
approximation for any quantity of the form "a thing times a series of
diminishing copies of itself."

A finite geometric sum is also useful in its own right:
\[
\sum_{n=0}^{N-1} r^{n} = \frac{1 - r^{N}}{1 - r} \quad (r \neq 1).
\]
This formula handles cases where the engineer wants the cumulative
effect after $N$ steps, regardless of whether the series ultimately
converges.

A worked engineering use makes the geometric series concrete. A
relay sends a packet that may echo: the receiver sees the direct
arrival at unit amplitude, then a reflected copy at amplitude $\rho$,
then a copy of the copy at $\rho^{2}$, and so on, where $\rho$ is the
reflection coefficient of a mismatched termination. The total received
energy, normalised to the direct arrival, is $\sum_{n=0}^{\infty}
\rho^{2n} = 1/(1 - \rho^{2})$. For a $20\%$ reflection ($\rho = 0.2$),
the multipath energy is $1/(1 - 0.04) = 1.0417$, a $4.2\%$ excess over
the direct path. For a badly matched $\rho = 0.5$, the excess is
$1/(1 - 0.25) = 1.333$, a third more energy than the direct arrival,
spread across echoes that the receiver must equalise. The same closed
form gives the steady-state output of a first-order recursive
(leaky-integrator) filter $y_{n} = x_{n} + \rho y_{n-1}$ driven by a
unit step: the steady-state gain is $1/(1 - \rho)$, the geometric sum
of the impulse response. The reflection problem and the filter problem
are the same geometric series wearing different clothes.

A second worked use settles a present-value question. An engineering
project promises a fixed annual saving $C$ for $T$ years, and the
analyst discounts future money at rate $i$ per year, so a saving $C$
realised $n$ years out is worth $C/(1 + i)^{n}$ today. The present
value of the stream is a finite geometric sum with ratio $r = 1/(1 +
i)$:
\[
\text{PV} = \sum_{n=1}^{T} \frac{C}{(1 + i)^{n}} = C \, r
\frac{1 - r^{T}}{1 - r} = C \, \frac{1 - (1 + i)^{-T}}{i}.
\]
For $C = \$10{,}000$ per year, $i = 0.05$, and $T = 20$ years, the
factor $(1 - 1.05^{-20})/0.05 = (1 - 0.3769)/0.05 = 12.46$, so the
present value is about $\$124{,}600$, well below the undiscounted
$\$200{,}000$. Letting $T \to \infty$ gives the perpetuity value $C/i =
\$200{,}000$, the limit of the partial sums. The same closed form
prices a maintenance contract, a bond coupon stream, and the
lifetime-cost term in a total-cost-of-ownership comparison. The
geometric series is the engineer's discounting engine, and the
distinction between the finite sum and its infinite limit is the
distinction between a fixed-term and a perpetual obligation.

\subsection{Telescoping series}

A \engterm{telescoping series} has terms that cancel in consecutive
pairs. If $a_{n} = b_{n} - b_{n+1}$, the partial sum collapses:
\[
S_{N} = \sum_{n=1}^{N} (b_{n} - b_{n+1}) = b_{1} - b_{N+1},
\]
and the infinite sum is $b_{1} - \lim_{n} b_{n+1}$ when the limit
exists. The recognition skill is to spot a partial-fraction
decomposition that produces the cancelling form. The canonical case is
\[
\sum_{n=1}^{\infty} \frac{1}{n(n+1)} = \sum_{n=1}^{\infty}
\left(\frac{1}{n} - \frac{1}{n+1}\right) = 1 - \lim_{N \to \infty}
\frac{1}{N+1} = 1.
\]
A less obvious case rewards the same trick. To sum $\sum 1/(n(n+2))$,
decompose $1/(n(n+2)) = \tfrac{1}{2}(1/n - 1/(n+2))$. The cancellation
now skips a term, leaving the first two unmatched:
\[
\sum_{n=1}^{\infty} \frac{1}{n(n+2)} = \frac{1}{2}\left(1 + \frac{1}{2}
- \lim_{N}\left(\frac{1}{N+1} + \frac{1}{N+2}\right)\right) = \frac{3}{4}.
\]
Telescoping is the one method in the chapter that returns an exact
closed-form sum rather than a convergence verdict, which is why it is
worth the recognition effort. Many sums that look intractable reduce to
telescoping after a partial-fraction step, including the variance of a
discrete-uniform distribution and several lattice-sum identities in
later volumes.

\subsection{The harmonic series}

For all the simplicity of its terms, the series
\[
\sum_{n=1}^{\infty} \frac{1}{n} = 1 + \frac{1}{2} + \frac{1}{3} +
\frac{1}{4} + \cdots
\]
diverges. The standard proof by Oresme groups terms:
\[
1 + \frac{1}{2} + \underbrace{\left(\frac{1}{3} + \frac{1}{4}\right)}_{>
1/2} + \underbrace{\left(\frac{1}{5} + \cdots + \frac{1}{8}\right)}_{>
1/2} + \cdots,
\]
each parenthesised group exceeds $1/2$, and there are infinitely many
groups, so the partial sums grow without bound. The growth is slow:
$S_{N} \approx \log N + \gamma$, where $\gamma \approx 0.5772$ is the
Euler-Mascheroni constant. To reach $S_{N} = 100$ takes roughly $N
\approx e^{100 - \gamma} \approx 1.5 \times 10^{43}$ terms.

The harmonic series is the canonical case of a divergent series whose
terms shrink to zero. The terms shrinking is necessary for convergence
but not sufficient. We return to this case in section 4.6 as the
chapter's failure mode.

\subsection{Alternating series}

A series of the form $\sum_{n=1}^{\infty} (-1)^{n+1} b_{n}$ with
$b_{n} > 0$ is \engterm{alternating}. The \engterm{alternating series
test} (Leibniz) states: if $b_{n}$ is monotone decreasing and
$b_{n} \to 0$, the series converges. The classical example is the
alternating harmonic series:
\[
\sum_{n=1}^{\infty} \frac{(-1)^{n+1}}{n} = 1 - \frac{1}{2} + \frac{1}{3}
- \frac{1}{4} + \cdots = \log 2.
\]
The cancellation between consecutive terms produces a convergent series
where the unsigned series diverged. The engineering interpretation: a
process that adds and subtracts alternately, with each contribution
slightly smaller than the last, has a finite cumulative effect even
when the unsigned process does not.

The alternating series test also gives a useful error bound. If $S_{N}$
is the $N$-th partial sum of a convergent alternating series and $S$ is
the true sum, then $\abs{S - S_{N}} \leq b_{N+1}$: the truncation error
is at most the magnitude of the first omitted term. This is the
simplest worked-bound formula in the chapter.

\begin{figure}[ht]
\centering
\begin{tikzpicture}[x=1.1cm, y=11cm, font=\small]

% Partial sums of alternating harmonic series converging to log 2 = 0.6931
% S1=1.0000, S2=0.5000, S3=0.8333, S4=0.5833, S5=0.7833, S6=0.6167,
% S7=0.7595, S8=0.6345, S9=0.7456, S10=0.6456
% y range ~ [0.45, 1.05]; subtract 0.45 baseline, scale handled by y unit.
\def\base{0.45}

% Axes
\draw[->, thick] (0.5, 0) -- (10.7, 0) node[right] {$N$};
\draw[->, thick] (0.5, 0) -- (0.5, 0.62) node[above] {$S_{N}$};

% y ticks at 0.5, 0.6931, 0.8, 1.0 (shifted by base)
\foreach \v/\lbl in {0.5/0.5, 0.6931/$\log 2$, 0.8/0.8, 1.0/1.0} {
  \pgfmathsetmacro\yy{\v-\base}
  \draw (0.45, \yy) -- (0.55, \yy) node[left, xshift=-0.1cm] {\lbl};
}
% x ticks
\foreach \n in {1,2,...,10} {
  \draw (\n, -0.005) -- (\n, 0.005);
}
\node[below] at (2, -0.01) {2};
\node[below] at (4, -0.01) {4};
\node[below] at (6, -0.01) {6};
\node[below] at (8, -0.01) {8};
\node[below] at (10, -0.01) {10};

% limit line
\draw[dashed, gray] (0.5, {0.6931-\base}) -- (10.6, {0.6931-\base});

% odd partial sums (approach from above) - red, decreasing
\draw[thick, engred, mark=*, mark size=1pt]
  plot coordinates {
  (1, {1.0000-\base}) (3, {0.8333-\base}) (5, {0.7833-\base})
  (7, {0.7595-\base}) (9, {0.7456-\base}) };
\node[engred, anchor=west] at (9.1, {0.7456-\base}) {odd $S_N$};

% even partial sums (approach from below) - blue, increasing
\draw[thick, engblue, mark=square*, mark size=1pt]
  plot coordinates {
  (2, {0.5000-\base}) (4, {0.5833-\base}) (6, {0.6167-\base})
  (8, {0.6345-\base}) (10, {0.6456-\base}) };
\node[engblue, anchor=west] at (10.1, {0.6456-\base}) {even $S_N$};

\end{tikzpicture}
\caption[Nested-interval envelope of alternating partial sums]{Partial
sums of the alternating harmonic series $1 - 1/2 + 1/3 - \cdots$
converging to $\log 2$. The odd partial sums descend from above and the
even partial sums climb from below, trapping the limit in a nested
sequence of intervals $[S_{2k}, S_{2k+1}]$ whose width is the next term
$b_{2k+1}$. The width is the alternating-series error bound made
visible: any partial sum lies within one term of the limit.}
\label{fig:vol02:ch04:alternating-envelope}
\end{figure}


The envelope picture in \autoref{fig:vol02:ch04:alternating-envelope}
shows why the bound holds. The partial sums of a convergent alternating
series straddle the limit: the odd partial sums approach from above, the
even partial sums from below, and the true sum is trapped in the
shrinking interval $[S_{2k}, S_{2k+1}]$. The width of the trapping
interval is exactly $b_{2k+1}$, the next term, so the distance from any
partial sum to the limit cannot exceed the next term. The bound is not
merely an inequality the engineer must trust; it is the visible width of
a nested-interval cage. A worked instance: for $\log 2 = 1 - 1/2 + 1/3 -
\cdots$, after the term $1/3$ the partial sum is $S_{3} = 0.8333$ and the
bound on the residual is $b_{4} = 1/4 = 0.25$. The true residual is
$\abs{0.6931 - 0.8333} = 0.1402$, within the bound, and on the upper
side as the odd-index rule predicts.

\subsection{Absolute and conditional convergence}

A series $\sum a_{n}$ \engterm{converges absolutely} if
$\sum \abs{a_{n}}$ converges. It \engterm{converges conditionally} if
$\sum a_{n}$ converges but $\sum \abs{a_{n}}$ does not.

Absolute convergence is the safer condition. A rearrangement of an
absolutely convergent series sums to the same total. A rearrangement
of a conditionally convergent series can be made to sum to any real
number (Riemann's rearrangement theorem). For engineers, the practical
content is: when manipulating a series (reordering, regrouping,
multiplying by another series), check that the operations preserve the
sum. Absolute convergence preserves sum under reordering; conditional
convergence does not.

The Riemann rearrangement theorem is not an abstraction without bite.
The alternating harmonic series sums to $\log 2 \approx 0.693$ in its
natural order. Take two positive terms for every one negative term,
$1 + 1/3 - 1/2 + 1/5 + 1/7 - 1/4 + \cdots$, and the same terms, every
one of them used exactly once, sum to $\tfrac{3}{2}\log 2 \approx
1.040$. A numerical routine that accumulates a conditionally convergent
series in an order set by cache locality or by parallel reduction rather
than by the mathematical index can land on a different total, and the
difference is not round-off. The defensive habit is to sum in magnitude
order, or to confirm absolute convergence before reordering for speed.

\begin{mastery}
\engterm{Summability of divergent series}. A divergent series can
still be assigned a finite value by a \engterm{summation method} that
agrees with ordinary summation on convergent series. \engterm{Ces\`{a}ro
summation} averages the partial sums: $\sigma_{N} = (S_{1} + \cdots +
S_{N})/N$, and a series is Ces\`{a}ro-summable to $\sigma$ if
$\sigma_{N} \to \sigma$. The series $1 - 1 + 1 - 1 + \cdots$ has
partial sums $1, 0, 1, 0, \ldots$ that do not converge, but the
Ces\`{a}ro averages converge to $1/2$. \engterm{Abel summation} assigns
$\lim_{x \to 1^{-}} \sum a_{n} x^{n}$ when the limit exists. These
methods are not mathematical curiosities for the engineer; Ces\`{a}ro
and Abel means appear in the convergence analysis of Fourier series
(where the Ces\`{a}ro mean of a Fourier series, the Fej\'{e}r kernel,
converges uniformly for a continuous function while the raw series may
not) and in signal-processing windowing. The hazard is to read a
summation method's finite answer as the value of the divergent series in
the ordinary sense; the regularised value is a different object, valid
only for the operations the method preserves.
\end{mastery}

\begin{figure}[ht]
\centering
\begin{tikzpicture}[x=1cm, y=1cm, font=\small]

% Axes: x = log10(N), y = partial sum value
\draw[->, thick] (-0.2, 0) -- (10.5, 0) node[right] {$N$};
\draw[->, thick] (0, -0.2) -- (0, 5.5) node[above] {$S_N$};

% x-axis ticks: log10 N from 1 to 6
\foreach \x/\lbl in {1/$10$, 2/$10^{2}$, 3/$10^{3}$, 4/$10^{4}$, 5/$10^{5}$, 6/$10^{6}$, 7/$10^{7}$, 8/$10^{8}$, 9/$10^{9}$} {
  \draw (\x, -0.1) -- (\x, 0.1) node[below, yshift=-0.15cm] {\lbl};
}
% y-axis ticks
\foreach \y in {1, 2, 3, 4, 5} {
  \draw (-0.1, \y) -- (0.1, \y) node[left, xshift=-0.15cm] {\y};
}

% Harmonic series: S_N ~ ln N + gamma. Approximate at decades.
% N=10: 2.929; N=100: 5.187; would exceed plot; so scale by 1/3
% Let's plot S_N/4 to fit; relabel y-axis accordingly
% Simpler: plot the geometric-series partial sum sum 0.5^n = 2 - 2*0.5^N -> converges to 2
% and alternating harmonic -> log 2 = 0.693
% and harmonic -> diverges

% Geometric r=0.5 sum: 2(1 - 0.5^N), bound 2
\draw[thick, engblue] plot coordinates {
  (0.3,1.0) (1,1.998) (2,2.0) (3,2.0) (4,2.0) (5,2.0) (6,2.0) (7,2.0) (8,2.0) (9,2.0)
};
\draw[dashed, engblue] (0, 2) -- (9.5, 2) node[right, engblue] {$=2$};
\node[engblue, anchor=west] at (4, 2.25) {geometric $r{=}0.5$};

% Alternating harmonic: -> log 2 = 0.693
\draw[thick, engred] plot coordinates {
  (0.3,0.583) (1,0.646) (2,0.688) (3,0.6927) (4,0.6931) (5,0.6931) (6,0.6931) (9,0.6931)
};
\draw[dashed, engred] (0, 0.693) -- (9.5, 0.693) node[right, engred] {$=\log 2$};
\node[engred, anchor=west] at (5.5, 0.4) {alternating harmonic};

% Harmonic divergence: log N + gamma; scale: plot value/4 so visible. Show through y=5
\draw[thick, engblack] plot coordinates {
  (0.3,0.155) (1,0.733) (2,1.297) (3,1.872) (4,2.448) (5,3.024) (6,3.601) (7,4.177) (8,4.753) (9,5.329)
};
\node[engblack, anchor=west] at (6.5, 4.5) {harmonic};
\node[engblack, anchor=west, font=\scriptsize] at (6.5, 4.1) {(divergent)};

\end{tikzpicture}
\caption[Partial sums of three canonical series]{Partial sums on a
semi-log $N$ axis for three series with shrinking terms. The geometric
series with $r = 0.5$ settles to $2$ after a handful of terms; the
alternating harmonic series settles to $\log 2$ after about a hundred;
the harmonic series grows as $\log N + \gamma$ and never settles. All
three have $a_n \to 0$, yet only two converge. The visual lesson
matches the section 4.6 principle: shrinking terms are necessary but
not sufficient for convergence. The harmonic curve passes $5$ at
$N \approx 10^{9}$.}
\label{fig:vol02:ch04:partial-sums}
\end{figure}


The three partial-sum trajectories in \autoref{fig:vol02:ch04:partial-sums}
sit alongside each other as the chapter's working visual reminder.
Two converge; one diverges; all three have $a_n \to 0$. The geometric
series settles in a handful of terms because its partial sums form a
contraction with ratio $r$. The alternating harmonic settles because
each step is bounded by the next omitted term, and that magnitude
shrinks to zero. The harmonic series grows as $\log N + \gamma$ at
a rate the eye reads as flat across any single screen of the plot,
and the engineer who trusts the visual without checking the analytic
bound truncates at a finite $N$ believing the tail is negligible. It
is not, and the chapter's failure section returns to this case.

\subsection{A worked Fourier-series partial sum}

The series engineers meet most often is not a series of numbers but a
series of functions: a Fourier series, in which the partial sum is a
trigonometric polynomial that approximates a periodic signal. The
convergence question for such a series is richer than for a number
series, because the partial sum can converge at one point at one rate
and at a neighbouring point at another, and because a discontinuity in
the target produces a permanent overshoot that no number of terms
removes. The square wave is the canonical worked case, and it pays the
chapter's machinery back in a single figure.

The square wave of unit amplitude and period $2\pi$, odd about the
origin, has the Fourier series
\[
s(x) = \frac{4}{\pi} \sum_{k=0}^{\infty} \frac{\sin\bigl((2k+1) x\bigr)}
{2k+1} = \frac{4}{\pi}\left(\sin x + \frac{\sin 3x}{3} +
\frac{\sin 5x}{5} + \cdots\right).
\]
The coefficients decay as $1/(2k+1)$, the harmonic rate, so the series
of coefficients on its own diverges. Convergence of the function series
is rescued by the cancellation between the oscillating sines, exactly
the mechanism that made the alternating harmonic series converge. At a
point of continuity the partial sum $S_{N}(x)$, keeping harmonics
through order $2N-1$, converges to $s(x)$, but the convergence is slow:
the pointwise error at a fixed interior point decays as $O(1/N)$,
inherited from the coefficient tail $\sum_{k > N} 1/(2k+1)$, which the
integral test bounds by $\tfrac{1}{2}\log(N/(N+1))^{-1} \approx
1/(4N)$.

The interesting failure is at the jump. The square wave steps from
$-1$ to $+1$ at $x = 0$. Near the jump the partial sum overshoots, and
the overshoot does not shrink as $N$ grows. It moves closer to the jump
and narrows, but its height settles at a fixed fraction above the true
value. The overshoot peaks at $x_{N} \approx \pi/(2N)$, the first
maximum of $S_{N}$ to the right of the jump, and its limiting height is
\[
\lim_{N \to \infty} S_{N}(x_{N}) = \frac{2}{\pi} \int_{0}^{\pi}
\frac{\sin t}{t} \, \dd t = \frac{2}{\pi}\,\mathrm{Si}(\pi) \approx
1.1790,
\]
an overshoot of about $9\%$ of the unit jump, where $\mathrm{Si}$ is
the sine integral. This is the \engterm{Gibbs phenomenon}. The
$9\%$ figure is independent of $N$ and of the particular signal; every
jump discontinuity reconstructed by a truncated Fourier series carries
the same proportional overshoot. The integral arises because the sum
$\sum \sin((2k+1)x_{N})/(2k+1)$ at the overshoot point is a Riemann sum
for $\int_{0}^{\pi} (\sin t)/t \, \dd t$, and $\mathrm{Si}(\pi) =
1.8519$ exceeds the value $\pi/2 = 1.5708$ that the converged series
would give.

\begin{figure}[ht]
\centering
\begin{tikzpicture}[x=1cm, y=1cm, font=\small]

% Square-wave Fourier partial sums S_N near the jump at x=0.
% Horizontal axis x in [-1.6, 1.6] mapped to cm by factor 2.4 (so x=1.6 -> 3.84 cm).
% Vertical axis amplitude in [-1.3, 1.3] mapped 1:1 (cm = amplitude * 2.0).
% target square wave: -1 for x<0, +1 for x>0.

% Axes
\draw[->, thick] (-4.2, 0) -- (4.6, 0) node[right] {$x$};
\draw[->, thick] (0, -2.9) -- (0, 2.9) node[above] {amplitude};

% amplitude ticks at +1, -1, and the Gibbs level 1.179
\foreach \a/\lbl in {2.0/$1$, -2.0/$-1$} {
  \draw (-0.1, \a) -- (0.1, \a) node[left, xshift=-0.1cm] {\lbl};
}
\draw[dotted, gray] (-4.0, 2.358) -- (4.0, 2.358)
  node[right, gray, yshift=0.15cm] {$1.179$};
\draw[dotted, gray] (-4.0, 2.0) -- (4.0, 2.0);

% true square wave (target): step at x=0
\draw[very thick, engblack] (-3.84, -2.0) -- (-0.02, -2.0);
\draw[very thick, engblack] (0.02, 2.0) -- (3.84, 2.0);
\node[engblack, anchor=west] at (2.6, 2.55) {target};

% S_5 (N=3 harmonics: sin x + sin3x/3 + sin5x/5, scaled 4/pi)
% sampled values amplitude*2.0; coordinates (x*2.4, amp*2.0). Right half only mirror.
% Precomputed S(x) = (4/pi)(sin x + sin3x/3 + sin5x/5) for selected x.
\draw[thick, engblue] plot[smooth] coordinates {
  (0,0) (0.24,0.55) (0.48,1.04) (0.72,1.45) (0.96,1.92)
  (1.2,2.34) (1.5,2.36) (1.92,1.78) (2.4,1.84) (2.88,2.36)
  (3.36,2.10) (3.84,1.72) };
\node[engblue, anchor=west] at (3.9, 1.72) {$N=3$};

% S_19 (many harmonics): tighter ring, overshoot near jump reaching ~1.179.
\draw[thick, engred] plot[smooth] coordinates {
  (0,0) (0.12,1.20) (0.24,2.10) (0.34,2.358) (0.5,2.10)
  (0.7,2.18) (1.0,1.92) (1.4,2.06) (2.0,1.98) (2.8,2.02)
  (3.6,2.00) (3.84,2.00) };
\node[engred, anchor=west] at (3.9, 2.10) {$N=19$};

% mirror the right-half curves to the left (odd symmetry): handled visually by
% labelling; for compactness only the right half and target are drawn in detail.

\end{tikzpicture}
\caption[Gibbs overshoot in a truncated Fourier reconstruction]{Partial
sums of the unit square wave, keeping harmonics through order $2N-1$, near
the jump at $x = 0$. The low-order sum ($N = 3$) rings broadly; the
high-order sum ($N = 19$) rings narrowly but overshoots to the same
limiting height $\tfrac{2}{\pi}\mathrm{Si}(\pi) \approx 1.179$, about
$9\%$ above the unit step. Adding harmonics moves the overshoot toward the
jump and narrows it without lowering its peak. The fixed overshoot is the
Gibbs phenomenon, a failure of uniform convergence at the discontinuity;
the partial sums converge pointwise to the square wave everywhere except
the jump, yet the worst-case error never falls below about $0.18$.}
\label{fig:vol02:ch04:gibbs}
\end{figure}


The Gibbs picture in \autoref{fig:vol02:ch04:gibbs} is the engineering
warning made visible. A truncated Fourier reconstruction of a signal
with a step (a digital edge, a switching transient, a control setpoint
change) rings near the edge with a fixed-height overshoot, and an
engineer who tightens the reconstruction by adding harmonics narrows
the ringing without lowering its peak. The defensive habits are to
window the series (a Fej\'{e}r or Lanczos sigma factor multiplies the
coefficients by a taper that suppresses the overshoot at the cost of a
gentler edge), to accept the $9\%$ as a hard floor when sharp edges
must be preserved, or to represent the discontinuity with a basis that
admits jumps. The convergence here is uniform on any closed interval
that excludes the jump and is never uniform across the jump, the
distinction the next mastery box sharpens. The classical analysis of
the phenomenon and of the windowing remedies is in
\cite{hist:fourier1822,text:v2ch03-stein-shakarchi-complex}.

\begin{mastery}
\engterm{Uniform versus pointwise convergence}. A sequence of functions
$f_{N}$ converges \engterm{pointwise} to $f$ if $f_{N}(x) \to f(x)$ for
each fixed $x$. It converges \engterm{uniformly} if the worst-case error
$\sup_{x} \abs{f_{N}(x) - f(x)}$ tends to zero, so a single $N$ serves
every $x$ at once. The distinction is the difference between "for each
$x$, eventually close" and "eventually close everywhere at once." The
Gibbs overshoot is precisely a failure of uniform convergence: the
partial sums of the square wave converge pointwise to $s(x)$ at every
$x \neq 0$, yet $\sup_{x} \abs{S_{N}(x) - s(x)} \approx 0.18$ for all
$N$, because the overshoot near the jump never shrinks. Uniform
convergence is the property the engineer needs to interchange limit
with integral or with derivative, to bound a function-series tail by a
single number, and to guarantee that the limit of continuous partial
sums is itself continuous. Pointwise convergence guarantees none of
these. The Weierstrass $M$-test certifies uniform convergence of a
function series $\sum g_{n}(x)$ when $\abs{g_{n}(x)} \leq M_{n}$ for all
$x$ and the number series $\sum M_{n}$ converges. The square wave fails
the $M$-test because its $M_{n} = 4/(\pi(2n-1))$ sum to the divergent
harmonic series, which is the analytic signature of the Gibbs trouble.
\end{mastery}

\section{Convergence tests engineers actually use}\pathtag{standard}

The following tests are the working tools. We state each in its
operational form, with a brief justification, and indicate the regime
where it is sharpest.

\subsection{Term test (necessary condition)}

If $\sum a_{n}$ converges, then $a_{n} \to 0$. Equivalently, if
$a_{n} \not\to 0$, the series diverges. This is the cheapest test and
catches the simplest failures: the partial sums cannot settle if the
last term added is not shrinking. The converse fails (the harmonic
series has $a_{n} \to 0$ but diverges); the term test rejects, it does
not approve.

\subsection{Comparison test}

If $0 \leq a_{n} \leq b_{n}$ for all sufficiently large $n$ and
$\sum b_{n}$ converges, then $\sum a_{n}$ converges. If $0 \leq b_{n}
\leq a_{n}$ and $\sum b_{n}$ diverges, then $\sum a_{n}$ diverges.

The working version compares the unknown series to a known series. The
catalog of known series carries the engineer's library: $\sum 1/n^{p}$
converges for $p > 1$ (the \engterm{p-series}), diverges for $p \leq 1$;
$\sum r^{n}$ converges for $\abs{r} < 1$. To test $\sum 1/(n^{2} + 1)$,
compare with $\sum 1/n^{2}$, which converges, so the original series
converges. To test $\sum 1/(2n + 1)$, compare with $\sum 1/(2n + 2n) =
(1/4) \sum 1/n$, which diverges, so the original diverges.

A useful refinement is the \engterm{limit comparison test}: if
$a_{n}, b_{n} > 0$ and $\lim_{n \to \infty} a_{n}/b_{n} = c$ for some
$0 < c < \infty$, then $\sum a_{n}$ and $\sum b_{n}$ either both
converge or both diverge. This handles cases where the inequality
$a_{n} \leq b_{n}$ holds only up to a constant factor.

\subsection{Ratio test}

For a series $\sum a_{n}$ with $a_{n} > 0$, compute
\[
L = \lim_{n \to \infty} \frac{a_{n+1}}{a_{n}}.
\]
The series converges if $L < 1$ and diverges if $L > 1$. The case
$L = 1$ is inconclusive.

The ratio test is the engineer's favourite for series whose terms
involve factorials or exponentials. The series $\sum n! / n^{n}$ has
ratio $a_{n+1}/a_{n} = (n+1)/((n+1)/n)^{n+1} \cdot 1/(n+1) = (n/(n+1))^{n}
\to 1/e < 1$, so the series converges. The series $\sum c^{n}/n!$ has
ratio $c/(n+1) \to 0$, so it converges absolutely for every $c$, which
underwrites the convergence of $e^{x}$ everywhere.

\subsection{Root test}

For a series $\sum a_{n}$ with $a_{n} \geq 0$, compute
\[
L = \lim_{n \to \infty} \sqrt[n]{a_{n}}.
\]
The series converges if $L < 1$ and diverges if $L > 1$. The case
$L = 1$ is inconclusive.

The root test is sharper than the ratio test in cases where the terms
involve $n$-th powers, such as $\sum (1 - 1/n)^{n^{2}}$. Here
$\sqrt[n]{a_{n}} = (1 - 1/n)^{n} \to 1/e$, so the series converges.
Where both tests apply and give a clean limit, the ratio test is
usually easier to compute; where the ratio test produces an
indeterminate ratio, the root test sometimes still gives a clean
$n$-th-root limit.

\subsection{Integral test (informal)}

If $f$ is positive, decreasing, and continuous on $[1, \infty)$ with
$f(n) = a_{n}$, then $\sum_{n=1}^{\infty} a_{n}$ converges if and only
if the improper integral $\int_{1}^{\infty} f(x) \, \dd x$ converges.
We will define this integral formally in Chapter 6; here we use the
test informally to motivate two key results.

The first is the p-series test. For $f(x) = 1/x^{p}$, the integral
$\int_{1}^{\infty} 1/x^{p} \, \dd x$ converges for $p > 1$ (giving
$1/(p-1)$) and diverges for $p \leq 1$. The series $\sum 1/n^{p}$
follows the same pattern.

The second is the harmonic divergence at the level of growth rate.
For $p = 1$, $\int_{1}^{N} 1/x \, \dd x = \log N$. The harmonic
partial sums grow as $\log N + \gamma + O(1/N)$, the integral test's
verdict in concrete form.

The integral test also delivers a truncation bound the engineer can use
directly. For a convergent $p$-series $\sum 1/n^{p}$ with $p > 1$, the
tail after $N$ terms is bounded by the corresponding integral:
\[
\sum_{n=N+1}^{\infty} \frac{1}{n^{p}} \leq \int_{N}^{\infty}
\frac{\dd x}{x^{p}} = \frac{N^{1-p}}{p - 1}.
\]
For $p = 2$, the tail after $N$ terms is at most $1/N$, so summing
$\sum 1/n^{2} = \pi^{2}/6 \approx 1.6449$ to within $10^{-3}$ takes
about $N = 1000$ terms by the crude bound, and the sharper midpoint
bound $\sum_{n > N} 1/n^{2} \approx 1/N - 1/(2N^{2})$ confirms the
order. The slow $1/N$ tail decay is why direct summation of $\sum
1/n^{2}$ is a poor way to compute $\pi^{2}/6$; an engineer who needs
the constant uses an accelerated form rather than the naive series. The
integral test thus does double duty: it settles convergence and it
bounds the truncation tail, the two questions the engineer actually
asks of any series.

\subsection{Picking a test}

A working order:

\begin{enumerate}[leftmargin=*]
  \item Try the term test. If $a_{n} \not\to 0$, the series diverges.
  \item If the term involves $r^{n}$, use the geometric series.
  \item If the term involves a factorial, use the ratio test.
  \item If the term is $\sim 1/n^{p}$ for large $n$, use the comparison
    or limit comparison test against $\sum 1/n^{p}$.
  \item If the terms are positive and the term is the value of a
    function $f(n)$ for which $\int f$ has a clean form, use the
    integral test.
  \item If the series alternates and the magnitude is monotone
    decreasing to zero, the alternating-series test gives convergence
    and a one-term error bound.
  \item If none of the above applies, the root test sometimes
    succeeds where the ratio test produces $L = 1$.
\end{enumerate}

The list is not exhaustive, but it covers the engineering catalog. A
reader who runs through it on a new series, in order, settles the
convergence question for nearly every series this volume produces
\cite{text:strang2016,text:v2ch04-stewart-calculus-2015}.

\begin{figure}[p]
\centering
\begin{tikzpicture}[
  x=1cm, y=1cm, font=\small,
  node distance=0.5cm,
  box/.style={draw, thick, rectangle, rounded corners=2pt, minimum width=4cm, minimum height=0.8cm, align=center, inner sep=4pt},
  decision/.style={draw, thick, diamond, aspect=2, minimum width=4cm, align=center, inner sep=1pt},
  leaf/.style={draw, thick, rectangle, fill=engblue!10, rounded corners=2pt, minimum width=3.5cm, minimum height=0.8cm, align=center, inner sep=4pt},
  fail/.style={draw, thick, rectangle, fill=engred!10, rounded corners=2pt, minimum width=3.5cm, minimum height=0.8cm, align=center, inner sep=4pt},
  arrow/.style={->, thick, >=stealth}
]

% Start
\node[box] (start) at (0, 14) {Given $\sum a_n$};

% Term test
\node[decision] (term) at (0, 12.5) {$a_n \to 0$?};
\draw[arrow] (start) -- (term);

\node[fail] (diverges1) at (-5, 12.5) {Diverges \\ (term test)};
\draw[arrow] (term) -- node[above, font=\scriptsize] {no} (diverges1);

% Form check: r^n?
\node[decision] (geom) at (0, 11) {form $r^n$?};
\draw[arrow] (term) -- node[right, font=\scriptsize] {yes} (geom);

\node[leaf] (geomtest) at (5, 11) {Geometric: \\ converges iff $|r| < 1$};
\draw[arrow] (geom) -- node[above, font=\scriptsize] {yes} (geomtest);

% Factorial?
\node[decision] (fact) at (0, 9.5) {factorial?};
\draw[arrow] (geom) -- node[right, font=\scriptsize] {no} (fact);

\node[leaf] (ratio) at (5, 9.5) {Ratio test};
\draw[arrow] (fact) -- node[above, font=\scriptsize] {yes} (ratio);

% Alternates?
\node[decision] (alt) at (0, 8) {alternates?};
\draw[arrow] (fact) -- node[right, font=\scriptsize] {no} (alt);

\node[leaf] (altseries) at (5, 8) {Alternating-series test \\ ($b_n \downarrow 0$ gives convergence)};
\draw[arrow] (alt) -- node[above, font=\scriptsize] {yes} (altseries);

% 1/n^p like?
\node[decision] (poly) at (0, 6.5) {$\sim 1/n^p$ large $n$?};
\draw[arrow] (alt) -- node[right, font=\scriptsize] {no} (poly);

\node[leaf] (polytest) at (5, 6.5) {(Limit) comparison \\ with $\sum 1/n^p$};
\draw[arrow] (poly) -- node[above, font=\scriptsize] {yes} (polytest);

% Integrable f(n)?
\node[decision] (intf) at (0, 5) {clean $\int f$ exists?};
\draw[arrow] (poly) -- node[right, font=\scriptsize] {no} (intf);

\node[leaf] (inttest) at (5, 5) {Integral test};
\draw[arrow] (intf) -- node[above, font=\scriptsize] {yes} (inttest);

% Root test
\node[decision] (rootq) at (0, 3.5) {$n$-th-root limit?};
\draw[arrow] (intf) -- node[right, font=\scriptsize] {no} (rootq);

\node[leaf] (rootest) at (5, 3.5) {Root test};
\draw[arrow] (rootq) -- node[above, font=\scriptsize] {yes} (rootest);

% Inconclusive
\node[fail] (incon) at (0, 2) {Inconclusive: \\ split into pieces, \\ or look up};
\draw[arrow] (rootq) -- node[right, font=\scriptsize] {no} (incon);

\end{tikzpicture}
\caption[Convergence-test decision tree]{The convergence-test ladder
as a decision tree. The engineer descends from cheapest test to most
specialised. The term test catches the simplest divergent series; the
geometric form is recognised on sight; the ratio test handles factorials;
the alternating-series test handles oscillating terms with monotone
magnitudes; comparison handles polynomial-tail series; the integral test
handles cases where $f(n)$ has a clean antiderivative; the root test is
the back-stop. The leaves at the right are the tests; the diamonds are
the recognition conditions. Engineers run the tree in order until a leaf
returns a verdict.}
\label{fig:vol02:ch04:test-decision-tree}
\end{figure}


The decision tree of \autoref{fig:vol02:ch04:test-decision-tree} arranges
the same ladder visually. The reader descends from cheap to expensive,
takes the first test that returns a verdict, and stops. The script
\texttt{code/series\_test\_picker.py} walks the same tree on a small
catalog of series and prints the test that settles each, as a check on
the reader's own walk through the tree.

\subsection{Common patterns to recognise on sight}

A working engineer should recognise the following five families
without going past the term test:

\begin{itemize}[leftmargin=*]
  \item \engterm{Geometric tail}, $\sum a r^{n}$: converges to
    $a/(1 - r)$ for $|r| < 1$. The sum of a discounted-cash-flow
    annuity, the steady-state output of an exponentially decaying
    impulse response, and the total energy in a stable resonator
    fall in this family.
  \item \engterm{$p$-series tail}, $\sum 1/n^{p}$: converges for
    $p > 1$, diverges for $p \leq 1$. Sums of squared residuals
    (when residuals decay polynomially), tail sums in number theory,
    and the convergence of zeta values fall here.
  \item \engterm{Telescoping}, $\sum (b_{n} - b_{n+1})$: equals
    $b_{1} - \lim_{n} b_{n}$ when the limit exists. Many partial-
    fraction decompositions reduce to a telescoping sum, including
    $\sum 1/(n(n+1)) = 1$.
  \item \engterm{Power-of-$n$ over factorial}, $\sum n^{k}/n!$ or
    $\sum c^{n}/n!$: converges for every $c$ and every $k$ by the
    ratio test, and converges quickly because the ratio falls to
    zero as $n$ grows.
  \item \engterm{Asymptotic-to-$\log$}, $S_{N} \sim c \log N$: the
    series diverges, and the engineer can lower-bound the truncation
    tail by the gap to the partial sum at any chosen $N$. The
    harmonic series and series of the form $\sum 1/(n \log n)$
    belong here.
\end{itemize}

The catalog above covers approximately ninety percent of the series
the engineer meets in this volume. The remaining ten percent are
handled by the ladder.

\subsection{Five worked convergence verdicts}

The tests earn their place by application. We settle five series, each
chosen to exercise a different rung of the ladder, and state the test
that does the work.

\engterm{Series 1}: $\displaystyle \sum_{n=1}^{\infty}
\frac{n^{2} + 3}{2 n^{4} - n + 1}$. The term behaves like $n^{2}/(2
n^{4}) = 1/(2 n^{2})$ for large $n$. Limit comparison with $b_{n} =
1/n^{2}$: the ratio $a_{n}/b_{n} = n^{2}(n^{2} + 3)/(2 n^{4} - n + 1)
\to 1/2$, a finite positive constant. Since $\sum 1/n^{2}$ converges
(a $p$-series with $p = 2 > 1$), the series converges.

\engterm{Series 2}: $\displaystyle \sum_{n=1}^{\infty}
\frac{1}{\sqrt{n^{2} + 1}}$. The term behaves like $1/n$. Limit
comparison with $b_{n} = 1/n$: the ratio $n/\sqrt{n^{2} + 1} \to 1$.
Since $\sum 1/n$ diverges, the series diverges. The deceptive
square-root masks a $p = 1$ tail.

\engterm{Series 3}: $\displaystyle \sum_{n=1}^{\infty}
\frac{3^{n}}{n!}$. Factorial in the denominator points to the ratio
test: $a_{n+1}/a_{n} = 3^{n+1}/(n+1)! \cdot n!/3^{n} = 3/(n+1) \to 0 <
1$. The series converges absolutely, and quickly. The sum is $e^{3} - 1
\approx 19.09$, by the exponential series.

\engterm{Series 4}: $\displaystyle \sum_{n=1}^{\infty}
\left(\frac{2n + 1}{3n + 2}\right)^{n}$. The $n$-th power points to the
root test: $\sqrt[n]{a_{n}} = (2n + 1)/(3n + 2) \to 2/3 < 1$. The
series converges. The ratio test here is awkward because the ratio of
consecutive terms does not simplify cleanly; the root test reads the
$n$-th-power structure directly.

\engterm{Series 5}: $\displaystyle \sum_{n=2}^{\infty}
\frac{1}{n (\log n)^{2}}$. The terms shrink slightly faster than
$1/n$, the borderline. The integral test settles it: with $u = \log
x$, $\int_{2}^{\infty} \dd x/(x (\log x)^{2}) = \int_{\log 2}^{\infty}
\dd u/u^{2} = 1/\log 2$, finite. The series converges. The neighbouring
series $\sum 1/(n \log n)$, with a single power of $\log$, diverges by
the same substitution producing $\int \dd u/u$. The exponent on the
logarithm is the whole story at this boundary.

The five verdicts share a method: identify the dominant large-$n$
behaviour, choose the cheapest test that reads that behaviour, and
apply it. The decision tree below automates the choice.

\subsection{Accelerating a slowly convergent series}

A convergent series settles the convergence question and then leaves the
engineer with a second problem: a series that converges at $O(1/N)$ or
$O(1/N^{2})$ needs an impractical number of terms for engineering
accuracy. The tail bounds of the integral test quantified the trouble,
$\sum 1/n^{2}$ needs a thousand terms for three decimals. Two classical
transforms extract the limit from a handful of partial sums by modelling
the tail rather than summing it.

\engterm{Aitken's $\Delta^{2}$ process}. Suppose the partial sums
approach their limit geometrically, $S_{n} \approx S + A r^{n}$ for some
ratio $r$, the behaviour of any linearly convergent sequence. Three
consecutive partial sums then determine all three unknowns. Eliminating
$A$ and $r$ gives the accelerated estimate
\[
\widehat{S}_{n} = S_{n} - \frac{(S_{n+1} - S_{n})^{2}}{S_{n+2} - 2
S_{n+1} + S_{n}} = S_{n} - \frac{(\Delta S_{n})^{2}}{\Delta^{2} S_{n}},
\]
where $\Delta S_{n} = S_{n+1} - S_{n}$ and $\Delta^{2} S_{n} = S_{n+2}
- 2 S_{n+1} + S_{n}$. The transformed sequence $\widehat{S}_{n}$
converges to the same limit, and faster whenever the geometric model is
a good fit. A worked case settles the gain. The alternating series for
$\pi/4 = 1 - 1/3 + 1/5 - 1/7 + \cdots$ converges at $O(1/N)$: the raw
partial sums give $S_{4} = 0.8349$, off the true $\pi/4 = 0.7854$ by
$0.050$. Apply Aitken to $S_{2}, S_{3}, S_{4} = 0.8667, 0.7238, 0.8349$:
the first difference is $\Delta S_{2} = 0.7238 - 0.8667 = -0.1429$ and
the second is $\Delta^{2} S_{2} = 0.8349 - 2(0.7238) + 0.8667 = 0.2540$,
so the accelerated estimate is $0.8667 - (-0.1429)^{2}/0.2540 = 0.8667 -
0.0804 = 0.7863$, off by only $9 \times 10^{-4}$, about fifty times
closer than the raw $S_{4}$ from the same three partial sums. Iterating
Aitken on its own output, the $\Delta^{2}$-of-$\Delta^{2}$ chain,
reaches six-decimal accuracy from the first eight raw terms, a task that
costs the raw series tens of thousands of terms. The transform fails when the convergence is not
asymptotically geometric, and it amplifies round-off through the
second-difference denominator, which can vanish near the limit. The
defensive rule is to stop accelerating once $\abs{\Delta^{2} S_{n}}$
drops to the round-off floor.

\engterm{Euler's transform}. For an alternating series $\sum (-1)^{n}
a_{n}$ with $a_{n} > 0$ slowly decreasing, Euler's transform rewrites
the sum in terms of the forward differences of the $a_{n}$:
\[
\sum_{n=0}^{\infty} (-1)^{n} a_{n} = \sum_{k=0}^{\infty}
\frac{(-1)^{k}}{2^{k+1}} \Delta^{k} a_{0}, \qquad \Delta^{k} a_{0} =
\sum_{j=0}^{k} (-1)^{j} \binom{k}{j} a_{j}.
\]
When the $a_{n}$ are smooth, their high differences $\Delta^{k} a_{0}$
shrink rapidly, and the transformed series, with its extra $2^{-(k+1)}$
factor, converges geometrically even though the original converged at
$O(1/N)$. For $\log 2 = \sum (-1)^{n}/(n+1)$, the forward differences of
$a_{n} = 1/(n+1)$ shrink steadily, and twelve Euler terms recover $\log
2$ to better than five decimals where the raw alternating series needs
about $10^{5}$ terms for the same accuracy; the higher-order Levin and
Weniger transforms push the same dozen terms to full double precision.
The two transforms together are the reason an engineer almost never sums
a slowly convergent series directly. The convergence test certifies the
limit exists; the acceleration transform delivers it cheaply. The
script \texttt{code/aitken\_euler.py} runs both transforms on
$\pi/4$ and $\log 2$ and tabulates the term count each route needs for
six-decimal accuracy (\texttt{data/acceleration\_compare.csv}); the
detailed theory of these and the higher Shanks and Levin transforms is
in \cite{text:v2ch04-bender-orszag-1999}.

\begin{figure}[ht]
\centering
\begin{tikzpicture}[x=1cm, y=1cm, font=\small]

% Error vs term count for raw pi/4 partial sums, single Aitken, Euler.
% Data from data/acceleration_compare.csv.
% x: term count N = 2..14 mapped x_cm = (N-2)*0.65.
% y: log10(error). range 10^0 (top) to 10^-6 (bottom). y_cm = 6 + log10(err).

\draw[->, thick] (-0.3, 0) -- (8.7, 0) node[right] {$N$ (terms used)};
\draw[->, thick] (0, -0.3) -- (0, 7.0) node[above] {error (log scale)};

\foreach \n in {2,4,6,8,10,12,14} {
  \pgfmathsetmacro\xc{(\n-2)*0.65}
  \draw (\xc, -0.1) -- (\xc, 0.1) node[below, yshift=-0.12cm] {\n};
}
\foreach \p in {0,-1,-2,-3,-4,-5,-6} {
  \pgfmathsetmacro\yc{6+\p}
  \draw (-0.1, \yc) -- (0.1, \yc) node[left, xshift=-0.1cm] {$10^{\p}$};
}

% Raw partial sums of pi/4: err ~ 0.12 down to 0.018 (log10 -0.92 .. -1.75)
% y_cm = 6 + log10(err): N=2 -0.93->5.07; N=4 -1.21->4.79; N=6 -1.38->4.62;
% N=8 -1.51->4.49; N=10 -1.60->4.40; N=12 -1.68->4.32; N=14 -1.75->4.25
\draw[thick, engblack, mark=*, mark size=1pt] plot coordinates {
  (0,5.07) (1.3,4.79) (2.6,4.62) (3.9,4.49) (5.2,4.40) (6.5,4.32) (7.8,4.25) };
\node[engblack, anchor=west] at (7.85, 4.25) {raw};

% Single Aitken: N=4 2.065e-3 ->3.31; N=6 4.78e-4 ->2.68; N=8 1.78e-4 ->2.25;
% N=10 8.45e-5 ->1.93; N=12 4.65e-5 ->1.67; N=14 2.82e-5 ->1.45
\draw[thick, engblue, mark=square*, mark size=1pt] plot coordinates {
  (1.3,3.31) (2.6,2.68) (3.9,2.25) (5.2,1.93) (6.5,1.67) (7.8,1.45) };
\node[engblue, anchor=west] at (7.85, 1.45) {Aitken};

% Euler transform: N=4 2.35e-2 ->4.37; N=6 5.02e-3 ->3.70; N=8 1.12e-3 ->3.05;
% N=10 2.54e-4 ->2.40; N=12 5.86e-5 ->1.77; N=14 1.37e-5 ->1.14
\draw[thick, engred, mark=triangle*, mark size=1pt] plot coordinates {
  (1.3,4.37) (2.6,3.70) (3.9,3.05) (5.2,2.40) (6.5,1.77) (7.8,1.14) };
\node[engred, anchor=west] at (7.85, 1.14) {Euler};

\end{tikzpicture}
\caption[Convergence acceleration of the Leibniz series for $\pi/4$]{Error
of the raw Leibniz partial sums for $\pi/4 = 1 - 1/3 + 1/5 - \cdots$
against the error of the single-pass Aitken-$\Delta^{2}$ and
Euler-transformed estimates, as a function of the number of raw terms
used. The raw series barely improves, tracing the slow $O(1/N)$ harmonic
rate. The accelerated estimates fall two to three decades below it on the
same term budget, and iterating the Aitken transform on its own output (not
shown) reaches six correct digits from the first eight raw terms. The
convergence test settles that the limit exists; the acceleration transform
delivers it cheaply.}
\label{fig:vol02:ch04:acceleration}
\end{figure}


The acceleration figure (\autoref{fig:vol02:ch04:acceleration}) plots
the error of the raw partial sums of $\pi/4$ against the error of the
Aitken-accelerated and Euler-transformed estimates, on a log scale
versus term count. The raw series barely improves, tracing the slow
$O(1/N)$ harmonic rate; the accelerated estimates fall two to three
decades below it at the same term count. The gap is the engineering
payoff: the same information, a handful of partial sums, yields four or
more correct digits after a single transform where the raw sum has one
or two, and iterating the transform reaches six digits from the first
eight raw terms.

\section{Power series and Taylor series}\pathtag{core}

A \engterm{power series} centred at $x_{0}$ is
\[
\sum_{n=0}^{\infty} c_{n} (x - x_{0})^{n} = c_{0} + c_{1} (x - x_{0})
+ c_{2} (x - x_{0})^{2} + \cdots,
\]
where the coefficients $c_{n}$ are constants and $x$ is a real (or
complex) variable. The set of $x$ for which the series converges is an
interval (or disk in the complex case) around $x_{0}$, called the
\engterm{interval (or disk) of convergence}. Its half-width is the
\engterm{radius of convergence} $R$, given by
\[
\frac{1}{R} = \limsup_{n \to \infty} \sqrt[n]{\abs{c_{n}}}.
\]
Inside the interval the series converges absolutely. Outside it
diverges. At the endpoints it may converge or diverge, and the cases
must be checked individually.

\begin{figure}[ht]
\centering
\begin{tikzpicture}[x=1cm, y=1cm, font=\small]

% Complex plane axes
\draw[->, thick] (-4.5, 0) -- (4.5, 0) node[right] {$\Re x$};
\draw[->, thick] (0, -3.5) -- (0, 3.5) node[above] {$\Im x$};

% Origin
\fill (0, 0) circle (1.5pt);
\node[anchor=north east] at (0, 0) {$0$};

% Radius of convergence disk R = 1
\draw[thick, engblue, fill=engblue!10] (0, 0) circle (2);
\node[engblue] at (0, 1.0) {disk of convergence};
\node[engblue, font=\scriptsize] at (0, 0.5) {$|x| < 1$};

% Mark singularities at +1 and -1 on real axis (e.g., log(1+x) at x=-1; 1/(1-x) at x=1)
\fill[engred] (2, 0) circle (3pt);
\node[engred, anchor=south west] at (2, 0.1) {singularity at $x=1$};

\fill[engred] (-2, 0) circle (3pt);
\node[engred, anchor=south east] at (-2, 0.1) {singularity at $x=-1$};

% Marker convergence inside, divergence outside
\draw[->, thick, gray] (0.4, -1.5) -- (1.0, -1.0);
\node[gray, anchor=north, align=center] at (0.4, -1.6) {convergent\\inside};

\draw[->, thick, gray] (3.5, -2.0) -- (3.0, -1.5);
\node[gray, anchor=north, align=center] at (3.5, -2.1) {divergent\\outside};

% Radius arrow
\draw[<->, thick] (0, -2.7) -- (2, -2.7);
\node[anchor=north] at (1, -2.7) {$R = 1$};

\end{tikzpicture}
\caption[Disk of convergence of a power series]{The set of $x$ for
which a power series $\sum c_n (x - x_0)^n$ converges is a disk in
the complex plane centred at $x_0$ (here $x_0 = 0$). Inside the disk
the series converges absolutely. Outside the disk it diverges. At the
boundary $|x| = R$ it may do either, and each boundary point must be
checked individually. Two singularities of the underlying function
sit at $x = \pm 1$ in this illustration (the configuration for
$\log(1+x)$, $1/(1-x)$, and the binomial series with non-integer
exponent). The radius of convergence is the distance from $x_0$ to
the nearest singularity, even when the real-axis behaviour suggests
otherwise.}
\label{fig:vol02:ch04:radius-of-convergence}
\end{figure}


The disk picture, illustrated in
\autoref{fig:vol02:ch04:radius-of-convergence}, carries a useful
warning. The radius of convergence of the real-valued series
$1/(1 + x^{2}) = \sum (-1)^{n} x^{2n}$ is $1$, even though the
function $1/(1 + x^{2})$ is perfectly smooth at every real $x$. The
explanation is in the complex plane: the function has poles at
$x = \pm i$, and the nearest pole to $x_{0} = 0$ sits at distance
$1$. The radius of convergence is set by the nearest singularity of
the function, real or complex. An engineer who tries to expand
$1/(1 + x^{2})$ at $x_{0} = 0$ and use the expansion at $x = 2$
will discover that no number of terms gives a finite answer; the
series diverges term by term. The fix is to expand at a different
working point, where the nearest singularity is farther away.

\subsection{Four worked radii of convergence}

The radius formula $1/R = \limsup \sqrt[n]{\abs{c_{n}}}$ is rarely the
fastest route; the ratio of coefficients usually settles the question
with less work. Four cases cover the patterns.

\engterm{Constant coefficients}: $\sum x^{n}$. The ratio
$\abs{c_{n+1}/c_{n}} = 1$, so $R = 1$. The geometric series converges
on $(-1, 1)$ and diverges at both endpoints.

\engterm{Factorial in the denominator}: $\sum x^{n}/n!$. The ratio
$\abs{c_{n+1}/c_{n}} = n!/(n+1)! = 1/(n+1) \to 0$, so $1/R = 0$ and
$R = \infty$. The exponential series converges on all of $\R$, the
reason the exponential is so well-behaved as a building block.

\engterm{Factorial in the numerator}: $\sum n! \, x^{n}$. The ratio
$\abs{c_{n+1}/c_{n}} = (n+1)!/n! = n+1 \to \infty$, so $R = 0$. The
series converges only at $x = 0$. A power series with $R = 0$ is useless
as a function representation, and it is the formal signature of an
asymptotic (divergent) series of the kind the failure section treats.

\engterm{Polynomial coefficients}: $\sum x^{n}/n^{2}$. The ratio
$n^{2}/(n+1)^{2} \to 1$, so $R = 1$. At the endpoints the series
becomes $\sum (\pm 1)^{n}/n^{2}$, which converges absolutely (a
$p$-series with $p = 2$), so this series converges on the closed
interval $[-1, 1]$. The endpoint behaviour is independent of the
radius and must be checked term by term, as here.

The pattern across the four: the factorial controls the radius. A
factorial in the denominator gives infinite radius; a factorial in the
numerator gives zero radius; polynomial or constant coefficients give a
finite nonzero radius set by the nearest singularity. The endpoints are
a separate question, settled by the $p$-series and alternating-series
tests of section 4.2 and 4.3.

\subsection{Taylor series}

A function $f$ that has derivatives of all orders at $x_{0}$ has a
\engterm{Taylor series} at $x_{0}$:
\[
f(x) = \sum_{n=0}^{\infty} \frac{f^{(n)}(x_{0})}{n!} (x - x_{0})^{n}.
\]
The coefficients are the function's successive derivatives at $x_{0}$,
divided by the appropriate factorial. The series at $x_{0} = 0$ is the
\engterm{Maclaurin series}.

We treat the Taylor series as a known result. The formal derivation
uses the differential calculus we take up in Chapter 5; for the present
chapter, the series is the working object, and the question is the
size of the error when we truncate it.

\subsection{Taylor's theorem with remainder}

For a function $f$ with $n+1$ continuous derivatives on an interval
containing $x_{0}$ and $x$, the truncated Taylor series gives
\[
f(x) = \sum_{k=0}^{n} \frac{f^{(k)}(x_{0})}{k!} (x - x_{0})^{k}
+ R_{n}(x),
\]
where the \engterm{remainder} $R_{n}(x)$ admits several useful forms.
The \engterm{Lagrange form} is
\[
R_{n}(x) = \frac{f^{(n+1)}(\xi)}{(n+1)!} (x - x_{0})^{n+1}
\]
for some $\xi$ between $x_{0}$ and $x$. The engineer rarely knows
$\xi$, but a bound on $\abs{f^{(n+1)}}$ over the interval gives a bound
on $\abs{R_{n}(x)}$:
\[
\abs{R_{n}(x)} \leq \frac{M_{n+1}}{(n+1)!} \abs{x - x_{0}}^{n+1},
\]
where $M_{n+1} = \sup_{\xi} \abs{f^{(n+1)}(\xi)}$ on the interval. This
inequality is the chapter's main quantitative tool.

\begin{archetype}[Uncertainty, formalised as truncation error]
Volume I introduced uncertainty as the engineer's accounting for what
is not known about a measured quantity. The Lagrange remainder gives
that accounting a closed form for the special case of approximation by
truncated series. When the engineer replaces a function by its
$n$-term Taylor polynomial, the discarded tail is an uncertainty in the
computed value, and the remainder bound $M_{n+1} \abs{x - x_{0}}^{n+1}/
(n+1)!$ is its worst-case magnitude. The bound has the same role a
calibration uncertainty plays in Volume I: it converts an approximation
from a hopeful gesture into a quantity carried through the rest of the
calculation with a known error bar. Two levers control the bound, the
distance $\abs{x - x_{0}}$ from the working point and the number of
terms kept, and the engineer trades them against a stated budget. Every
worked example in this chapter that ends with an explicit error figure
is an instance of this archetype.
\end{archetype}

\subsection{The canonical engineering Taylor series}

Five Maclaurin series carry the bulk of engineering work. We list them
with their radii of convergence and a one-line note on what each is
used for.

\engterm{Exponential}.
\[
e^{x} = \sum_{n=0}^{\infty} \frac{x^{n}}{n!} = 1 + x + \frac{x^{2}}{2}
+ \frac{x^{3}}{6} + \frac{x^{4}}{24} + \cdots
\]
Radius of convergence: infinite. The exponential converges everywhere
on $\R$ (and on $\C$). It is the basic series of growth, decay, and
oscillation when combined with imaginary arguments.

\engterm{Sine}.
\[
\sin x = \sum_{n=0}^{\infty} \frac{(-1)^{n} x^{2n+1}}{(2n+1)!} = x
- \frac{x^{3}}{6} + \frac{x^{5}}{120} - \frac{x^{7}}{5040} + \cdots
\]
Radius of convergence: infinite. Engineers use the leading-order term
$\sin x \approx x$ for small $x$ in radians; the next term, $-x^{3}/6$,
sets the error of the small-angle approximation.

\engterm{Cosine}.
\[
\cos x = \sum_{n=0}^{\infty} \frac{(-1)^{n} x^{2n}}{(2n)!} = 1
- \frac{x^{2}}{2} + \frac{x^{4}}{24} - \frac{x^{6}}{720} + \cdots
\]
Radius of convergence: infinite. The leading correction $\cos x
\approx 1 - x^{2}/2$ is the second-order pendulum and small-deflection
formula.

\engterm{Logarithm}.
\[
\log(1 + x) = \sum_{n=1}^{\infty} \frac{(-1)^{n+1} x^{n}}{n} = x
- \frac{x^{2}}{2} + \frac{x^{3}}{3} - \frac{x^{4}}{4} + \cdots
\]
Radius of convergence: $1$. Converges for $-1 < x \leq 1$. The
leading-order $\log(1 + x) \approx x$ is the workhorse for relative
deviations: a $5\%$ change in a quantity contributes about $0.05$ to
the logarithm.

\engterm{Binomial}.
\[
(1 + x)^{\alpha} = \sum_{n=0}^{\infty} \binom{\alpha}{n} x^{n} = 1
+ \alpha x + \frac{\alpha (\alpha - 1)}{2} x^{2} + \frac{\alpha
(\alpha - 1)(\alpha - 2)}{6} x^{3} + \cdots,
\]
where $\binom{\alpha}{n} = \alpha(\alpha - 1) \cdots (\alpha - n +
1)/n!$ is the generalised binomial coefficient. Radius of convergence:
$1$ (when $\alpha$ is not a non-negative integer; otherwise the series
terminates). The binomial expansion supplies $\sqrt{1 + x} \approx 1 +
x/2 - x^{2}/8$ and $1/(1 + x) \approx 1 - x + x^{2} - x^{3} + \cdots$
as special cases.

These five series, together with the manipulations of section 4.4.5,
cover most of the smooth functions an engineer evaluates near a
working point.

A worked truncation-budget calculation puts the Lagrange bound to work.
Suppose the engineer must evaluate $\cos x$ for $x \in [-0.5, 0.5]$ rad
to an absolute accuracy of $10^{-8}$, using the fewest terms. The
remainder after keeping terms through degree $2m$ is bounded by
$\abs{R} \leq \abs{x}^{2m+2}/(2m+2)!$, since every derivative of
$\cos$ has magnitude at most $1$. The worst case is $\abs{x} = 0.5$.
Test orders: at $2m = 4$ (keeping $1 - x^{2}/2 + x^{4}/24$), the bound
is $0.5^{6}/6! = 0.015625/720 \approx 2.2 \times 10^{-5}$, too large;
at $2m = 6$, the bound is $0.5^{8}/8! \approx 3.9 \times
10^{-3}/40320 \approx 9.7 \times 10^{-8}$, still above budget; at
$2m = 8$, the bound is $0.5^{10}/10! \approx 9.8 \times
10^{-4}/3.63 \times 10^{6} \approx 2.7 \times 10^{-10}$, comfortably
below $10^{-8}$. Five nonzero terms (degrees $0, 2, 4, 6, 8$) meet the
budget across the whole interval. The calculation is the entire
engineering content of "how many terms": pick the worst-case point,
write the Lagrange bound, and step the order until the bound clears the
budget. The figure on convergence rate makes the per-term gain visible.

\begin{figure}[ht]
\centering
\begin{tikzpicture}[x=0.62cm, y=0.55cm, font=\small]

% Plot log10(|x|^n / n!) vs n for x = 2, 5, 10.
% y in log10; range about [-12, 4]. Map y_cm = (log10val + 12)/16 * height.
% We'll plot raw log10 values, axis spanning -12..4 (16 decades).
% Use a coordinate: x-axis n=0..16, y-axis log10 from -12 (bottom) to 4 (top).

% Axes
\draw[->, thick] (-0.3, -12) -- (16.5, -12) node[right] {$n$};
\draw[->, thick] (-0.3, -12.3) -- (-0.3, 4.8) node[above] {$\log_{10}(|x|^{n}/n!)$};

% y ticks every 4 decades
\foreach \y in {-12,-8,-4,0,4} {
  \draw (-0.45, \y) -- (-0.15, \y) node[left, xshift=-0.1cm] {$10^{\y}$};
}
% x ticks
\foreach \n in {0,4,8,12,16} {
  \draw (\n, -12.15) -- (\n, -11.85) node[below, yshift=-0.1cm] {\n};
}

% x = 2: log10(2^n/n!), exact
\draw[thick, engblack, mark=triangle*, mark size=1pt] plot coordinates {
  (0,0) (2,0.301) (4,-0.176) (6,-1.051) (8,-2.197) (10,-3.549)
  (12,-5.068) (14,-6.726) (16,-8.504) };
\node[engblack, anchor=west] at (16.1, -8.504) {$x=2$};

% x = 5: log10(5^n/n!), exact, peak near n=4..5
\draw[thick, engblue, mark=*, mark size=1pt] plot coordinates {
  (0,0) (2,1.097) (4,1.416) (6,1.336) (8,0.986) (10,0.430)
  (12,-0.293) (14,-1.155) (16,-2.137) };
\node[engblue, anchor=west] at (16.1, -2.137) {$x=5$};

% x = 10: log10(10^n/n!), exact, peak near n=9..10
\draw[thick, engred, mark=square*, mark size=1pt] plot coordinates {
  (0,0) (2,1.699) (4,2.620) (6,3.143) (8,3.394) (10,3.440)
  (12,3.320) (14,3.060) (16,2.679) };
\node[engred, anchor=west] at (16.1, 2.679) {$x=10$};

\end{tikzpicture}
\caption[Term magnitudes of the exponential series peak then collapse]{Term
magnitudes $\abs{x}^{n}/n!$ of the exponential series for $x = 2, 5,
10$, on a log scale. Each curve rises while $n < \abs{x}$, peaks near
$n \approx \abs{x}$, and then falls faster than any geometric sequence
because the consecutive-term ratio $\abs{x}/(n+1)$ keeps shrinking. The
post-peak collapse is the factorial winning the growth hierarchy of
section 4.1, and it is why a fixed-accuracy truncation of the
exponential, sine, or cosine series needs only a handful of terms once
$\abs{x}$ is modest.}
\label{fig:vol02:ch04:factorial-decay}
\end{figure}


The factorial-decay figure
(\autoref{fig:vol02:ch04:factorial-decay}) plots the term magnitudes
$\abs{x}^{n}/n!$ against $n$ for several fixed $x$. Each curve rises
while $n < \abs{x}$, peaks near $n \approx \abs{x}$, and then falls
faster than any geometric sequence because the ratio of consecutive
terms, $\abs{x}/(n+1)$, keeps shrinking. The shape is why the
exponential, sine, and cosine series converge everywhere and why a
fixed-budget truncation needs only a handful of terms once $\abs{x}$ is
modest. The post-peak collapse is the factorial winning the growth
hierarchy of section 4.1.

A sixth series, the \engterm{tangent}, is worth listing because the
project block uses it:
\[
\tan x = x + \frac{x^{3}}{3} + \frac{2 x^{5}}{15} + \frac{17 x^{7}}{315}
+ \cdots
\]
The coefficients are $2(2^{2n} - 1) B_{2n}/(2n)!$ where $B_{2n}$ is
the $2n$-th Bernoulli number; the closed form is in
\cite{text:v2ch04-abramowitz-stegun-1972}. The radius of convergence
is $\pi/2$, set by the poles of $\tan$ at $\pm \pi/2$. The same
restriction makes the tangent's small-angle approximation noticeably
narrower in safe range than the sine's: by the time $\theta = \pi/4
\approx 45^{\circ}$, the leading-order tangent has accumulated about
$27\%$ relative error.

\begin{mastery}
\engterm{Asymptotic series vs convergent series}. A series can be
asymptotic to a function without converging to it. The series
$\sum_{n=0}^{\infty} (-1)^{n} n! / x^{n+1}$, called the asymptotic
expansion of the exponential integral $E_{1}(x)$, diverges for every
real $x$, yet truncating it at the $N$-th term and using the bound
$|R_{N}| \leq |a_{N+1}|$ gives an excellent approximation for large
$x$ when $N$ is chosen optimally (typically $N \approx x$). The
optimal-$N$ rule is to stop adding terms when they stop shrinking;
the smallest term in an asymptotic series gives the best achievable
error. The detailed treatment is in
\cite{text:v2ch04-bender-orszag-1999}. Engineers meet asymptotic
series in singular-perturbation problems, boundary-layer
calculations, and the WKB approximation; the chapter's failure
section gives the warning, and the volume's later perturbation
chapter formalises the technique.
\end{mastery}

\subsection{Manipulating Taylor series}

Inside their disks of convergence, power series can be added,
subtracted, multiplied, divided (subject to the leading coefficient
not vanishing), composed, and integrated or differentiated term by
term. The composition rule justifies, for example, $e^{\sin x}$ near
$x = 0$:
\begin{align*}
e^{\sin x} &= 1 + \sin x + \frac{\sin^{2} x}{2} + \frac{\sin^{3} x}{6}
+ \cdots \\
&= 1 + \left(x - \frac{x^{3}}{6}\right) + \frac{1}{2}\left(x -
\frac{x^{3}}{6}\right)^{2} + \cdots \\
&= 1 + x + \frac{x^{2}}{2} - \frac{x^{4}}{8} + \cdots
\end{align*}
The reader collects terms by total power of $x$; the result is the
Maclaurin series of $e^{\sin x}$ to the order kept. The same
manipulations underwrite the standard small-angle composites
(amplitude responses, frequency-shifted oscillations, perturbed
expansions) that appear later in the chapter.

Division of power series is the manipulation engineers most often get
wrong, so a worked case is warranted. To find the series of $\tan x =
\sin x / \cos x$ near $0$ without quoting the Bernoulli-number formula,
divide the sine series by the cosine series:
\[
\tan x = \frac{x - x^{3}/6 + x^{5}/120 - \cdots}{1 - x^{2}/2 +
x^{4}/24 - \cdots}.
\]
Write $\tan x = c_{1} x + c_{3} x^{3} + c_{5} x^{5} + \cdots$ (only odd
powers, by the odd symmetry of $\tan$) and multiply through by the
cosine series:
\[
\left(c_{1} x + c_{3} x^{3} + c_{5} x^{5} + \cdots\right)
\left(1 - \frac{x^{2}}{2} + \frac{x^{4}}{24} - \cdots\right) = x -
\frac{x^{3}}{6} + \frac{x^{5}}{120} - \cdots.
\]
Match coefficients. The $x^{1}$ term gives $c_{1} = 1$. The $x^{3}$
term gives $c_{3} - c_{1}/2 = -1/6$, so $c_{3} = -1/6 + 1/2 = 1/3$. The
$x^{5}$ term gives $c_{5} - c_{3}/2 + c_{1}/24 = 1/120$, so $c_{5} =
1/120 + (1/3)/2 - 1/24 = 1/120 + 1/6 - 1/24 = 2/15$. The result
$\tan x = x + x^{3}/3 + 2 x^{5}/15 + \cdots$ matches the listed
tangent series, derived now by division rather than by quoting. The
method generalises to any quotient whose denominator has a nonzero
constant term, and it is the safe route when no closed form for the
coefficients is at hand.

\begin{mastery}
\engterm{Power series as ODE solutions}. A power series can solve a
differential equation by undetermined coefficients, the technique
formalised in Chapter 7. The preview: to solve $y' = y$ with $y(0) =
1$, write $y = \sum_{n=0}^{\infty} c_{n} x^{n}$, differentiate term by
term to get $y' = \sum_{n=1}^{\infty} n c_{n} x^{n-1} =
\sum_{m=0}^{\infty} (m+1) c_{m+1} x^{m}$, and match against $y =
\sum c_{m} x^{m}$. Coefficient matching gives the recursion $(m+1)
c_{m+1} = c_{m}$, that is $c_{m+1} = c_{m}/(m+1)$, with $c_{0} = 1$ from
the initial condition. The recursion unwinds to $c_{m} = 1/m!$, and the
series is $\sum x^{m}/m! = e^{x}$, the known solution. The same machine
solves equations with no elementary closed-form solution, including the
Airy equation $y'' = x y$ whose series solution defines the Airy
functions used in optics and quantum tunnelling. The power-series method
turns a differential equation into an algebraic recursion on the
coefficients, and the radius of convergence of the resulting series is
set by the nearest singularity of the equation's coefficients in the
complex plane.
\end{mastery}

\subsection{Pad\'{e} approximants: a rational alternative to truncation}

A truncated Taylor series is a polynomial, and a polynomial diverges to
$\pm\infty$ as its argument grows. When the function being approximated
levels off, has a pole, or must be evaluated near the edge of the
Taylor radius, a ratio of two polynomials does far better with the same
number of coefficients. The \engterm{Pad\'{e} approximant} of order
$[m/n]$ is the rational function
\[
R_{m/n}(x) = \frac{p_{0} + p_{1} x + \cdots + p_{m} x^{m}}{1 + q_{1} x
+ \cdots + q_{n} x^{n}}
\]
whose Taylor expansion at $x = 0$ matches the function's Taylor series
through order $m + n$. The coefficients are fixed by that matching
condition: cross-multiply, equate the series of $R_{m/n}$ with the
target series term by term up to $x^{m+n}$, and solve the resulting
linear system for the $q_{j}$ first, then the $p_{i}$.

A worked $[2/2]$ approximant of the exponential shows the gain. The
Taylor series is $e^{x} = 1 + x + x^{2}/2 + x^{3}/6 + x^{4}/24 +
\cdots$. Write $R_{2/2}(x) = (p_{0} + p_{1} x + p_{2} x^{2})/(1 + q_{1}
x + q_{2} x^{2})$ and require $(1 + q_{1} x + q_{2} x^{2}) e^{x} = p_{0}
+ p_{1} x + p_{2} x^{2} + O(x^{5})$. Matching coefficients of $x^{0}$
through $x^{4}$:
\begin{align*}
x^{0}&: \; p_{0} = 1, & x^{3}&: \; 0 = \tfrac{1}{6} +
\tfrac{1}{2} q_{1} + q_{2}, \\
x^{1}&: \; p_{1} = 1 + q_{1}, & x^{4}&: \; 0 = \tfrac{1}{24} +
\tfrac{1}{6} q_{1} + \tfrac{1}{2} q_{2}. \\
x^{2}&: \; p_{2} = \tfrac{1}{2} + q_{1} + q_{2},
\end{align*}
The last two equations give $q_{1} = -1/2$, $q_{2} = 1/12$, and then
$p_{1} = 1/2$, $p_{2} = 1/12$. The approximant is
\[
e^{x} \approx \frac{1 + x/2 + x^{2}/12}{1 - x/2 + x^{2}/12},
\]
which is the diagonal Pad\'{e} that underlies the implicit-midpoint and
Crank-Nicolson time integrators of Chapter 7. Its Taylor expansion
matches $e^{x}$ through $x^{4}$ (the first discrepancy is at $x^{5}$,
where the approximant gives $x^{5}/144$ against the true $x^{5}/120$),
the same matching order as the degree-four Taylor polynomial, from the
same five coefficients. At $x = 1$ the Pad\'{e} returns $2.71429$
against $e = 2.71828$, a relative error of $1.5 \times 10^{-3}$, while
the degree-four Taylor polynomial $1 + 1 + 1/2 + 1/6 + 1/24 = 2.70833$
is off by $3.7 \times 10^{-3}$, about two and a half times worse. The
rational form's real advantage shows on functions that level off or
have poles: for $f(x) = \log(1 + x)/x$, which the Taylor series
represents only on $\abs{x} < 1$, the $[1/1]$ Pad\'{e} stays accurate
well past the Taylor radius, because the rational form spends its
coefficients on the shape of the function rather than on a single
polynomial trend, and it can place a pole where the function has one.
The diagonal $e^{x}$ Pad\'{e} above is prized in time integration not
for large-$x$ accuracy (where its denominator's complex roots make it
mediocre on the real axis) but because $\abs{R_{2/2}(i\omega)} = 1$
exactly, so the Crank-Nicolson scheme it generates neither grows nor
damps a pure oscillation, the stability property an integrator needs. The hazard is the spurious pole: the
denominator of a Pad\'{e} approximant can vanish at a point where the
true function is finite, producing a \engterm{Froissart doublet}, a
near-cancelling pole-zero pair that a careless evaluation reads as a
singularity. The treatment of Pad\'{e} theory and its convergence is in
\cite{text:v2ch04-bender-orszag-1999}.

\subsection{A worked power-series ODE solution: the Airy equation}

The mastery box above previewed the power-series method on $y' = y$,
whose answer was already known. The method earns its place on an
equation with no elementary solution. The \engterm{Airy equation}
\[
y'' = x y
\]
governs the transition between oscillatory and exponential behaviour at
a turning point, and it appears in optics (the intensity near a caustic),
in quantum mechanics (a particle near a linear potential barrier), and
in the asymptotics of the WKB approximation. Its solutions, the Airy
functions, are defined by their power series.

Write $y = \sum_{n=0}^{\infty} c_{n} x^{n}$. Differentiating twice,
$y'' = \sum_{n=2}^{\infty} n(n-1) c_{n} x^{n-2} = \sum_{m=0}^{\infty}
(m+2)(m+1) c_{m+2} x^{m}$, while $x y = \sum_{m=1}^{\infty} c_{m-1}
x^{m}$. Matching coefficients of $x^{m}$:
\[
(m+2)(m+1) c_{m+2} = c_{m-1} \quad (m \geq 1), \qquad
2 c_{2} = 0.
\]
The recursion $c_{m+2} = c_{m-1}/((m+2)(m+1))$ links coefficients three
indices apart, so the series splits into three families seeded by
$c_{0}$, $c_{1}$, and $c_{2}$. The condition $c_{2} = 0$ kills the
third family, leaving two independent solutions, one even-dominated
(seeded by $c_{0}$) and one odd-dominated (seeded by $c_{1}$). Unwinding
the recursion from $c_{0} = 1, c_{1} = 0$:
\[
c_{3} = \frac{c_{0}}{3 \cdot 2} = \frac{1}{6}, \quad c_{6} =
\frac{c_{3}}{6 \cdot 5} = \frac{1}{180}, \quad c_{9} = \frac{c_{6}}
{9 \cdot 8} = \frac{1}{12960}, \ldots
\]
giving the first solution
\[
y_{1}(x) = 1 + \frac{x^{3}}{6} + \frac{x^{6}}{180} + \frac{x^{9}}
{12960} + \cdots = \sum_{k=0}^{\infty} \frac{x^{3k}}{(3k)!}
\prod_{j=1}^{k} (3j - 2),
\]
and from $c_{0} = 0, c_{1} = 1$ the second solution $y_{2}(x) = x +
x^{4}/12 + x^{7}/504 + \cdots$. The ratio test on the first series:
the term ratio is $x^{3}/((3k+3)(3k+2)) \to 0$ for every fixed $x$, so
the radius of convergence is infinite. Every coefficient is computed by
a single division, and the series converges everywhere, which is what
makes the power-series method the practical route to the Airy functions
for moderate $\abs{x}$. For large $\abs{x}$ the convergent series loses
to catastrophic cancellation (the terms grow before they decay, exactly
the $e^{-20}$ trap of the failure section), and the engineer switches to
the asymptotic expansion, the complementary tool the mastery box on
asymptotic series named. The standard Airy functions $\mathrm{Ai}$ and
$\mathrm{Bi}$ are the particular linear combinations of $y_{1}$ and
$y_{2}$ fixed by their behaviour at infinity; the coefficients and the
connection formulae are tabulated in
\cite{text:v2ch04-abramowitz-stegun-1972}. The power-series method
turned a differential equation with no elementary solution into a
three-term recursion that a calculator can run, and the radius of
convergence came for free from the ratio test on that recursion.

\subsection{Euler's identity}

A clean payoff. Substitute $x = i\theta$ into the exponential series:
\[
e^{i\theta} = \sum_{n=0}^{\infty} \frac{(i\theta)^{n}}{n!}.
\]
The powers of $i$ cycle: $i^{0} = 1$, $i^{1} = i$, $i^{2} = -1$,
$i^{3} = -i$, $i^{4} = 1$, $\ldots$. Splitting the sum into even and
odd indices:
\[
e^{i\theta} = \left(1 - \frac{\theta^{2}}{2} + \frac{\theta^{4}}{24}
- \cdots\right) + i \left(\theta - \frac{\theta^{3}}{6}
+ \frac{\theta^{5}}{120} - \cdots\right) = \cos\theta + i \sin\theta.
\]
The identity $e^{i\theta} = \cos\theta + i \sin\theta$ unifies the
exponential, sine, and cosine series. Setting $\theta = \pi$ gives
$e^{i\pi} = -1$, the relation between five fundamental constants
known as Euler's identity. Volume II Chapter 3 introduced the rotation
interpretation of complex multiplication; the Taylor series machinery
gives the algebraic backing for the same identity.

\section{Limits: epsilon-delta informally; engineering use}\pathtag{standard}

The limit of a function captures the value the function approaches as
the input approaches a target. Limits are the formal machinery of
"what happens at the edge of the working range," and they underwrite
every derivative and every integral that follows in this volume.

\subsection{The intuitive picture}

For a function $f$ defined near a point $x_{0}$ (but not necessarily
at $x_{0}$), we say
\[
\lim_{x \to x_{0}} f(x) = L
\]
if $f(x)$ can be made as close to $L$ as we like by taking $x$
sufficiently close to $x_{0}$ (with $x \neq x_{0}$).

The qualifier "with $x \neq x_{0}$" matters. The function value
$f(x_{0})$, if defined, is allowed to disagree with the limit. The
function $f(x) = (x^{2} - 1)/(x - 1)$ has $f(x) = x + 1$ for $x \neq
1$, so $\lim_{x \to 1} f(x) = 2$, even though $f(1)$ is not defined
(the formula has $0/0$ at $x = 1$). The limit answers a question about
behaviour near a point, not at the point.

\subsection{Epsilon-delta as motivating discipline}

The formal definition makes "as close as we like" and "sufficiently
close" precise.

\engterm{Epsilon-delta limit}. We say $\lim_{x \to x_{0}} f(x) = L$ if,
for every $\varepsilon > 0$, there exists a $\delta > 0$ such that
$0 < \abs{x - x_{0}} < \delta$ implies $\abs{f(x) - L} < \varepsilon$.

The structure of the definition is a dialogue: the user specifies a
tolerance $\varepsilon$ on the output; the prover supplies a
corresponding tolerance $\delta$ on the input that guarantees the
output stays within tolerance. The limit exists when, for every demand,
a working $\delta$ can be produced.

We are using this definition as motivating discipline rather than as a
working tool. An engineer rarely computes limits by writing out
epsilon-delta arguments. The engineer uses limit laws (sums, products,
quotients of limits go to sums, products, and quotients of the limits,
when defined), continuity (a continuous function has $\lim_{x \to x_{0}}
f(x) = f(x_{0})$, by definition of continuity), the standard limits
(such as $\lim_{x \to 0} \sin x / x = 1$), and L'H\^{o}pital's rule
(treated in Chapter 5). The epsilon-delta definition is the
definition's truth condition; the working tools live above it.

The discipline the definition imposes is that "behaviour near a point"
has a precise meaning, and that statements like "the function tends to
$L$" are decidable rather than aesthetic. An engineer who has internalised
this discipline reads "the system reaches steady state" as a claim with a
specific epsilon-delta content.

\subsection{Worked limits and the indeterminate forms}

Five worked limits exercise the standard tools and name the traps.

\engterm{The standard small-angle limit}. $\lim_{x \to 0} \sin x / x =
1$. The Taylor expansion gives $\sin x / x = 1 - x^{2}/6 + \cdots$,
whose limit at $x = 0$ is $1$ by inspection. This limit is the reason
the small-angle approximation has unit slope, and it underwrites the
derivative of $\sin$ in Chapter 5.

\engterm{A removable singularity}. $\lim_{x \to 0} (1 - \cos x)/x^{2}$.
The numerator and denominator both vanish, a $0/0$ indeterminate form.
Expand: $1 - \cos x = x^{2}/2 - x^{4}/24 + \cdots$, so $(1 - \cos
x)/x^{2} = 1/2 - x^{2}/24 + \cdots \to 1/2$. The series turns the
indeterminate form into a clean leading coefficient. This limit sets
the curvature of the cosine at the origin and appears in the
second-order small-deflection formula.

\engterm{The exponential definition}. $\lim_{n \to \infty} (1 + x/n)^{n}
= e^{x}$. Take the logarithm: $n \log(1 + x/n) = n (x/n - x^{2}/(2
n^{2}) + \cdots) = x - x^{2}/(2n) + \cdots \to x$, so the limit of the
logarithm is $x$ and the limit of the expression is $e^{x}$. This is the
continuous-compounding limit and the bridge between discrete and
continuous growth in Chapter 7.

\engterm{A limit at infinity}. $\lim_{x \to \infty} x (\sqrt{x^{2} + 1}
- x)$. The form is $\infty \cdot 0$. Rationalise: $\sqrt{x^{2} + 1} - x
= 1/(\sqrt{x^{2} + 1} + x)$, so $x (\sqrt{x^{2} + 1} - x) = x/(\sqrt{x^{2}
+ 1} + x) \to 1/2$ as $x \to \infty$. The same rationalisation appears
whenever an engineer subtracts two nearly equal large quantities, and it
is the symbolic cousin of the catastrophic-cancellation fix in the
failure section.

\engterm{A one-sided limit that fails}. $\lim_{x \to 0} \sin(1/x)$ does
not exist. As $x \to 0$ the argument $1/x$ runs through all of its
period infinitely often, so $\sin(1/x)$ oscillates between $-1$ and $1$
without settling. No epsilon-delta $\delta$ can confine the output to a
small interval, because every neighbourhood of $0$ contains inputs
giving outputs at both extremes. The example is the standard warning
that "the input approaches a target" does not by itself guarantee the
output approaches anything.

\engterm{The squeeze theorem}, used once: $\lim_{x \to 0} x^{2}
\sin(1/x) = 0$. The oscillating factor is bounded, $\abs{\sin(1/x)}
\leq 1$, so $-x^{2} \leq x^{2} \sin(1/x) \leq x^{2}$. Both bounding
sequences tend to $0$, so the trapped quantity does too. The squeeze
theorem is the formal tool for limits where a bounded but
non-convergent factor multiplies a vanishing one, a pattern common in
windowed signals and damped oscillations.

\subsection{One-sided limits, infinite limits, limits at infinity}

Three variants extend the basic definition.

\engterm{One-sided limit}. $\lim_{x \to x_{0}^{+}} f(x) = L$ if $f(x)
\to L$ as $x$ approaches $x_{0}$ from the right (that is, with
$x > x_{0}$). The left-hand limit $\lim_{x \to x_{0}^{-}} f(x)$ is
defined similarly. The two-sided limit exists if and only if both
one-sided limits exist and are equal.

\engterm{Infinite limit}. $\lim_{x \to x_{0}} f(x) = +\infty$ if $f(x)$
exceeds every prescribed $M > 0$ for $x$ sufficiently close to $x_{0}$.
The function does not converge to a finite value; it grows without
bound as $x$ approaches $x_{0}$. The same idea defines $-\infty$ as a
limit and limits as $x \to \pm \infty$.

\engterm{Limit at infinity}. $\lim_{x \to \infty} f(x) = L$ if $f(x)
\to L$ as $x$ grows without bound. The behaviour of a sequence at
infinity is the case where $x$ ranges over the integers.

\subsection{Engineering uses of limits}

Five engineering questions are limit questions in disguise:

\begin{itemize}[leftmargin=*]
  \item \engterm{Steady state}. The voltage on a charging capacitor is
    $V(t) = V_{0}(1 - e^{-t/\tau})$. The steady-state voltage is
    $\lim_{t \to \infty} V(t) = V_{0}$. The time required to reach a
    fraction of $V_{0}$ is set by the time constant $\tau$. Volume II
    Chapter 7 treats the underlying ODE; the limit is the answer to
    "where does this system end up?"
  \item \engterm{Linearity at small input}. A nonlinear device behaves
    approximately linearly at small input. The constant of
    proportionality is $\lim_{x \to 0} f(x)/x$ when this limit exists.
    For $f(x) = \sin x$, the limit is $1$; for $f(x) = \tan x$, the
    limit is $1$; for $f(x) = e^{x} - 1$, the limit is $1$. The
    common-leading-coefficient pattern is the small-signal regime.
  \item \engterm{Asymptotic behaviour}. A drag force $F(v) = c_{1} v +
    c_{2} v^{2}$ has linear-drag behaviour for small $v$ (where the
    quadratic term is negligible) and quadratic-drag behaviour for
    large $v$ (where the linear term is negligible). The crossover is
    where $c_{1} v \sim c_{2} v^{2}$, that is, $v \sim c_{1}/c_{2}$.
    The behaviour at the edges of the working range is the limit
    behaviour at $v \to 0$ and $v \to \infty$.
  \item \engterm{Continuity at a boundary}. A property of a material
    has one expression in one regime and another expression in the
    next. The two expressions match at the boundary if and only if the
    one-sided limits agree. Engineers chase this matching condition at
    contact surfaces, phase transitions, and method-of-stripes joints.
  \item \engterm{Convergence of an iteration}. A numerical iteration
    is useful when the sequence of iterates has a limit. The
    convergence question becomes a limit question.
\end{itemize}

The unifying engineering use is the question \emph{what happens at the
edge}, where "edge" means at zero, at infinity, at a working-point
deviation, at a boundary, or at the limit of a long-running process.
Limits are the language in which that question is asked.

\subsection{Continuity, the intermediate value theorem, and the engineer's
working tools}

Three results bundle a great deal of the working content of limits.

\engterm{Continuity}. A function $f$ is \engterm{continuous} at $x_{0}$
if $\lim_{x \to x_{0}} f(x) = f(x_{0})$. A function continuous at every
point of an interval is continuous on the interval. The arithmetic
operations of sum, product, quotient (where the denominator is
nonzero), and composition preserve continuity. Polynomials are
continuous everywhere; rational functions are continuous wherever the
denominator does not vanish; $\sin$, $\cos$, $\exp$, $\log$ are
continuous on their domains.

\engterm{Intermediate value theorem}. If $f$ is continuous on $[a, b]$
and $y$ is a value between $f(a)$ and $f(b)$, there is some
$c \in [a, b]$ with $f(c) = y$. The engineering use is the bisection
root finder: a sign change of $f$ between $a$ and $b$ guarantees a
root in the interval, and successive halving locates it. The same
theorem underwrites the existence of a continuous boundary value
between any two regimes of a continuously varying parameter.

\engterm{Extreme value theorem}. A continuous function on a closed,
bounded interval attains its maximum and its minimum. The engineering
use is in optimisation: a continuous objective on a closed feasible
set has a maximiser and a minimiser, even when the gradient never
vanishes (the extrema may sit at the boundary). The theorem fails on
open intervals; an engineer who optimises on $(0, 1]$ instead of
$[0, 1]$ may find the supremum is not attained.

A worked bisection run shows the intermediate value theorem doing
engineering work. To locate the root of $f(x) = x^{3} - x - 2$ on
$[1, 2]$, note $f(1) = -2 < 0$ and $f(2) = 4 > 0$, so by the
intermediate value theorem a root sits in $[1, 2]$. Bisect: $f(1.5) =
-0.125 < 0$, so the root is in $[1.5, 2]$; $f(1.75) = 1.61 > 0$, so the
root is in $[1.5, 1.75]$; $f(1.625) = 0.666 > 0$, root in $[1.5,
1.625]$; $f(1.5625) = 0.252 > 0$, root in $[1.5, 1.5625]$. Each step
halves the bracket, so the interval width after $n$ steps is
$2^{-n}$ times the original. To reach an absolute tolerance of
$10^{-6}$ on a unit bracket takes $\lceil \log_{2}(10^{6}) \rceil = 20$
bisections, a sequence of nested intervals whose common limit is the
root. The convergence is geometric with ratio $1/2$, the slowest of the
root finders but the only one that needs nothing beyond a sign change
and continuity. The root here is $x \approx 1.5214$.

These three results compose into a working toolkit. The bisection
root finder uses the intermediate value theorem to certify the
existence of a root and the continuity of $f$ to guarantee that
halving the interval at the sign-change midpoint preserves the
certificate. The Newton root finder uses the differentiability of
$f$ (a strictly stronger condition than continuity, treated in
Chapter 5) to achieve quadratic convergence. The two algorithms
share the same epsilon-delta truth condition; the difference is what
they cost per step.

The cost difference is worth tabulating, because it drives the
engineer's choice of method. To reach an error of $10^{-12}$ from a
starting error of $1$, bisection needs $40$ steps (each halving the
bracket, $2^{-40} \approx 9 \times 10^{-13}$); a fixed-point iteration
with rate $0.5$ needs the same $40$ steps; a fixed-point iteration with
rate $0.1$ needs $12$ steps; the secant method, with order $\phi
\approx 1.618$, needs about $7$ steps; Newton's method, with order $2$,
needs about $5$ steps ($1, 10^{-1}, 10^{-2}, 10^{-4}, 10^{-8},
10^{-16}$, the correct-digit count doubling each step once it is in the
quadratic regime). The ordering is the convergence-rate hierarchy of
the chapter, and the comparator script
\texttt{code/convergence\_comparator.py} prints the full table. The
engineering judgment is that Newton's five steps cost a derivative
evaluation each, the secant's seven steps cost only a function
evaluation each, and bisection's forty steps cost the least per step
but the most in total, so the right method depends on the relative cost
of a function evaluation, a derivative evaluation, and a guaranteed
bracket. Quadratic convergence is fastest in iteration count and is not
always cheapest in wall-clock time.

\subsection{A worked convergence-rate analysis: when Newton slows down}

Newton's iteration converges quadratically at a simple root and only
linearly at a multiple root, and the engineer who does not know which
case is in hand will misread the iteration's residuals. The analysis
that distinguishes the two is worth a full worked pass, because it is
the prototype for reading any iteration's convergence rate from its
residual sequence.

Take $f(x) = (x - 2)^{2}$, which has a double root at $x = 2$. Newton's
iteration is $x_{n+1} = x_{n} - f(x_{n})/f'(x_{n}) = x_{n} - (x_{n} -
2)^{2}/(2(x_{n} - 2)) = x_{n} - (x_{n} - 2)/2 = (x_{n} + 2)/2$. The
error $e_{n} = x_{n} - 2$ then satisfies $e_{n+1} = e_{n}/2$ exactly:
linear convergence with rate $1/2$, not quadratic. Starting from $x_{0}
= 3$, the errors run $1, 0.5, 0.25, 0.125, \ldots$, halving each step,
and reaching $10^{-12}$ takes about forty iterations rather than the
five that a simple root would need. The general result is that at a root
of multiplicity $m$, Newton converges linearly with rate $1 - 1/m$, so
the double root gives rate $1/2$, a triple root rate $2/3$, and so on,
each higher multiplicity slower.

The diagnostic the engineer reads is the ratio $e_{n+1}/e_{n}$. For a
simple root this ratio tends to zero (the next error is a vanishing
fraction of the current one, the signature of superlinear convergence);
for a multiplicity-$m$ root it tends to the constant $1 - 1/m$. A
residual sequence whose successive ratios settle to a constant near
$0.5$ rather than collapsing to zero is the fingerprint of a double
root, and the eye that expected quadratic convergence reads it as a
broken solver when the solver is working correctly on a harder problem.

Two fixes restore fast convergence. The \engterm{modified Newton}
iteration $x_{n+1} = x_{n} - m \, f(x_{n})/f'(x_{n})$, with the known
multiplicity $m$ folded in, recovers quadratic convergence: for the
double root, the factor $m = 2$ turns the update into $x_{n+1} = x_{n}
- (x_{n} - 2)$, which lands on the root in one step. When the
multiplicity is unknown, \engterm{Steffensen's method} accelerates a
linearly convergent fixed-point iteration $a_{n+1} = g(a_{n})$ to
quadratic order without derivatives, by applying Aitken's $\Delta^{2}$
transform of section 4.3 to each triple of iterates:
\[
a_{n+1} = a_{n} - \frac{(g(a_{n}) - a_{n})^{2}}{g(g(a_{n})) - 2 g(a_{n})
+ a_{n}}.
\]
The same $\Delta^{2}$ that accelerated a slowly convergent series
accelerates a slowly convergent iteration, because both are sequences
approaching a limit geometrically. Applied to the Dottie-number
iteration $g(a) = \cos a$ of section 4.1, which converged linearly at
rate $0.674$, Steffensen reaches the fixed point $0.739085$ to
machine precision in four iterations rather than the thirty-five the
raw iteration needed. The convergence-rate ladder, the diagnostic
ratios, and the acceleration fixes are one connected body of technique:
read the rate from the residual ratios, identify the obstacle (multiple
root, slow contraction), and apply the matching accelerator. The
analysis is treated alongside the stiff-ODE convergence theory in
\cite{text:v2ch04-hairer-wanner-2008}.

\begin{mastery}
\engterm{Interval arithmetic for rigorous tails}. The truncation bounds
of this chapter are inequalities the engineer trusts; interval
arithmetic makes them machine-checkable certificates. In interval
arithmetic every quantity is a pair $[\underline{x}, \overline{x}]$
bracketing the true value, and every operation is rounded outward so
that the computed interval provably contains the exact result.
Evaluating a truncated series in interval arithmetic, and adding an
interval $[-B, B]$ for the rigorously bounded tail $B$, returns an
interval guaranteed to contain the true sum, round-off included. The
method converts the Lagrange remainder bound from a hand inequality
into a verified enclosure: a Taylor model carries a polynomial plus an
interval remainder through a whole computation, and the final interval
width is a proof of accuracy rather than an estimate of it. Engineers
use interval and Taylor-model arithmetic where a wrong answer is
catastrophic and a verified one is mandatory, in spacecraft trajectory
validation, in computer-assisted proofs (the proof of the Kepler
conjecture and the verified bounds on the Lorenz attractor both rest on
interval enclosures), and in safety-critical control certification. The
cost is that intervals widen pessimistically under repeated operations,
the \engterm{dependency problem}, so naive interval evaluation of a long
expression can return a uselessly wide bracket; Taylor models and affine
arithmetic are the engineering remedies that track correlations and keep
the enclosure tight.
\end{mastery}

\begin{estimation}
\textbf{Question.} An iterative scheme has linear convergence with
rate $r = 0.7$ per iteration. The starting error is $e_{0} = 1$.
A second scheme has quadratic convergence with constant $C = 1$ and
the same starting error. Estimate the iteration count each scheme
needs to reach $e_{n} \leq 10^{-9}$.

\textbf{Estimate.} The linear scheme satisfies $e_{n} \approx r^{n}
e_{0}$, so $n \approx -\log_{10}(10^{-9})/\log_{10}(1/0.7) = 9/0.155
\approx 58$ iterations. The quadratic scheme satisfies $e_{n}
\approx e_{0}^{2^{n}}$, so for $e_{0} = 1$ the error squares each
step: starting from a perturbed $e_{0} = 0.5$, we get $0.5, 0.25,
0.0625, 3.9 \times 10^{-3}, 1.5 \times 10^{-5}, 2.3 \times 10^{-10}$.
About six iterations of Newton from a half-decent starting point
suffice. The quadratic scheme is ten times faster than the linear
scheme when the starting error is moderate. The breakeven analysis
is recovered by the comparator script
\texttt{code/convergence\_comparator.py}.
\end{estimation}

\section{Failure: divergent series that look convergent}\pathtag{core}

The chapter's failure mode is a series whose terms shrink, whose
partial sums seem to be settling, but which never reaches a finite
limit. The harmonic series is the canonical case.

\subsection{The harmonic-series trap}

A reader confronted with $\sum 1/n$ for the first time often guesses
that it converges. The terms shrink. Each new term is smaller than the
last. The partial sums grow more slowly with each step. By visual
inspection they look like they are approaching some finite value.

They are not. The partial sums grow as $\log N + \gamma$, which is
unbounded. The growth is just slow enough that the engineer's eye does
not notice the trend without checking. After ten thousand terms, the
sum is about $9.79$. After a million, $14.39$. After a billion, $21.30$.
After Avogadro's number of terms, $55.99$. The growth is slow, but
unbounded.

The engineering hazard is that an iteration whose residual decreases
like $1/n$ rather than like $1/n^{2}$ does not converge in the sense
that allows the engineer to truncate at finite $n$ with a controlled
error. The truncation residual is the harmonic tail $\sum_{k > N} 1/k$,
which itself diverges.

\subsection{Diagnostics}

Three diagnostics help the working engineer separate convergent from
divergent series with shrinking terms.

\engterm{Term test for divergence}. If the terms do not shrink to zero,
the series diverges. This is the cheap rejection.

\engterm{Comparison with $1/n^{p}$}. If the terms behave asymptotically
like $1/n^{p}$, the series converges for $p > 1$ and diverges for
$p \leq 1$. The difference between $p = 1.0$ and $p = 1.01$ is the
difference between divergence and convergence. The $p = 1$ boundary is
sharp.

\engterm{Look at the partial sums on a log scale}. A truly convergent
series has partial sums that approach a horizontal asymptote when
plotted on a linear scale. A logarithmically divergent series, plotted
on a log-log scale of $S_{N}$ versus $N$, has $S_{N}$ growing as
$\log \log N$ asymptotically, which still goes to infinity but does so
visibly slowly. A power-law-divergent series has $S_{N}$ growing as a
power of $N$ on the same plot. The visual diagnostic is fragile; the
analytical diagnostic above is more reliable.

\subsection{Why the boundary is sharp}

The boundary between convergence and divergence is set by the
borderline behaviour of $\sum 1/(n \log n)$, $\sum 1/(n \log n \log
\log n)$, and so on, each marginally divergent. Asymptotic analyses
of borderline series are useful in queue-theoretic and renewal-theoretic
calculations and in some convergence analyses for stochastic
algorithms. For working engineers, the take-away is procedural: a
series whose terms shrink as $1/n$ or slower is suspect, and the
engineer should switch to a comparison or integral-test argument
before assuming convergence.

\subsection{The asymptotic-series trap: divergence that helps, then hurts}

The harmonic series diverges and looks convergent. The complementary
trap is a series that diverges yet gives an excellent answer if
truncated at the right place, and a worse answer if pushed further. The
asymptotic expansion of the exponential integral,
\[
E_{1}(x) = \int_{x}^{\infty} \frac{e^{-t}}{t} \, \dd t \sim
\frac{e^{-x}}{x}\left(1 - \frac{1}{x} + \frac{2!}{x^{2}} -
\frac{3!}{x^{3}} + \cdots\right),
\]
has terms of magnitude $n!/x^{n}$ that shrink while $n < x$ and grow
once $n > x$, because the ratio of consecutive terms is $(n+1)/x$. For
$x = 5$, the term magnitudes (times $e^{-x}/x$) run $1, 0.2, 0.08,
0.048, 0.0384, 0.0384, 0.046, \ldots$: they fall to a minimum near
$n = 5$ and then climb. The best truncation stops at the smallest term,
near $n \approx x$, and the achievable error is roughly that smallest
term, about $e^{-x}\sqrt{2\pi/x}$ by Stirling, which for $x = 10$ is at
the $10^{-5}$ level. Adding more terms past the minimum makes the
answer worse, then catastrophically wrong. An engineer who treats the
asymptotic series like a convergent one, summing as many terms as
patience allows, computes a number that diverges from the truth. The
rule, stated in the mastery box of section 4.4, is to stop at the
smallest term. The detailed analysis of optimal truncation and
superasymptotics is in \cite{text:v2ch04-bender-orszag-1999}.

\subsection{Catastrophic cancellation in series evaluation}

A convergent series, summed in the wrong order or with the wrong sign
structure, can lose all precision to round-off even though the
mathematics is sound. The example is $e^{-x}$ for large positive $x$
evaluated by its Maclaurin series $\sum (-x)^{n}/n!$ directly. At $x =
20$, the largest terms reach magnitude $20^{20}/20! \approx 4.3 \times
10^{7}$, while the true value $e^{-20} \approx 2.06 \times 10^{-9}$ is
seventeen orders of magnitude smaller. The sum is a tower of large
terms of alternating sign that must cancel to sixteen digits to leave
the tiny result, and double precision carries only about sixteen
significant digits in total. The cancellation consumes every available
digit, and the computed answer is dominated by round-off, often wrong in
the first significant figure or returning a negative number where the
true value is positive. The fix is algebraic: compute $e^{+20}$ by the
all-positive series, where no cancellation occurs, and take the
reciprocal. Series convergence is necessary for a correct sum and is not
sufficient for a correct \emph{computed} sum; the floating-point
arithmetic between the terms is the other half of the engineering. The
script \texttt{code/cancellation\_exp.py} tabulates both routes
(\texttt{data/cancellation\_exp.csv}): at $x = 20$ the direct series
returns a relative error near $200\%$ while the reciprocal route is
correct to machine precision.

\subsection{The principle}

\begin{principle}[Shrinking terms do not guarantee a convergent sum]
For a series $\sum a_{n}$, the condition $a_{n} \to 0$ is necessary
for convergence but not sufficient. A series whose partial sums
appear to be settling may still diverge if the terms shrink no faster
than $1/n$. Before truncating an infinite series at finite $N$ and
treating the tail as negligible, the engineer verifies convergence by
the comparison, ratio, root, or integral test, and bounds the
truncation tail explicitly.
\end{principle}

The same principle reappears later in the volume in the contexts of
Fourier series convergence, asymptotic expansions, perturbation
theory, and improper integrals. The harmonic-series trap is the
chapter's first encounter with a pattern the engineer will meet again.

\section{Case study: computing a constant to a guaranteed accuracy}\pathtag{standard}

The chapter's machinery, convergence verdicts, truncation bounds,
acceleration, and round-off discipline, comes together when an engineer
must compute a transcendental constant to a stated accuracy with a
defensible error budget. We take the concrete task of computing $\pi$ to
ten correct decimal places and carry it from the choice of series
through the final certified digits, treating every error source the way
a measurement engineer treats an uncertainty budget in Volume I. The
constant is chosen because the answer is known, so every claimed bound
can be checked, and because the same discipline transfers verbatim to
constants whose values are not known in advance: a material property
defined by an infinite series, a special-function value in a control
law, a normalisation integral in a probability model.

\subsection{Step one: reject the naive series}

The first arctangent series a reader reaches for is the Leibniz series
$\pi/4 = \arctan 1 = 1 - 1/3 + 1/5 - 1/7 + \cdots$, evaluated at $x =
1$. It converges, by the alternating series test, and its truncation
error after $N$ terms is bounded by the first omitted term $1/(2N+1)$.
To reach ten decimals the engineer needs $1/(2N+1) \leq 5 \times
10^{-11}$, that is $N \geq 10^{10}$ terms. Ten billion terms for ten
digits is the harmonic-rate trap wearing a respectable coefficient: the
series converges, but at $O(1/N)$, and the cost is prohibitive. The
term test and the alternating-series bound together deliver the verdict
in one line, and the verdict is to find a faster series before writing
any code.

\subsection{Step two: choose a fast-converging representation}

The arctangent series $\arctan x = x - x^{3}/3 + x^{5}/5 - \cdots$
converges geometrically at rate $x^{2}$ for $\abs{x} < 1$, so a small
argument converges fast. The classical device is a \engterm{Machin-like
formula} that writes $\pi/4$ as a combination of arctangents of small
rational arguments. Machin's 1706 identity is
\[
\frac{\pi}{4} = 4 \arctan\frac{1}{5} - \arctan\frac{1}{239}.
\]
The identity is verified by the tangent addition formula: $\tan(4
\arctan\tfrac{1}{5})$ works out to $120/119$, and $\tan(4
\arctan\tfrac{1}{5} - \arctan\tfrac{1}{239}) = 1$, so the combination
equals $\pi/4$. The two arguments $1/5$ and $1/239$ are small, so each
arctangent series converges at geometric rate $1/25$ and $1/57121$
respectively. The choice of representation, not the choice of
programming language, is what turns a $10^{10}$-term computation into a
fifteen-term one.

\subsection{Step three: the truncation budget, term by term}

The error budget allots the total tolerance across the two series. We
demand each arctangent be computed to an absolute error below $2 \times
10^{-11}$, so the combined error in $\pi/4$ stays under $5 \times
10^{-11}$ even after the factor of four on the first term, and $\pi$
itself lands within $2 \times 10^{-10}$, comfortably inside the
ten-decimal target with margin for the round-off allotment of step four.

For $\arctan(1/5)$, the alternating series has $k$-th term magnitude
$(1/5)^{2k+1}/(2k+1)$. The first omitted term after keeping terms
through $k = K$ is $(1/5)^{2K+3}/(2K+3)$. Requiring this below $2 \times
10^{-11}/4 = 5 \times 10^{-12}$ (the factor of four because this
arctangent is multiplied by four): at $K = 7$ the next term is
$(1/5)^{17}/17 = 1/(7.6 \times 10^{11} \cdot 17) \approx 7.7 \times
10^{-14}$, comfortably below budget; at $K = 6$ it is $(1/5)^{15}/15
\approx 2.2 \times 10^{-12}$, marginal. Eight terms ($k = 0$ through
$7$) settle $\arctan(1/5)$ within budget.

For $\arctan(1/239)$, the geometric rate $1/239^{2} \approx 1.75 \times
10^{-5}$ is so fast that the first term $1/239 \approx 4.18 \times
10^{-3}$ is already most of the answer. The next term is
$(1/239)^{3}/3 \approx 2.4 \times 10^{-8}$, and the term after that is
$(1/239)^{5}/5 \approx 8.6 \times 10^{-14}$, below budget. Two terms
($k = 0, 1$) settle $\arctan(1/239)$, because the bound on the first
omitted term at $k = 2$ already clears the allotment. The asymmetry in
term counts, eight against two, is the geometric rate at work: the
smaller argument buys a much faster series.

\begin{figure}[ht]
\centering
\begin{tikzpicture}[x=1cm, y=1cm, font=\small]

% Term magnitudes of the two arctan series, log scale, vs term index k.
% arctan(1/5): term_k = (1/5)^(2k+1)/(2k+1). log10 values:
%  k=0: 1/5=0.2 -> -0.70; k=1: (0.2)^3/3=2.67e-3 -> -2.57; k=2: (0.2)^5/5=6.4e-5 -> -4.19;
%  k=3: (0.2)^7/7=1.83e-6 -> -5.74; k=4: 5.69e-8 -> -7.24; k=5: 1.86e-9 -> -8.73;
%  k=6: 6.3e-11 -> -10.2; k=7: 2.17e-12 -> -11.66
% arctan(1/239): term_k = (1/239)^(2k+1)/(2k+1):
%  k=0: 4.18e-3 -> -2.38; k=1: 2.44e-8 -> -7.61; k=2: 8.6e-14 -> -13.07
% y scale: log10 from 0 (top) to -14 (bottom). y_cm = (val+14)*0.55 ; map.

% Axes (y from -14 bottom to 0 top, total 14 decades; use y_cm = val (negative) shifted)
\draw[->, thick] (-0.3, -14) -- (8.0, -14) node[right] {term index $k$};
\draw[->, thick] (-0.3, -14.3) -- (-0.3, 0.8) node[above] {term magnitude (log scale)};

% y ticks every 2 decades
\foreach \p in {0,-2,-4,-6,-8,-10,-12,-14} {
  \draw (-0.45, \p) -- (-0.15, \p) node[left, xshift=-0.1cm] {$10^{\p}$};
}
% x ticks at k=0..7 mapped x_cm = k*1.0
\foreach \k in {0,1,2,3,4,5,6,7} {
  \draw (\k, -14.15) -- (\k, -13.85) node[below, yshift=-0.1cm] {\k};
}

% per-series tolerance line at 5e-12 -> log10 = -11.30
\draw[dashed, gray] (-0.3, -11.30) -- (7.6, -11.30)
  node[right, gray, yshift=0.18cm] {budget $5\times10^{-12}$};

% arctan(1/5) series
\draw[thick, engblue, mark=*, mark size=1pt] plot coordinates {
  (0,-0.70) (1,-2.57) (2,-4.19) (3,-5.74) (4,-7.24) (5,-8.73)
  (6,-10.2) (7,-11.66) };
\node[engblue, anchor=west] at (5.1, -8.0) {$\arctan\tfrac{1}{5}$};

% arctan(1/239) series
\draw[thick, engred, mark=square*, mark size=1pt] plot coordinates {
  (0,-2.38) (1,-7.61) (2,-13.07) };
\node[engred, anchor=west] at (2.1, -13.07) {$\arctan\tfrac{1}{239}$};

\end{tikzpicture}
\caption[Truncation budget for the Machin computation of $\pi$]{Term
magnitudes of the two arctangent series in Machin's formula $\pi/4 = 4
\arctan\tfrac{1}{5} - \arctan\tfrac{1}{239}$, on a log scale against term
index. The dashed line is the per-series truncation budget. Each series is
cut once the next-term bound drops below the line: at index seven for the
slowly converging $1/5$ series, at index one for the fast $1/239$ series
(whose index-two term already sits far below budget). The asymmetry is the
geometric convergence rate at work, $1/25$ per step against $1/57121$ per
step. The figure is the case study's error budget made visual.}
\label{fig:vol02:ch04:machin-budget}
\end{figure}


The budget figure (\autoref{fig:vol02:ch04:machin-budget}) plots the
term magnitudes of the two arctangent series on a log scale against term
index, with the per-series tolerance drawn as a horizontal line. The
term where each curve crosses below the line is the truncation point,
read directly off the figure: term seven for the $1/5$ series, term two
for the $1/239$ series. The figure is the case study's error budget made
visual, the same way a Volume I uncertainty bar chart shows which
component dominates.

\subsection{Step four: the round-off budget}

Truncation is half the budget; floating-point round-off is the other
half, and a computation that controls only the truncation error fails
the moment the arithmetic between the terms loses digits. Two features of
the Machin computation keep round-off small. First, the terms are all the
same sign within each arctangent series after factoring out the
alternation (the series is summed as a difference of grouped positive
quantities, with no catastrophic cancellation of the kind that wrecked
the $e^{-20}$ computation of the failure section). Second, the terms are
summed smallest-first, so the running sum accumulates the tiny tail terms
before they are swamped by the leading term, which preserves the
low-order bits that smallest-last summation discards.

In IEEE double precision, about sixteen significant decimals are
available, and the dominant round-off in summing fifteen terms is at
the level of the largest term times the unit round-off $2^{-53} \approx
1.1 \times 10^{-16}$, accumulated over the term count, so on the order
of $15 \times 4 \times 10^{-3} \times 1.1 \times 10^{-16} \approx 7
\times 10^{-18}$. The round-off contribution is six decades below the
truncation budget, so double precision is more than adequate for ten
decimals; the truncation error dominates the budget, which is the
healthy regime. An engineer chasing thirty decimals would exhaust double
precision and would switch to an extended-precision arithmetic, at which
point the round-off allotment would have to be re-budgeted against the
new unit round-off.

\subsection{Step five: assemble and certify}

Summing the budgeted terms:
\begin{align*}
\arctan\tfrac{1}{5} &= 0.197395559849\ldots, \\
\arctan\tfrac{1}{239} &= 0.004184076002\ldots, \\
\tfrac{\pi}{4} &= 4(0.197395559849) - 0.004184076002 =
0.785398163397\ldots, \\
\pi &= 3.141592653590\ldots
\end{align*}
The computed value agrees with the reference $\pi =
3.1415926535897932\ldots$ to twelve decimals, two better than the
ten-decimal target, the margin the budget reserved. The certificate is
the budget itself: truncation bounded below $2 \times 10^{-10}$ by the
alternating-series bounds of step three, round-off bounded below $10^{-17}$
by the analysis of step four, and the two summed to a total error bound
of $2 \times 10^{-10}$, inside the target. The script
\texttt{code/machin\_pi.py} runs the computation, prints each term, the
running partial sums, and the empirical error against a high-precision
reference, and writes the budget table to
\texttt{data/machin\_pi\_budget.csv}.

\begin{principle}[A computed constant carries two error budgets]
The accuracy of a constant computed from an infinite series is governed
by two independent budgets that must both be controlled: the truncation
error, set by where the series is cut and bounded by a convergence-test
tail bound, and the round-off error, set by the arithmetic precision and
the summation order. A computation that controls only one budget can fail
the other silently. The engineer states both bounds, sums them, and
checks the total against the target before trusting a single digit.
\end{principle}

The case study rehearses the whole chapter on one task. The naive series
was rejected by the term test and the alternating-series bound. The fast
series was chosen for its geometric rate. The truncation point was set by
the per-term bound against an allotted budget. The round-off was bounded
by the summation-order discipline of the failure section. The final
digits were certified by the sum of the two budgets, not by comparison
with a known answer, which is the only certificate available when the
constant's value is not known in advance.

\section{Worked examples: small-angle approximations and perturbation expansions}\pathtag{core}

We work two extended examples, each producing a usable formula and an
explicit error bound.

\subsection{The small-angle approximations}

For $\theta$ small in radians, the Taylor series of $\sin$, $\cos$, and
$\tan$ at $\theta = 0$ give:
\begin{align*}
\sin \theta &= \theta - \frac{\theta^{3}}{6} + \frac{\theta^{5}}{120}
- \cdots, \\
\cos \theta &= 1 - \frac{\theta^{2}}{2} + \frac{\theta^{4}}{24}
- \cdots, \\
\tan \theta &= \theta + \frac{\theta^{3}}{3} + \frac{2 \theta^{5}}{15}
+ \cdots.
\end{align*}
Truncating at the leading term gives the standard approximations:
\[
\sin \theta \approx \theta, \quad \cos \theta \approx 1, \quad
\tan \theta \approx \theta,
\]
together with the next-order refinements
\[
\sin \theta \approx \theta - \frac{\theta^{3}}{6}, \quad \cos \theta
\approx 1 - \frac{\theta^{2}}{2}, \quad \tan \theta \approx \theta
+ \frac{\theta^{3}}{3}.
\]

The truncation error for the leading-term sine is bounded by Taylor's
theorem (Lagrange form) by $M_{3} \theta^{3}/3!$ where $M_{3}$ bounds
$\abs{f^{(3)}}$ on the interval. For $\sin$, $f^{(3)} = -\cos$, so
$M_{3} \leq 1$, and the bound becomes
\[
\abs{\sin \theta - \theta} \leq \frac{\theta^{3}}{6}.
\]
The bound is sharp at $\theta \to 0$. We tabulate it (with all
quantities in radians, and the approximated value $\theta$ given in
degrees alongside for orientation) at the four working angles named in
the project (\autoref{tab:smallangle}).

\begin{table}[h]
\centering
\begin{tabular}{lllll}
\toprule
$\theta$ (deg) & $\theta$ (rad) & $\sin \theta$ & $\theta$ (approx) &
$\abs{\sin \theta - \theta}/\sin \theta$ \\
\midrule
$5^{\circ}$ & $0.0873$ & $0.0872$ & $0.0873$ & $0.13\%$ \\
$10^{\circ}$ & $0.1745$ & $0.1736$ & $0.1745$ & $0.51\%$ \\
$25^{\circ}$ & $0.4363$ & $0.4226$ & $0.4363$ & $3.2\%$ \\
$45^{\circ}$ & $0.7854$ & $0.7071$ & $0.7854$ & $11.1\%$ \\
\bottomrule
\end{tabular}
\caption{The leading-order small-angle sine approximation $\sin \theta
\approx \theta$, with relative error at four working angles. The error
grows as $\theta^{2}/6$ in the relative form, so the same table at
$\theta = 1^{\circ}$ would give roughly $0.005\%$ error.}
\label{tab:smallangle}
\end{table}

The cosine error follows the same pattern with $\theta^{2}/2$ as the
leading term and $\theta^{4}/24$ as the next correction. At
$\theta = 25^{\circ}$, the relative error in $\cos \theta \approx 1$ is
roughly $9.4\%$, and the relative error in $\cos \theta \approx 1 -
\theta^{2}/2$ drops to about $0.4\%$. At $\theta = 45^{\circ}$, the
leading-order cosine approximation is off by $29\%$, and the
quadratic-correction approximation is off by $4.0\%$.

The engineering decision the table makes operational is at what angle
the leading-order approximation has been pushed past its budget. For a
$1\%$ error budget on the sine, the leading-order approximation works
at $\theta \leq 14^{\circ}$ approximately; the quadratic correction
extends the budget to $\theta \leq 35^{\circ}$. For a $0.1\%$ budget,
the leading-order approximation is good only to about $4.5^{\circ}$.
The reader who carries a tolerance into the small-angle decision can
look up the bound rather than guessing.

\begin{figure}[ht]
\centering
\begin{tikzpicture}[x=1cm, y=1cm, font=\small]

% Plot |sin x - T_n(x)| on log scale vs x in [0, 2]
% Axes
\draw[->, thick] (-0.2, 0) -- (8.5, 0) node[right] {$\theta$ (rad)};
\draw[->, thick] (0, -0.2) -- (0, 6.2) node[above] {$|\sin\theta - T_n(\theta)|$ (log scale)};

% x-ticks at 0.25, 0.5, 1.0, 1.5, 2.0 mapped to 1, 2, 4, 6, 8 cm
\foreach \x/\lbl in {1/0.25, 2/0.5, 4/1.0, 6/1.5, 8/2.0} {
  \draw (\x, -0.1) -- (\x, 0.1) node[below, yshift=-0.15cm] {\lbl};
}
% y-ticks: log10 decades, y = 6 at 10^0, y=5 at 10^-1, etc., down to y=0 at 10^-6
\foreach \y/\lbl in {0/$10^{-6}$, 1/$10^{-5}$, 2/$10^{-4}$, 3/$10^{-3}$, 4/$10^{-2}$, 5/$10^{-1}$, 6/$10^{0}$} {
  \draw (-0.1, \y) -- (0.1, \y) node[left, xshift=-0.15cm] {\lbl};
}

% T_1 (linear): error ~ theta^3/6. log10(theta^3/6) at theta values, mapped to y = 6 + log10(...)
% theta=0.25 -> theta^3/6 = 0.0026 -> log10 = -2.58 -> y = 6 - 2.58 = 3.42; but our scale is 1 cm per decade from y=6 (10^0) down to y=0 (10^-6)
% So y = 6 + log10(value).  -2.58 -> 3.42
\draw[thick, engblue] plot coordinates {
  (0.5, 1.40)   % theta=0.125: theta^3/6 = 3.26e-4 -> log = -3.49 -> y = 2.51; but at x=0.5 mapped to theta=0.125
  (1, 2.40)     % theta=0.25: 2.60e-3 -> -2.58 -> y=3.42, but plotting steeper... use real T_1 errors below
};
% Recompute T1 error = |sin theta - theta| exact values
% theta=0.1: 1.67e-4 -> y = 6 - 3.78 = 2.22; x at 0.1 -> x_cm = 0.4
% theta=0.25: 2.60e-3 -> y = 3.42; x_cm=1
% theta=0.5: 2.06e-2 -> y = 4.31; x_cm=2
% theta=1.0: 0.1585 -> y = 5.20; x_cm=4
% theta=1.5: 0.5025 -> y = 5.70; x_cm=6
% theta=2.0: 1.091 -> y = 6.04 (clip at 6); x_cm=8

\draw[thick, engblue, mark=*, mark size=1pt] plot coordinates {
  (0.4, 2.22) (1, 3.42) (2, 4.31) (4, 5.20) (6, 5.70) (8, 6.04)
};
\node[engblue, anchor=west] at (8.1, 6.04) {$T_1$};

% T_3 (cubic): error = |sin theta - (theta - theta^3/6)|
% theta=0.1: 8.33e-8 -> y = -1.08; off-chart at bottom
% theta=0.25: 8.13e-6 -> y = 0.91; x_cm=1
% theta=0.5: 2.59e-4 -> y = 2.41; x_cm=2
% theta=1.0: 8.33e-3 -> y = 3.92; x_cm=4
% theta=1.5: 6.46e-2 -> y = 4.81; x_cm=6
% theta=2.0: 2.42e-1 -> y = 5.38; x_cm=8
\draw[thick, engred, mark=square*, mark size=1pt] plot coordinates {
  (1, 0.91) (2, 2.41) (4, 3.92) (6, 4.81) (8, 5.38)
};
\node[engred, anchor=west] at (8.1, 5.38) {$T_3$};

% T_5 (quintic): error theta=0.5: 1.55e-6 -> y=0.19; theta=1: 1.96e-4 -> y=2.29; theta=2: 0.02 -> y=4.30
\draw[thick, engblack, mark=triangle*, mark size=1pt] plot coordinates {
  (2, 0.19) (4, 2.29) (6, 3.43) (8, 4.30)
};
\node[engblack, anchor=west] at (8.1, 4.30) {$T_5$};

% Tolerance line at 1e-3
\draw[dashed, gray] (0, 3) -- (8.3, 3) node[right, gray] {$10^{-3}$};

\end{tikzpicture}
\caption[Taylor-truncation error for $\sin\theta$ at orders 1, 3, 5]{Absolute
truncation error $|\sin\theta - T_n(\theta)|$ for Taylor polynomials of
$\sin\theta$ at $0$ of degrees $1$, $3$, and $5$. Each curve has slope
$n+1$ in log-log when $\theta$ is small, matching the Lagrange remainder
bound $\theta^{n+1}/(n+1)!$. The dashed line at $10^{-3}$ shows the
budget the project's small-angle rule has to meet: $T_1$ holds it to
$\theta \approx 0.18$ rad ($10^{\circ}$); $T_3$ to $\theta \approx 1.0$
rad ($57^{\circ}$); $T_5$ comfortably to $\theta = 2$ rad. Each
additional order extends the safe range by roughly a factor of $\sqrt[3]{}$
to $\sqrt[5]{}$ of the tolerance ratio.}
\label{fig:vol02:ch04:taylor-error}
\end{figure}


The Taylor-error figure (\autoref{fig:vol02:ch04:taylor-error})
expresses the same trade visually, across three orders. The
log-log slopes match the Lagrange remainder bound: $T_{1}$ has slope
$3$, $T_{3}$ has slope $5$, $T_{5}$ has slope $7$. Each extra order
buys an asymptotically steeper error curve at small $\theta$,
trading two factorial-divided coefficients for the next derivative
bound. The factorial in the denominator dominates as the order
climbs, which is why $T_{9}$ of $\sin$ on $[0, 2]$ is already at
the noise floor of double-precision arithmetic.

\subsection{Perturbation expansion: a quadratic equation with a small parameter}

A standard pattern in engineering analysis: a system has a clean
solution when a small parameter is set to zero, and the engineer wants
the correction to first order in the small parameter. Consider the
quadratic equation
\[
x^{2} + 2 \varepsilon x - 1 = 0
\]
with $\varepsilon$ small. The unperturbed equation ($\varepsilon = 0$)
is $x^{2} = 1$, with roots $x = \pm 1$. We look for the perturbed root
near $x = 1$ in the form
\[
x(\varepsilon) = 1 + a_{1} \varepsilon + a_{2} \varepsilon^{2}
+ a_{3} \varepsilon^{3} + \cdots.
\]
Substitute and expand:
\[
(1 + a_{1} \varepsilon + a_{2} \varepsilon^{2} + \cdots)^{2}
+ 2 \varepsilon (1 + a_{1} \varepsilon + \cdots) - 1 = 0.
\]
Collect by powers of $\varepsilon$. The constant term: $1 - 1 = 0$,
satisfied. The $\varepsilon^{1}$ term: $2 a_{1} + 2 = 0$, giving $a_{1}
= -1$. The $\varepsilon^{2}$ term: $2 a_{2} + a_{1}^{2} + 2 a_{1} = 0$,
giving $a_{2} = (2 a_{1} - a_{1}^{2})/(-2) \cdot (-1) = (2 - 1)/2$;
correctly, $a_{2} = 1/2$. The expansion to second order is
\[
x_{+}(\varepsilon) = 1 - \varepsilon + \frac{\varepsilon^{2}}{2} +
O(\varepsilon^{3}).
\]
The exact root is $x_{+}(\varepsilon) = -\varepsilon + \sqrt{1 +
\varepsilon^{2}}$, which by binomial expansion gives $-\varepsilon + 1
+ \varepsilon^{2}/2 - \varepsilon^{4}/8 + \cdots$, in agreement.

The same procedure produces $x_{-}(\varepsilon) = -1 - \varepsilon -
\varepsilon^{2}/2 + O(\varepsilon^{3})$ for the root near $-1$. The
engineering payoff: a closed-form first-order correction without
solving the quadratic at every value of $\varepsilon$. This pattern
generalises across nonlinear equations, ODEs, and PDEs, and underwrites
much of the asymptotic-methods machinery later in the volume
\cite{text:mahajan2014}.

\subsection{The pendulum period as a worked series expansion}

The third worked example produces the amplitude correction the project
asks the reader to verify. The exact period of a simple pendulum at
amplitude $\theta_{0}$ involves the complete elliptic integral of the
first kind,
\[
T(\theta_{0}) = 4 \sqrt{\frac{L}{g}} \int_{0}^{\pi/2}
\frac{\dd \phi}{\sqrt{1 - k^{2} \sin^{2} \phi}}, \qquad k =
\sin\frac{\theta_{0}}{2}.
\]
Expand the integrand by the binomial series $(1 - u)^{-1/2} = 1 +
\tfrac{1}{2} u + \tfrac{3}{8} u^{2} + \cdots$ with $u = k^{2} \sin^{2}
\phi$:
\[
\frac{1}{\sqrt{1 - k^{2}\sin^{2}\phi}} = 1 + \frac{1}{2} k^{2}
\sin^{2}\phi + \frac{3}{8} k^{4} \sin^{4}\phi + \cdots.
\]
Integrate term by term over $[0, \pi/2]$, using the Wallis integrals
$\int_{0}^{\pi/2} \sin^{2}\phi \, \dd\phi = \pi/4$ and
$\int_{0}^{\pi/2} \sin^{4}\phi \, \dd\phi = 3\pi/16$:
\[
\int_{0}^{\pi/2} \frac{\dd\phi}{\sqrt{1 - k^{2}\sin^{2}\phi}} =
\frac{\pi}{2}\left(1 + \frac{1}{4} k^{2} + \frac{9}{64} k^{4} +
\cdots\right).
\]
Therefore $T = 2\pi\sqrt{L/g}\,(1 + \tfrac{1}{4} k^{2} + \tfrac{9}{64}
k^{4} + \cdots)$. Now expand $k = \sin(\theta_{0}/2) = \theta_{0}/2 -
\theta_{0}^{3}/48 + \cdots$, so $k^{2} = \theta_{0}^{2}/4 -
\theta_{0}^{4}/48 + \cdots$ and $k^{4} = \theta_{0}^{4}/16 + \cdots$.
Substituting and collecting by power of $\theta_{0}$:
\[
T(\theta_{0}) = 2\pi\sqrt{\frac{L}{g}}\left(1 + \frac{\theta_{0}^{2}}{16}
+ \frac{11\,\theta_{0}^{4}}{3072} + \cdots\right),
\]
matching the formula quoted in the project. The leading correction
$\theta_{0}^{2}/16$ says a pendulum released from $23^{\circ}$ runs
about $1\%$ slow relative to the small-amplitude period, the kind of
systematic error a pendulum-clock designer compensates with the escape
mechanism. The worked expansion combines three of the chapter's tools:
the binomial series, term-by-term integration inside the radius of
convergence, and composition of one series into another.

\section{Project}\pathtag{core}

\begin{project}[Small-angle approximations with explicit error bounds]
The reader derives the standard small-angle approximations from their
Taylor series, computes explicit error bounds at four working angles,
compares to the exact values, and writes a 1500-word reflection on
when each approximation is safe to use and when it should be replaced
with the exact form.

\textbf{Track}. Analysis only. \textbf{Hazard class}: none.
\textbf{Tools}. Paper, pencil, calculator. No computer or simulation
required, although the reader may use a computer-algebra system or
spreadsheet to confirm the arithmetic.

\textbf{Procedure}.

\begin{enumerate}[leftmargin=*]
  \item \engterm{Sine and cosine}. Derive the leading-order
    approximations $\sin \theta \approx \theta$ and $\cos \theta
    \approx 1$, and their next-order refinements $\sin \theta \approx
    \theta - \theta^{3}/6$ and $\cos \theta \approx 1 - \theta^{2}/2$,
    by writing out the Maclaurin series and truncating at the chosen
    order. State the Lagrange remainder bound for each approximation,
    using $M = 1$ for the bound on the next derivative magnitude.
  \item \engterm{Tangent}. Derive $\tan \theta \approx \theta$ and
    $\tan \theta \approx \theta + \theta^{3}/3$ from the Maclaurin
    series of the tangent. State the radius of convergence of the
    series and identify the obstacle to a global Taylor expansion of
    $\tan$.
  \item \engterm{Pendulum-period correction}. The exact period of a
    simple pendulum at amplitude $\theta_{0}$ is
    \[
    T(\theta_{0}) = 4 \sqrt{\frac{L}{g}} \int_{0}^{\pi/2}
    \frac{\dd \phi}{\sqrt{1 - \sin^{2}(\theta_{0}/2) \sin^{2} \phi}}.
    \]
    The integral on the right is the complete elliptic integral of
    the first kind, $K(\sin(\theta_{0}/2))$. Expand the integrand in
    powers of $\sin(\theta_{0}/2)$ and integrate term by term to
    obtain
    \[
    T(\theta_{0}) = 2 \pi \sqrt{\frac{L}{g}}\left(1
    + \frac{\theta_{0}^{2}}{16} + \frac{11 \theta_{0}^{4}}{3072}
    + \cdots\right).
    \]
    The reader may take the elliptic-integral identity as given (a
    standard result) and verify the first two terms of the expansion;
    the derivation of the elliptic-integral form belongs to Volume II
    Chapter 6.
  \item \engterm{Relativistic kinetic energy at low speed}. The
    relativistic kinetic energy is $K = (\gamma - 1) m c^{2}$ with
    $\gamma = 1/\sqrt{1 - v^{2}/c^{2}}$. Expand $\gamma$ in powers of
    $\beta^{2} = v^{2}/c^{2}$ and show that, to leading order in
    $\beta^{2}$,
    \[
    K \approx \frac{1}{2} m v^{2} + \frac{3}{8} m v^{2} \beta^{2}
    + \cdots.
    \]
    The leading-order term is the Newtonian kinetic energy. The
    next-order term is the leading relativistic correction. Identify
    the speed at which the relativistic correction is $1\%$ of the
    Newtonian result.
  \item \engterm{Binomial expansion}. Derive
    \[
    (1 + x)^{1/2} \approx 1 + \frac{x}{2} - \frac{x^{2}}{8}
    + \frac{x^{3}}{16} - \cdots, \qquad
    (1 + x)^{-1} \approx 1 - x + x^{2} - x^{3} + \cdots,
    \]
    from the binomial theorem with non-integer $\alpha$. State the
    radius of convergence ($\abs{x} < 1$) and identify three
    engineering quantities for which one of these expansions is the
    standard linearisation tool.
\end{enumerate}

\textbf{Error tabulation}. For each approximation, compute the
absolute error at $\theta = 5^{\circ}, 10^{\circ}, 25^{\circ},
45^{\circ}$ (where applicable, and at the corresponding $\beta$ values
for the relativistic case). Tabulate both the leading-order error and
the next-order error.

\textbf{Reflection (1500 words)}. The reader writes a 1500-word
narrative addressing:

\begin{itemize}[leftmargin=*]
  \item Which small-angle approximation has the most generous error
    budget, and why. (The answer turns on which derivative bound is
    smallest, given that $\sin$, $\cos$, and $\tan$ have very
    different higher-derivative growth.)
  \item For each approximation, the angle at which the leading-order
    formula has $1\%$ relative error and the angle at which it has
    $5\%$ error. State whether these crossover angles match what the
    reader had implicitly assumed before doing the calculation.
  \item One engineering context in which the leading-order
    approximation is reliable, one in which the next-order correction
    is required, and one in which the exact form should be used. For
    each, name the regime and the dominant uncertainty source.
  \item Whether the small-angle approximation is, in the reader's
    judgment, a piece of mathematical convenience or an irreducible
    feature of the engineering work. Defend the position with a
    specific example.
  \item One measurement or analysis in the reader's prior experience
    where a small-angle (or analogous small-parameter) approximation
    was applied. Identify the implicit error tolerance the writer
    appears to have had, whether the tolerance was met, and what the
    next-order correction would have changed.
\end{itemize}

\textbf{Deliverable}. The five derivations, the error table at the
specified angles, and the reflection. The general editor's reference
set of derivations and tables is published online; the reader is
encouraged to compare \emph{after} completing their own.

\textbf{Time}. Approximately five to seven hours of derivation, two
to three hours of error tabulation, and three to four hours of
reflection. Total approximately ten to fourteen hours, distributed
across two or three sittings.

\textbf{Phases and success criteria}.

\begin{enumerate}[leftmargin=*]
  \item \emph{Series derivation phase} (3--4 hr). Success criterion:
    each of the five Maclaurin expansions is written out to the
    stated order, with the Lagrange remainder bound stated in
    closed form and the $M_{n+1}$ constant identified. Common
    failure: writing the series and stopping without the remainder
    bound. The reader who omits the bound has the formula but not
    the licence to use it inside a tolerance budget.
  \item \emph{Pendulum and relativistic-energy phase} (2--3 hr).
    Success criterion: the elliptic-integral identity is taken as
    given and the first two terms of the period expansion are
    verified by direct substitution; the relativistic correction
    is derived from $(1 - \beta^{2})^{-1/2}$ via the binomial
    series and the leading-order Newtonian limit is recovered.
    Common failure: confusing $\beta = v/c$ with $\beta^{2}$ or
    dropping the $1/2$ that converts $\gamma - 1$ to the kinetic
    energy.
  \item \emph{Error tabulation phase} (2--3 hr). Success criterion:
    the absolute and relative errors at $5^{\circ}, 10^{\circ},
    25^{\circ}, 45^{\circ}$ match the values in
    \autoref{tab:smallangle} to three significant figures for the
    sine, with analogous tables for cosine, tangent, pendulum
    period, and relativistic energy. Common failure: forgetting
    to convert degrees to radians before evaluating the series
    (the formulas all expect radians).
  \item \emph{Reflection phase} (3--4 hr). Success criterion: each
    of the five bullets in the reflection list receives at least
    one paragraph of concrete content, with a named example or
    measurement. Common failure: writing the reflection in the
    abstract, without an example. A reflection that names no
    specific engineering context has not engaged the question.
\end{enumerate}

\textbf{What the writer is meant to take away}. The technical
machinery (five expansions, four error tables, two physical
applications) is the proximate output. The deeper takeaway is the
discipline of carrying an explicit error budget through every
approximation, deciding the truncation order from the budget rather
than from convention, and recognising when the budget has been
exceeded. The reflection earns its place by forcing the writer to
state where on their own working surface this discipline has not yet
been internalised, and what would change if it were.
\end{project}

\section{Exercises}

The exercises train calculation, derivation, estimation, simulation,
design, diagnosis, failure analysis, and judgment. Calculation and
derivation problems have full solutions; open-ended problems have
rubrics or reference write-ups. Exercises are tagged with their
reader-path tier.

\subsection*{Calculation}\pathtag{core}

\begin{exercise}
Compute the sum of the geometric series
$\sum_{n=0}^{\infty} (0.9)^{n}$ and the partial sum
$\sum_{n=0}^{19} (0.9)^{n}$. State the absolute and relative
truncation error.

\paragraph{Solution sketch.} The infinite sum is $1/(1 - 0.9) = 10$.
The finite sum is $(1 - 0.9^{20})/(1 - 0.9) = 10 (1 - 0.9^{20})$.
Computing $0.9^{20} \approx 0.1216$ gives $S_{20} \approx 8.784$.
The truncation tail equals the rest of the geometric series starting
at $n = 20$, which sums to $0.9^{20}/(1 - 0.9) = 10 \cdot 0.9^{20}
\approx 1.216$. The absolute error is approximately $1.216$ and the
relative error is approximately $12.16\%$. The geometric formula
$\text{tail} = r^{N}/(1 - r) \cdot S_{0}$ is exact for the geometric
series and is the cheap way to obtain truncation tails when the
series is in geometric form.
\end{exercise}

\begin{exercise}
Find the radius of convergence of the power series
$\sum_{n=1}^{\infty} (n!) \, x^{n} / n^{n}$. State the test you used.

\paragraph{Solution sketch.} Apply the ratio test on the coefficient
magnitudes $c_{n} = n!/n^{n}$:
\[
\frac{c_{n+1}}{c_{n}} = \frac{(n+1)!}{(n+1)^{n+1}} \cdot \frac{n^{n}}{n!}
= \frac{n^{n}}{(n+1)^{n}} = \left(\frac{n}{n+1}\right)^{n} = \left(1 -
\frac{1}{n+1}\right)^{n} \to \frac{1}{e}.
\]
The radius of convergence is $1/L = e$ where $L = \lim c_{n+1}/c_{n}
= 1/e$. The series converges absolutely for $|x| < e$ and diverges
for $|x| > e$. The endpoints $x = \pm e$ require separate testing.
The test used is the ratio test on the coefficients
$c_{n}$, equivalent to applying the ratio test to the absolute terms
$|c_{n} x^{n}|$ at fixed $x$.
\end{exercise}

\begin{exercise}
Compute the first four nonzero terms of the Maclaurin series of
$f(x) = \log(1 + x^{2})$. State the radius of convergence.

\paragraph{Solution sketch.} Substitute $u = x^{2}$ into the
Maclaurin series for $\log(1 + u)$:
\[
\log(1 + u) = u - \frac{u^{2}}{2} + \frac{u^{3}}{3} - \frac{u^{4}}{4}
+ \cdots
\]
Substituting back gives the first four nonzero terms:
\[
\log(1 + x^{2}) = x^{2} - \frac{x^{4}}{2} + \frac{x^{6}}{3} -
\frac{x^{8}}{4} + \cdots
\]
The series for $\log(1 + u)$ converges for $-1 < u \leq 1$; with
$u = x^{2}$ this becomes $0 \leq x^{2} \leq 1$, that is, $|x| \leq 1$
(with strict inequality at $x^{2} = -1$, vacuous in $\R$). The
radius of convergence is $1$, with convergence at both endpoints
$x = \pm 1$ since at those values $u = 1$ and $\log 2$ converges.
\end{exercise}

\begin{exercise}
Use the alternating-series error bound to determine how many terms of
the Maclaurin series of $\sin(0.5)$ are needed to compute the value
to within $10^{-6}$.

\paragraph{Solution sketch.} The alternating-series bound says the
truncation error is at most the magnitude of the first omitted term.
The $n$-th term magnitude (for the $(2n+1)$-th order term) is
$0.5^{2n+1}/(2n+1)!$. Set this $\leq 10^{-6}$. At $n = 3$ the term
is $0.5^{7}/7! = 1/(128 \cdot 5040) \approx 1.55 \times 10^{-6}$,
which fails the bound. At $n = 4$ the term is $0.5^{9}/9! \approx
1.4 \times 10^{-9}$, which satisfies. Therefore four nonzero terms
suffice: $\sin(0.5) \approx 0.5 - 0.5^{3}/6 + 0.5^{5}/120 -
0.5^{7}/5040 \approx 0.47943$. The exact value is $0.479425539$;
the truncation error is below $10^{-6}$ as predicted.
\end{exercise}

\begin{exercise}
Compute $e^{0.1}$ using the Maclaurin series truncated after the
quadratic term, and state the Lagrange remainder bound at $\xi$ in
$[0, 0.1]$.

\paragraph{Solution sketch.} The truncated value is $1 + 0.1 +
0.1^{2}/2 = 1 + 0.1 + 0.005 = 1.105$. The Lagrange remainder bound is
\[
|R_{2}(0.1)| = \left|\frac{e^{\xi}}{3!} (0.1)^{3}\right| \leq
\frac{e^{0.1}}{6} \cdot 10^{-3} \leq \frac{1.106}{6} \cdot 10^{-3}
\approx 1.84 \times 10^{-4}.
\]
The exact value is $e^{0.1} = 1.105171$. The actual error is
$|1.105171 - 1.105| \approx 1.71 \times 10^{-4}$, well within the
bound. A tighter bound uses $e^{\xi} \leq e^{0.1}$ on the interval
$[0, 0.1]$; a looser bound uses $e^{\xi} \leq e^{1} = 2.72$ for any
$\xi \in [0, 1]$.
\end{exercise}

\subsection*{Derivation}\pathtag{standard}

\begin{exercise}
Prove that the harmonic series $\sum_{n=1}^{\infty} 1/n$ diverges by
the grouping argument of Oresme. Then prove that the partial sums
satisfy $S_{N} \geq 1 + (\log_{2} N)/2$ for $N$ a power of $2$.

\paragraph{Solution sketch.} Group the terms into blocks of doubling
size starting after $n = 1$: $\{1/2\}$, $\{1/3, 1/4\}$, $\{1/5, \ldots,
1/8\}$, and so on, with the $k$-th block containing terms from
$1/(2^{k-1} + 1)$ to $1/2^{k}$, that is, $2^{k-1}$ terms each at
least $1/2^{k}$ in size. The block sum is at least $2^{k-1}/2^{k} =
1/2$. The sum up through $N = 2^{m}$ thus contains the initial term
$1$ plus $m$ blocks each contributing at least $1/2$, giving
\[
S_{2^{m}} \geq 1 + \frac{m}{2} = 1 + \frac{\log_{2} N}{2}.
\]
The right-hand side is unbounded as $m \to \infty$, so the partial
sums grow without bound and the series diverges. The bound $1 +
(\log_{2} N)/2$ is loose; the asymptotic $S_{N} \sim \log N + \gamma$
gives a sharper estimate at large $N$, but the loose bound suffices
to prove divergence.
\end{exercise}

\begin{exercise}
Prove the alternating-series test: if $b_{n}$ is monotone decreasing
to zero, the partial sums of $\sum (-1)^{n+1} b_{n}$ form a Cauchy
sequence, and so the series converges. State the truncation-error
bound and prove it.

\paragraph{Solution sketch.} Let $S_{N} = \sum_{n=1}^{N} (-1)^{n+1}
b_{n}$. The differences $S_{2k}$ and $S_{2k+2}$ satisfy $S_{2k+2}
- S_{2k} = b_{2k+1} - b_{2k+2} \geq 0$ by monotonicity, so the even
partial sums are monotone increasing. Similarly the odd partial sums
$S_{2k+1}$ are monotone decreasing. Both are bounded: $S_{2k} \leq
S_{1} = b_{1}$ and $S_{2k+1} \geq S_{2} = b_{1} - b_{2} \geq 0$. By
the monotone convergence theorem each subsequence has a limit. The
gap $S_{2k+1} - S_{2k} = b_{2k+1} \to 0$ forces the two subsequence
limits to coincide. Therefore $S_{N}$ converges to a common limit
$S$.

For the error bound, the limit $S$ sits between any two consecutive
partial sums: $S_{2k} \leq S \leq S_{2k+1}$, hence $|S - S_{N}| \leq
|S_{N+1} - S_{N}| = b_{N+1}$. The truncation error is at most the
magnitude of the first omitted term.
\end{exercise}

\begin{exercise}
Derive Euler's identity $e^{i\theta} = \cos \theta + i \sin \theta$
by substituting $x = i\theta$ into the Maclaurin series for $e^{x}$
and separating real and imaginary parts. State the radius of
convergence.

\paragraph{Solution sketch.} Substituting $x = i\theta$ in the
Maclaurin series gives
\[
e^{i\theta} = \sum_{n=0}^{\infty} \frac{(i\theta)^{n}}{n!}.
\]
The powers $i^{n}$ cycle through $1, i, -1, -i, 1, \ldots$ with
period four. Split the sum by parity of $n$. For even $n = 2k$,
$i^{2k} = (-1)^{k}$ and the contribution is $\sum_{k} (-1)^{k}
\theta^{2k}/(2k)! = \cos\theta$. For odd $n = 2k+1$, $i^{2k+1} = i
\cdot (-1)^{k}$ and the contribution is $i \sum_{k} (-1)^{k}
\theta^{2k+1}/(2k+1)! = i \sin\theta$. Adding gives $e^{i\theta} =
\cos\theta + i \sin\theta$.

The exponential series has infinite radius of convergence on the
complex plane, so the manipulation is justified for every
$\theta \in \R$ (indeed for every complex argument). The
rearrangement of an absolutely convergent series of complex terms
into its real and imaginary subsums uses absolute convergence,
which is precisely the safer condition emphasised in section
4.2.4.
\end{exercise}

\begin{exercise}
Derive the Maclaurin expansion of $1/(1 - x)^{2}$ by differentiating
the geometric series $1/(1 - x) = \sum x^{n}$ term by term. State the
radius of convergence and verify the result against the binomial
expansion.

\paragraph{Solution sketch.} Differentiate term by term:
\[
\frac{\dd}{\dd x}\left[\sum_{n=0}^{\infty} x^{n}\right] =
\sum_{n=1}^{\infty} n x^{n-1} = \sum_{m=0}^{\infty} (m + 1) x^{m},
\]
which equals $1/(1 - x)^{2}$ since $\frac{\dd}{\dd x} 1/(1 - x)
= 1/(1 - x)^{2}$. Therefore
\[
\frac{1}{(1 - x)^{2}} = 1 + 2 x + 3 x^{2} + 4 x^{3} + \cdots
= \sum_{n=0}^{\infty} (n + 1) x^{n}.
\]
Term-by-term differentiation of a power series is valid inside the
radius of convergence, so the new series converges for $|x| < 1$,
the radius of convergence of the original geometric series.

Verification via the binomial expansion. The binomial series for
$(1 - x)^{-2} = (1 + (-x))^{-2}$ has coefficients
$\binom{-2}{n}(-1)^{n} = (-1)^{n} \cdot (-2)(-3) \cdots (-2 - n + 1)
/ n! = (-1)^{n} \cdot (-1)^{n} (n+1)!/n! / 1 = (n + 1)$, matching
the derivative result.
\end{exercise}

\begin{exercise}
Show that the Taylor series of $f(x) = e^{-1/x^{2}}$ at $x_{0} = 0$
(with $f(0) = 0$ by convention) has every coefficient equal to zero,
and yet $f(x) \neq 0$ for $x \neq 0$. Discuss what this example shows
about the relationship between a function and its Taylor series.

\paragraph{Solution sketch.} For $x \neq 0$, $f(x) = e^{-1/x^{2}}$ is
strictly positive but extremely small near zero. By l'H\^{o}pital or
direct substitution, $\lim_{x \to 0} f(x) = 0$, so the extension
$f(0) = 0$ is continuous. Each derivative $f^{(n)}(x)$ for $x \neq 0$
is a rational function of $x$ times $e^{-1/x^{2}}$. The exponential
factor dominates any polynomial in $1/x$, so $\lim_{x \to 0} f^{(n)}(x)
= 0$ for every $n$. Setting $f^{(n)}(0) = 0$ then preserves continuity
of every derivative. The Taylor series at $x_{0} = 0$ is
\[
\sum_{n=0}^{\infty} \frac{f^{(n)}(0)}{n!} x^{n} = 0 + 0 + 0 + \cdots
= 0,
\]
which converges everywhere to the zero function. Yet $f(x) > 0$ for
$x \neq 0$, so the function does not equal its Taylor series anywhere
except at $x = 0$ itself.

\paragraph{Discussion.} A smooth function ($C^{\infty}$) is not the
same as an analytic function (equal to its Taylor series in a
neighbourhood). The example shows that smoothness is strictly weaker
than analyticity. Engineers who substitute "infinitely differentiable"
for "analytic" risk drawing conclusions that depend on the stronger
property. The class of bump functions used to build smooth partitions
of unity is constructed from this counterexample; the price of
having a smooth function with compact support is exactly that the
function is not analytic.
\end{exercise}

\subsection*{Estimation}\pathtag{core}

\begin{exercise}
Estimate the angle below which the small-angle sine approximation
$\sin \theta \approx \theta$ has error below $0.1\%$. Then estimate
the corresponding angle for $0.01\%$ error. State the order-of-magnitude
scaling of the boundary with the error tolerance.

\paragraph{Solution sketch.} The relative error of $\sin\theta
\approx \theta$ at leading order is $\theta^{2}/6$. Setting this
equal to $10^{-3}$ gives $\theta^{2} \approx 6 \times 10^{-3}$, so
$\theta \approx 0.0775$ rad $\approx 4.4^{\circ}$. Setting it equal
to $10^{-4}$ gives $\theta \approx \sqrt{6 \times 10^{-4}} \approx
0.0245$ rad $\approx 1.4^{\circ}$. The scaling of the boundary
$\theta_{\max}$ with relative tolerance $\eta$ is
$\theta_{\max} \sim \sqrt{6 \eta}$: each two decades of tighter
tolerance corresponds to one decade of tighter angle. A working rule:
$\sin\theta \approx \theta$ is safe to relative tolerance $\eta$
when $\theta \lesssim 2.4\sqrt{\eta}$ in radians.
\end{exercise}

\begin{exercise}
Estimate the number of terms of the Maclaurin series for $e^{x}$
needed to compute $e^{1}$ to within $10^{-9}$, using the Lagrange
remainder bound with $M = e \approx 2.72$. Compare to the number of
terms needed to compute $e^{0.1}$ to the same tolerance.
\paragraph{Solution sketch.} The Lagrange bound at $x = 1$ is
$|R_{n}(1)| \leq M \cdot 1^{n+1}/(n+1)! = e/(n+1)!$. Setting
$e/(n+1)! \leq 10^{-9}$ requires $(n+1)! \geq 2.72 \times 10^{9}$.
Since $12! \approx 4.79 \times 10^{8}$ and $13! \approx 6.23 \times
10^{9}$, the bound is met at $n + 1 = 13$, that is $n = 12$, so
thirteen terms (indices $0$ through $12$) suffice.

For $e^{0.1}$ the bound is $e^{0.1} \cdot (0.1)^{n+1}/(n+1)!$. Each
extra term multiplies the bound by $0.1/(n+2)$, so the bound shrinks
roughly an order of magnitude with each term added. At $n = 6$ the
bound is $1.1 \cdot 10^{-7}/5040 \approx 2.2 \times 10^{-11}$, well
below $10^{-9}$; at $n = 5$ it is $1.1 \cdot 10^{-6}/720 \approx
1.5 \times 10^{-9}$, marginal. Six terms suffice for $e^{0.1}$.

The factor of about two in series length to extend from $x = 0.1$
to $x = 1$ reflects the geometric prefactor $x^{n+1}$ in the
remainder bound. The factorial denominator is what makes the
exponential cheap to evaluate at any fixed $x$; only its competition
with $x^{n+1}$ sets the truncation order.
\end{exercise}

\begin{exercise}
Estimate the relative error in the small-angle pendulum period $T = 2
\pi \sqrt{L/g}$ at amplitude $\theta_{0} = 30^{\circ}$, using the
correction $T(\theta_{0})/T_{0} \approx 1 + \theta_{0}^{2}/16$. State
the implicit assumption about the next-order correction.

\paragraph{Solution sketch.} Convert $30^{\circ}$ to radians:
$\theta_{0} = \pi/6 \approx 0.5236$ rad. The leading correction
gives $T/T_{0} \approx 1 + (0.5236)^{2}/16 = 1 + 0.01714 \approx
1.0171$, so the leading-order error in using $T_{0}$ is about
$1.7\%$ at $30^{\circ}$.

The implicit assumption is that the next-order correction term
$11 \theta_{0}^{4}/3072$ is negligible. At $\theta_{0} = 0.5236$ rad
this term equals $11 \cdot (0.5236)^{4}/3072 = 11 \cdot 0.0751/3072
\approx 2.69 \times 10^{-4}$, that is about $0.027\%$. The
next-order correction shifts the result by less than $0.03\%$, so
the leading-order correction $\theta_{0}^{2}/16$ is reliable to about
the third significant figure at this amplitude. A reader who
restricts attention to $\theta_{0} \leq 30^{\circ}$ can use the
two-term formula confidently to one-percent accuracy and the
one-term formula only at amplitudes below about $8^{\circ}$ if the
same one-percent tolerance applies.
\end{exercise}

\begin{exercise}
Estimate the speed at which the leading relativistic correction to
the kinetic energy contributes $1\%$ of the Newtonian value. Express
the answer as a fraction of the speed of light and as a speed in
$\si{\meter\per\second}$.

\paragraph{Solution sketch.} The Newtonian kinetic energy is
$K_{\text{N}} = m v^{2}/2$. The leading relativistic correction
from the expansion is $(3/8) m v^{2} \beta^{2}$ with $\beta = v/c$.
The ratio of correction to Newtonian is $(3/8) m v^{2} \beta^{2}
\div (m v^{2}/2) = (3/4) \beta^{2}$. Setting this equal to $0.01$
gives $\beta^{2} = 4/300 \approx 0.01333$, $\beta \approx 0.1155$.
The speed is $v = \beta c = 0.1155 \times 2.998 \times 10^{8}\,
\si{\meter\per\second} \approx 3.46 \times 10^{7}\,
\si{\meter\per\second}$, that is, about $11.5\%$ of the speed of
light or $35{,}000\,\si{\kilo\meter\per\second}$. Below this speed
the Newtonian kinetic energy is accurate to one percent; above it
the relativistic correction is the engineer's first responsibility
to include.
\end{exercise}

\subsection*{Simulation}\pathtag{mastery}

\begin{exercise}
In a language of the reader's choice, compute the partial sums of the
harmonic series at $N = 10, 100, 10^{3}, 10^{4}, 10^{5}, 10^{6}$, and
plot them against $\log N$. Compare to the asymptotic $\log N + \gamma$.
Identify the regime in which the asymptotic is accurate to better than
$1\%$.

\paragraph{Solution sketch.} Running \texttt{code/harmonic\_partial\_sums.py}
produces the table in \texttt{data/harmonic\_decades.csv}: $H_{10} =
2.929$ vs asymptotic $2.880$ ($1.7\%$ gap); $H_{100} = 5.187$ vs
$5.182$ ($0.1\%$); $H_{1000} = 7.485$ vs $7.485$ ($0.007\%$). The
gap is approximately $1/(2N)$ (from the next term in the asymptotic
expansion $H_{N} = \log N + \gamma + 1/(2N) + O(1/N^{2})$, given in
\cite{text:v2ch04-knuth-taocp1-1997}), which falls below $1\%$ once
$N \gtrsim 50$, that is, the asymptotic is accurate to better than
$1\%$ for $N \geq 50$ approximately. The asymptotic improves linearly
in $1/N$, so each decade of larger $N$ gives one decade of tighter
agreement.
\end{exercise}

\begin{exercise}
Numerically evaluate the partial sums of the alternating harmonic
series at $N = 10, 100, 10^{3}, 10^{4}$, and compare to $\log 2$.
Verify that the truncation error is bounded by the magnitude of the
first omitted term.

\paragraph{Solution sketch.} Running \texttt{code/alternating\_harmonic.py}
produces, with $\log 2 = 0.693147$:
\begin{itemize}[leftmargin=*]
  \item $N = 10$: $S_{10} = 0.6456$, gap $4.77 \times 10^{-2}$,
    bound $1/11 \approx 9.09 \times 10^{-2}$.
  \item $N = 100$: $S_{100} = 0.6882$, gap $4.95 \times 10^{-3}$,
    bound $1/101 \approx 9.90 \times 10^{-3}$.
  \item $N = 1000$: $S_{1000} = 0.6927$, gap $5.00 \times 10^{-4}$,
    bound $\approx 1 \times 10^{-3}$.
  \item $N = 10\,000$: $S_{10000} = 0.6931$, gap $\approx 5 \times
    10^{-5}$, bound $\approx 1 \times 10^{-4}$.
\end{itemize}
The empirical gap is consistently about half the bound (the
alternating-series bound is loose by a factor of two when $b_{n}$
shrinks slowly), but in every row the gap is well within the bound,
confirming the alternating-series error inequality.
\end{exercise}

\begin{exercise}
Compute the Taylor polynomial of $\sin x$ at $x_{0} = 0$ to orders
$1, 3, 5, 7, 9$ and tabulate the absolute error at $x = 0.5, 1.0,
1.5, 2.0$. Compare the empirical error to the Lagrange remainder
bound.

\paragraph{Solution sketch.} Running \texttt{code/taylor\_sin\_error.py}
produces \texttt{data/taylor\_sin\_orders.csv}. Sample rows: at
$x = 1.0$, the empirical errors are $1.59 \times 10^{-1}$, $8.33 \times
10^{-3}$, $1.98 \times 10^{-4}$, $2.76 \times 10^{-6}$, $2.51 \times
10^{-8}$ for orders $1, 3, 5, 7, 9$. The Lagrange bound $x^{n+1}/(n+1)!$
gives $1.67 \times 10^{-1}$, $8.33 \times 10^{-3}$, $1.98 \times
10^{-4}$, $2.76 \times 10^{-6}$, $2.51 \times 10^{-8}$. The match is
nearly tight: the empirical error is within $5\%$ of the bound at
every order for $|x| \leq 2$, because the bound uses $M_{n+1} = 1$
and the derivatives of $\sin$ all have magnitude exactly $1$ at the
extremising $\xi$ for small enough $x$. The figure
\autoref{fig:vol02:ch04:taylor-error} plots the same data on log
scale.
\end{exercise}

\begin{exercise}
Implement Newton's iteration $a_{n+1} = a_{n} - (a_{n}^{2} - 2)/(2
a_{n})$ for $\sqrt{2}$ starting from $a_{0} = 1$. Tabulate $a_{n}$
and the residual $a_{n}^{2} - 2$ for $n = 0, 1, \ldots, 6$. Verify
that the residual squares at each step (quadratic convergence).

\paragraph{Solution sketch.} Running \texttt{code/newton\_sqrt2.py}
gives $a_{0} = 1.0000$, $a_{1} = 1.5000$, $a_{2} = 1.4167$,
$a_{3} = 1.41422$, $a_{4} = 1.4142136$, $a_{5} = 1.41421356237$,
$a_{6} = 1.41421356237$ (double precision converged). The
residuals are $a_{n}^{2} - 2 = -1.0, 0.25, 6.94 \times 10^{-3},
6.01 \times 10^{-6}, 4.51 \times 10^{-12}, 2.0 \times 10^{-16}, \ldots$.
The successive residuals do not exactly square (the Newton constant
involves $1/(2 a_{n}) \approx 1/(2\sqrt{2}) \approx 0.354$), but
their orders of magnitude approximately double in negative powers of
ten: $-1, -1, -3, -6, -12, -16$, where the last is at the
floating-point noise floor. The pattern is the signature of quadratic
convergence.
\end{exercise}

\subsection*{Design}\pathtag{standard}

\begin{exercise}
Design a series-truncation criterion for evaluating $e^{x}$ on a
microcontroller without floating-point hardware, using the Maclaurin
series and a stated relative error tolerance of $10^{-6}$ on the
range $x \in [0, 1]$. Specify the truncation order and the storage
requirement (number of coefficients).

\paragraph{Discussion.} The Lagrange bound on $[0, 1]$ is
$|R_{n}(x)| \leq e \cdot x^{n+1}/(n+1)!$. Worst-case over $x \in
[0, 1]$ is at $x = 1$: $e/(n+1)! \leq 10^{-6}$ requires $(n+1)!
\geq 2.72 \times 10^{6}$, met at $n + 1 = 10$, that is $n = 9$.
Ten coefficients (orders $0$ through $9$) and ten multiply-add
operations evaluate the polynomial via Horner's scheme. The
coefficients are $1, 1, 1/2, 1/6, 1/24, 1/120, 1/720, 1/5040,
1/40320, 1/362880$; these can be stored as integers (numerators
$1$, with denominators $0!, 1!, \ldots, 9!$ for the integer-arithmetic
implementation) or as fixed-point with the integer multiplier
$362880$ as the common denominator.

A range-reduction step extends the design beyond $[0, 1]$: write
$x = k \log 2 + r$ with $r \in [0, \log 2)$, evaluate $e^{r}$ by
the series, and multiply by $2^{k}$. The reduced range $[0, \log 2)$
needs $n = 6$ (seven coefficients), reducing the work by about a
third at the cost of one branchless reduction. A production
microcontroller library typically uses this two-step pattern.
\end{exercise}

\begin{exercise}
Design a numerical scheme that uses the geometric series formula $1/(1
- r) = \sum r^{n}$ to evaluate $1/x$ for $x$ in $[0.5, 2]$ without
division. Specify the substitution $x = 1 - r$, the convergence
condition, and the truncation order required for $10^{-4}$ relative
accuracy.

\paragraph{Discussion.} The substitution $x = 1 - r$ gives $r = 1 -
x$. The series $1/(1 - r) = \sum r^{n}$ converges absolutely for
$|r| < 1$, which translates to $0 < x < 2$. Over $x \in [0.5, 2]$,
the maximum $|r|$ is at the endpoints: at $x = 0.5$, $r = 0.5$; at
$x = 2$, $r = -1$ (boundary, divergent). The scheme as stated fails
near $x = 2$.

A common fix splits the range: for $x \in [0.5, 1]$ use $r = 1 - x
\in [0, 0.5]$ directly; for $x \in (1, 2]$ use $x = 1 / y$ with
$y \in [0.5, 1)$, so $1/x = y$ is exactly what one was originally
trying to compute. The roundabout suggests another substitution
$x = 2(1 - s)$, $s = 1 - x/2$, giving $1/x = (1/2)/(1 - s) =
(1/2) \sum s^{n}$ with $s \in (0, 0.5]$ for $x \in [1, 2)$. Either
substitution keeps $|r|, |s| \leq 0.5$.

For $10^{-4}$ relative accuracy at the worst-case $|r| = 0.5$, the
truncation tail is $r^{N}/(1 - r) / S_{\infty} = r^{N}$; require
$0.5^{N} \leq 10^{-4}$, that is $N \geq 14$. Fourteen multiply-adds
without division gives $1/x$ to $10^{-4}$ relative accuracy across
$[0.5, 2]$ after the two-piece range reduction. The microcontroller
saves a costly hardware divider at the cost of fourteen
multiplications.
\end{exercise}

\begin{exercise}
Design a small-angle decision rule for a device that mixes
approximate-trigonometry computation (faster, less accurate) and
exact-trigonometry computation (slower, accurate). State the angle
threshold above which the device must use the exact form, given an
error budget of $0.5\%$ on the sine.

\paragraph{Discussion.} Relative error of the leading-order $\sin
\theta \approx \theta$ is $\theta^{2}/6$. Setting $\theta^{2}/6 \leq
5 \times 10^{-3}$ gives $\theta \leq \sqrt{0.03} \approx 0.173$ rad
$\approx 9.9^{\circ}$. The threshold is approximately $10^{\circ}$.

A working rule reads:

\begin{itemize}[leftmargin=*]
  \item If $|\theta| \leq 10^{\circ}$: use $\sin\theta \approx
    \theta$. Cost is one floating-point read.
  \item If $10^{\circ} < |\theta| \leq 30^{\circ}$: use the
    cubic-correction form $\theta - \theta^{3}/6$. Cost is three
    multiplies and one subtract; the budget is still met because
    the next-order relative error is bounded by $\theta^{4}/120
    \leq 0.0023\%$ at $30^{\circ}$.
  \item If $|\theta| > 30^{\circ}$: call the exact-trigonometry
    library routine. Cost is the library call (often a CORDIC
    iteration on resource-constrained hardware).
\end{itemize}

The two-threshold rule covers an extended safe range with only the
extra cubic term, paying for the next-order arithmetic only where
the leading-order budget runs out. A single-threshold rule (using
the exact routine above $10^{\circ}$) is simpler but cedes accuracy
on the middle range where the cubic correction is nearly free.
\end{exercise}

\subsection*{Diagnosis and reverse engineering}\pathtag{core}

\begin{exercise}
A colleague reports that the series $\sum 1/(n \log n)$ converges
because its terms decrease faster than $1/n$. State whether the
conclusion is correct and identify the test that settles the question.

\paragraph{Solution sketch.} The conclusion is wrong. The terms
$1/(n \log n)$ do shrink faster than $1/n$, but the series diverges.
The integral test settles it: with $f(x) = 1/(x \log x)$, the
substitution $u = \log x$, $\dd u = \dd x / x$ gives
\[
\int_{2}^{N} \frac{\dd x}{x \log x} = \int_{\log 2}^{\log N}
\frac{\dd u}{u} = \log \log N - \log \log 2,
\]
which grows without bound as $N \to \infty$. Therefore the
corresponding series diverges. The shrinking-faster-than-$1/n$ test
is necessary for convergence (the term test, in reverse) but is
emphatically not sufficient. The borderline cases near $1/n$ require
the integral test or comparison to a known borderline series. The
next series $\sum 1/(n (\log n)^{2})$ does converge, by the same
integral test with $u = \log x$ producing $\int \dd u/u^{2}$, finite.
\end{exercise}

\begin{exercise}
A computation of $\sin(45^{\circ})$ on a calculator returns
$0.7071067811$. The reader uses the Maclaurin series $\sin x \approx x
- x^{3}/6 + x^{5}/120 - x^{7}/5040 + x^{9}/362880$ and gets a value
that agrees to within roughly $4 \times 10^{-8}$. Identify the
expected truncation error from the Lagrange remainder bound and
verify it against the empirical residual.

\paragraph{Solution sketch.} At $\theta = 45^{\circ} = \pi/4 \approx
0.7854$ rad, the Lagrange bound at order $9$ (five nonzero terms,
last kept term is $x^{9}/362880$) is
\[
|R_{9}(\theta)| \leq \frac{\theta^{11}}{11!} =
\frac{0.7854^{11}}{39\,916\,800} \approx \frac{0.0742}{3.99 \times
10^{7}} \approx 1.86 \times 10^{-9}.
\]
The next-term magnitude (alternating-series bound) is $\theta^{11}/11!
\approx 1.86 \times 10^{-9}$. The empirical residual of $4 \times
10^{-8}$ is about twenty times the bound, which suggests either an
arithmetic round-off contribution in the calculator's polynomial
evaluation, or that the reader is measuring against a less precise
reference for $\sin(45^{\circ}) = 1/\sqrt{2}$. The Lagrange bound at
this order predicts the truncation error is at the $10^{-9}$ level;
any larger discrepancy is round-off rather than truncation.
\end{exercise}

\begin{exercise}
A spreadsheet evaluates $f(x) = (e^{x} - 1)/x$ at $x = 10^{-12}$ by
direct computation and returns $0$. Identify the failure mechanism
and propose a Taylor-series formula for $f(x)$ that is numerically
stable near $x = 0$.

\paragraph{Solution sketch.} The failure is catastrophic
cancellation. At $x = 10^{-12}$, $e^{x} = 1 + x + O(x^{2})
\approx 1 + 10^{-12}$. Double precision stores $1 + 10^{-12}$ as
$1.0000000000010 \ldots$ with about three significant digits to
spare past the implicit $1$. The subtraction $e^{x} - 1$ then yields
roughly $10^{-12}$ but with no preserved precision; depending on the
rounding mode, the result may be exactly zero (when $e^{x}$ rounds
to exactly $1$ at quad precision). Dividing the resulting zero or
near-zero by $x = 10^{-12}$ gives $0$ or a meaningless value.

The Taylor series for $f(x) = (e^{x} - 1)/x$ is
\[
f(x) = \sum_{n=0}^{\infty} \frac{x^{n}}{(n+1)!} = 1 + \frac{x}{2} +
\frac{x^{2}}{6} + \frac{x^{3}}{24} + \cdots
\]
This series evaluates without cancellation: at $x = 10^{-12}$, the
leading term is $1$ and the correction is $5 \times 10^{-13}$, well
within floating-point precision. The branched implementation uses
the series for $|x| < \epsilon$ (with $\epsilon$ chosen so the cubic
term is below floating-point precision) and the direct formula for
larger $|x|$. The series gives a stable answer where the direct
formula fails. Cleve Moler's \emph{Numerical Computing with
MATLAB} \cite{text:moler-numerical} uses this exact case as a
canonical example of cancellation and its series-based fix.
\end{exercise}

\subsection*{Failure analysis}\pathtag{core}

\begin{exercise}
A control loop uses the small-angle approximation $\sin \theta
\approx \theta$ and operates with $\theta$ ranging up to $40^{\circ}$.
Argue, in 200 to 300 words, whether the approximation is safe in this
range, identify the angle at which the next-order correction becomes
mandatory, and state the consequence of using the leading-order
approximation past its budget.

\paragraph{Solution sketch.} The leading-order relative error of
$\sin\theta \approx \theta$ scales as $\theta^{2}/6$. At
$\theta = 40^{\circ} \approx 0.698$ rad, the relative error is
$0.698^{2}/6 \approx 8\%$. A control loop that nominally regulates
to within a few percent has been pushed well past the approximation
budget at the upper end of its range.

The next-order correction $\sin\theta \approx \theta - \theta^{3}/6$
has residual $\theta^{5}/120$, which at $40^{\circ}$ is $0.698^{5}/120
\approx 1.4 \times 10^{-3}$, that is about $0.16\%$ relative. The
cubic-correction form is within a one-percent budget across the full
$0$--$40^{\circ}$ operating range. The crossover where the next-order
correction becomes mandatory (relative error of leading-order $\geq
1\%$) is around $\theta = \sqrt{6 \cdot 0.01} \approx 0.245$ rad
$\approx 14^{\circ}$.

The consequence of using the leading-order form past its budget in a
control loop is a systematic gain error: at $40^{\circ}$ the
controller thinks it is applying a command of magnitude $0.698$ rad
of effective angle but is actually producing $\sin(40^{\circ}) =
0.643$ rad of physical effect. The loop overshoots its target by
about $8\%$ at the upper end of the range, with the error scaling
quadratically toward the corners. In a stability analysis the
phase-margin calculation built around the linearised model is
mispredicted by a similar fraction, and a feedback loop tuned for
$30^{\circ}$ may oscillate at $40^{\circ}$. The fix is the
cubic-correction form throughout, or a piecewise rule that switches
at $14^{\circ}$.
\end{exercise}

\begin{exercise}
A signal-processing routine uses a power-series approximation of a
nonlinear sensor characteristic, truncated at the cubic term. The
sensor is then operated outside its calibrated range, where the
quartic and higher terms become significant. Argue, in 200 to 300
words, what the user observes, what the diagnostic for the regime
shift looks like, and the procedural fix.

\paragraph{Solution sketch.} The truncated cubic model $y \approx
c_{0} + c_{1} x + c_{2} x^{2} + c_{3} x^{3}$ has bounded error inside
the calibrated range $[x_{\min}, x_{\max}]$ where the higher-order
terms $c_{4} x^{4} + \cdots$ contribute below the calibration's
target tolerance. Outside this range, the higher-order terms grow
faster than the cubic, and the model error grows polynomially with
the deviation from the calibration boundary.

The user observes that the sensor readings are systematically wrong
when the physical input strays past calibration boundaries: at twice
the calibrated maximum, a model that was good to $0.1\%$ inside is
typically wrong by tens of percent. The bias is one-sided (sensor
nonlinearities are rarely symmetric), so the error appears as a
drift rather than as zero-mean noise. Outputs near the upper boundary
of the calibrated range drift smoothly into bias-dominated readings
past it; the user may not notice until two independent measurements
disagree.

The diagnostic for the regime shift is a residual plot of measured
vs predicted, or measured vs a reference instrument, taken across
the operating range. A clean residual inside the calibrated range
that fans out past the boundary identifies the problem. The
procedural fix is recalibration over the extended range (adding
calibration points beyond the original boundary), refitting with a
higher-order polynomial (a quintic or septic) that contains the
quartic and higher terms, and updating the field-deployed coefficient
table. The lesson, more generally: a truncated Taylor or polynomial
expansion is only as valid as the calibrated range over which it
was fitted, and operating past that range silently corrupts the
output.
\end{exercise}

\begin{exercise}
A numerical iteration produces residuals that decrease but do not
square at each step. Argue, in 200 to 300 words, what this tells the
analyst about the convergence rate of the iteration, what the
expected total iteration count is to reach a target tolerance, and
which of three named iterations (Newton, secant, fixed-point) is
consistent with the observed behaviour.

\paragraph{Solution sketch.} Quadratic convergence (Newton) shows
residual squaring: the number of correct digits doubles each
iteration, and the residual at step $n+1$ is approximately $C$ times
the square of the residual at step $n$. Failure to observe squaring
implies the iteration is not exhibiting Newton-style quadratic
convergence.

Three patterns are consistent with non-squaring residuals.

\textit{Fixed-point (linear) convergence}: $e_{n+1} \approx r e_{n}$
with $|r| < 1$. The residual shrinks by a constant factor per step.
The number of correct digits grows linearly in $n$. Iteration count
to tolerance $\eta$ from starting error $e_{0}$ is approximately
$\log(e_{0}/\eta) / \log(1/|r|)$. For $r = 0.5$ this is $3.3 \times
\log_{10}(e_{0}/\eta)$, that is about thirty iterations for a
nine-decade tolerance reduction.

\textit{Superlinear (secant) convergence}: $e_{n+1} \approx C
e_{n}^{\phi}$ with $\phi = (1 + \sqrt 5)/2 \approx 1.618$. The
secant method achieves this order without computing derivatives.
Residuals shrink faster than linear but slower than quadratic; the
ratio $e_{n+1}/e_{n}^{2}$ tends to zero (not bounded above), while
$e_{n+1}/e_{n}$ tends to zero (faster than linear). Iteration count
to nine-decade reduction is typically eight to twelve, intermediate
between linear and quadratic.

\textit{Newton with a multiple root}: at a root of multiplicity $m$
the Newton iteration degrades from quadratic to linear convergence
with rate $1 - 1/m$. A reader observing residual shrinkage by a
constant factor per step at a known Newton iteration should suspect
a multiple root; the fix is modified Newton with multiplicity
adjustment.

The observed behaviour is consistent with fixed-point iteration if
the rate is roughly constant, with the secant method if it
super-linearises but never quadrupes correct digits, and with
Newton-at-a-multiple-root if the convergence is degraded from a
known second-order scheme. The convergence-rate ladder is treated
in \cite{text:v2ch04-hairer-wanner-2008}.
\end{exercise}

\subsection*{Judgment}\pathtag{standard}

\begin{exercise}
A junior engineer says: "Any approximation is good enough as long as
the next term is smaller than the error budget." Write a one-paragraph
response that takes the junior engineer seriously, identifies what is
correct in the statement, identifies what is missing (the role of
absolute versus relative error and the role of the next-derivative
bound), and proposes a sharper rule.

\paragraph{Discussion.} The statement is correct for an alternating
series with monotone-decreasing magnitudes: the alternating-series
test gives exactly the bound the junior engineer cites. It is also
correct, approximately, for a Taylor series at a point well inside
its radius of convergence, where the next term is a reliable proxy
for the Lagrange remainder. What is missing is twofold. First,
absolute and relative budgets behave differently: the next term may
be smaller than an absolute budget in absolute terms while still
being a large fraction of the quantity being computed, so a relative
budget must be tested against the next term divided by the current
estimate. Second, the next-term proxy is only safe when the next
derivative bound is well-behaved over the interval; if $|f^{(n+1)}|$
spikes far from the working point, the next term in the series
understates the remainder and the budget rule fails. A sharper rule:
the truncation is acceptable when the Lagrange remainder bound
$M_{n+1} |x - x_{0}|^{n+1}/(n+1)!$ is below the budget, where
$M_{n+1}$ is an explicit upper bound on $|f^{(n+1)}|$ across the
interval of interest. The next-term proxy is the fast check; the
Lagrange bound is the rigorous one.
\end{exercise}

\begin{exercise}
A safety analysis uses the leading-order small-angle approximation
$\sin \theta \approx \theta$ in deriving an upper bound on a hazard.
Argue, in 200 to 300 words, whether this is a conservative or an
optimistic choice, when the leading-order approximation overstates
the bound and when it understates it, and what the safety analyst
should do to make the approximation defensible.

\paragraph{Discussion.} The approximation $\sin\theta \approx \theta$
overstates $\sin\theta$ for $\theta > 0$, because the next nonzero
term of the Taylor series is $-\theta^{3}/6 < 0$ at $\theta > 0$.
For a safety bound proportional to $\sin\theta$, using $\theta$
instead of the true value gives a value larger than the true one,
which is conservative (the bound exceeds the actual hazard
contribution).

Whether this conservative direction is the safe direction depends on
what the bound is used for.

\textit{If the bound is an upper bound on something the analyst
wants to limit} (peak stress, displacement, exposure), using
$\theta$ gives a conservative (safe) upper estimate: the actual
value is less than the computed bound, so a design that satisfies
the bound is safer than computed.

\textit{If the bound is a lower bound on something the analyst
wants to guarantee} (clearance, margin, separation), using $\theta$
overstates the quantity, so the apparent margin is larger than the
actual margin. This is optimistic, not conservative; the design may
have less clearance than computed.

\textit{If the bound enters a quadratic or higher-power
expression}, the direction of bias may flip; carrying the bias
through the algebra is required.

The defensible procedure is: state the budget explicitly (relative
tolerance $\eta$); compute the maximum operating angle
$\theta_{\max}$ implied by the budget ($\theta_{\max} = \sqrt{6
\eta}$); state in the safety case which way the approximation
biases the bound at $\theta_{\max}$; and either include the
next-order term or document that the approximation is conservative
in the relevant direction. Safety analysts who write "we use the
small-angle approximation throughout" without naming the bias
direction give up the documentation discipline the standards
demand.
\end{exercise}

\begin{exercise}
Argue, in one paragraph, whether engineering training should
emphasise the formal epsilon-delta definition of a limit or the
working limit-laws-and-Taylor-series machinery. Take a position and
defend it with a specific example from a working engineering problem.

\paragraph{Discussion.} The engineering position is that working
limit-laws-and-Taylor-series machinery is the daily tool and the
epsilon-delta definition is the truth condition that the working
tools inherit. An engineer who can write a Taylor expansion with a
remainder bound, compute the limit of a composite function from the
limit laws, and recognise that a tolerance argument has an implicit
epsilon-delta content has the right grasp of limits for design and
analysis. The specific example: a feedback-loop stability analysis
asks whether the closed-loop response stays within $1\%$ of the
desired step response for all gain settings within a stated range.
This is an epsilon-delta question by structure (for every operating
tolerance, find a gain tolerance that keeps the response within
spec), but the working analysis uses the Bode-magnitude limit
machinery and the small-signal Taylor expansion of the actuator
nonlinearity. A student who can do the Bode analysis without
explicit epsilon-delta arguments but who can recognise when the
small-signal expansion has been pushed past its budget is doing the
work an engineer needs to do. The epsilon-delta discipline matters
when the student must read or write a proof that supports a
liability-bearing claim, or when the working tools fail and the
problem must be reduced to first principles. The training, then,
emphasises the working tools, with the formal definition presented
as the certifier and reserved for cases where the working tools do
not settle the question.
\end{exercise}
